% ============================================================
%  R. Nielsen, Forelæsninger over „Philosophisk Propædeutik“
%  fra Universitetsaaret 1860--61.
%  Kjøbenhavn: Forlagt af den Gyldendalske Boghandling (F. Hegel);
%  Thieles Bogtrykkeri, 1862.  (482 pp., Fraktur)
%  TRANSCRIPTION (Danish)
%  Source: Google scan of the Koninklijke Bibliotheek (The Hague) copy,
%          bibliotek/Nielsen, Rasmus/1862-forelaesninger_over_Philosophisk_propaed-2.pdf
%          (byte-identical duplicate there: propaedeutik-1861.pdf).
%          A second Google scan (British Library copy) is at
%          https://books.google.com/books?id=83xZAAAAcAAJ and serves as a
%          second witness where the working scan is doubtful.
%  PAGE MAP:  see pagemap.py — uniform, printed 1--482 = PDF + 12.
%          Verified by sweeping the header numeral of every PDF page 13--494
%          via the scan's text layer: no offset change, no duplicated leaves.
%          Endpoints to be re-confirmed against the page images by the first
%          and last batch agents.
%          Front matter (unpaginated): PDF 3--8 title leaves (two settings,
%          both 1862, the fuller naming F. Hegel), PDF 9--10 Forord,
%          PDF 11--12 Indhold.
%  STRUCTURE (from the Indhold; page numbers of the three parts verified
%          against the header sweep, sub-sections still to be verified
%          against the printed heads as transcription reaches them):
%            I.   Philosophiens Begreb                      pp. 1--29
%            II.  Philosophien og de særskilte Videnskaber  pp. 30--476
%                 (A. Mathematiken: Uendelighedsproblemet, from p. 39;
%                  B. Chemien: Forandringsproblemet, from p. 81;
%                  Resultat, from p. 167;
%                  C. Physiologien: Livsproblemet, from p. 177)
%            III. Philosophien og de akademiske Studier     pp. 477--482
%  CONVENTIONS: subdivision marks in this book run a) b) c), then Greek
%          α) β) γ) and doubled αα) ββ) γγ) — the OCR reads the Greek as
%          a/6/7/B/y; restore the Greek (textalpha is loaded).
%          The Darwin passage pp. 326--328 was previously transcribed as
%          the excerpt edition texts/nielsen/propaedeutik-darwin/.
%  FOOTNOTE MARKS: *) for a page's first note, **) for its second (both
%          occur on p. 58). Rendered via the footnote counter with an
%          explicit \setcounter{footnote}{0} after the % --- p. N ---
%          marker of every printed page carrying a footnote; see the
%          preamble note. Footnotes so far: pp. 15 (runs over to 16),
%          34, 48, 55, 58 (two).
%  PRINTER'S DEFECTS: none logged yet.
% ============================================================
\documentclass[12pt,a4paper]{book}
\usepackage[utf8]{inputenc}
\usepackage[T1]{fontenc}
\usepackage{libertinus}
\usepackage{libertinust1math}
\usepackage{textalpha}   % Greek typed directly (α) β) γ) subdivision marks)
\usepackage{graphicx}    % for \reflectbox (the old Danish »ɔ:« id-est mark)
\DeclareUnicodeCharacter{0254}{\reflectbox{c}}  % ɔ (U+0254) = reversed c
\usepackage[danish]{babel}
\usepackage[protrusion=true,expansion=true]{microtype}
\usepackage{setspace}
\usepackage[margin=1.2in]{geometry}
\usepackage{enumitem}
\usepackage{fancyhdr}
\setlength{\headheight}{28pt}
\addtolength{\topmargin}{-13.5pt}

% Footnotes as printed: *) for a page's first note, **) for its second
% (p. 58 carries both). Not \fnsymbol (math mode, U+2217) and not
% footmisc[perpage] (resets on the TYPESET page, not the printed one).
% Instead the mark follows the footnote counter, and the counter is reset
% by an explicit \setcounter{footnote}{0} placed after the % --- p. N ---
% marker of EVERY printed page that carries a footnote. Batch agents must
% include that reset line.
\renewcommand{\thefootnote}{\ifcase\value{footnote}\or *)\or **)\or ***)\fi}

\usepackage{hyperref}

\hypersetup{
  pdftitle={Forelæsninger over „Philosophisk Propædeutik“ fra Universitetsaaret 1860--61},
  pdfauthor={R. Nielsen},
  hidelinks
}

\pagestyle{fancy}
\fancyhf{}
\fancyhead[LE,RO]{\thepage}
\fancyhead[RE]{\textit{Forelæsninger over „Philosophisk Propædeutik“}}
\fancyhead[LO]{\textit{\leftmark}}

\setstretch{1.3}

\title{\textbf{Forelæsninger}\\[4pt]
  \large over\\[4pt]
  \textbf{\large „Philosophisk Propædeutik“}\\[6pt]
  \large fra Universitetsaaret 1860--61}
\author{R.\ Nielsen}
\date{Kjøbenhavn: Forlagt af den Gyldendalske Boghandling (F.\ Hegel), 1862}

\begin{document}
\maketitle
\thispagestyle{empty}
\newpage

\tableofcontents
\newpage

%% ============================================================
%%  FORORD (unpaginated, PDF 9--10)
%% ============================================================
\chapter*{Forord}
\addcontentsline{toc}{chapter}{Forord}
\label{ch:forord}

% The pp. 1--14 fragment below carries: the Forord text (marked
% "% --- Forord, PDF 9 ---" and "% --- Forord, PDF 10 ---"), then the
% chapter head for I. Philosophiens Begreb as printed on p. 1, then the
% body of pp. 1--14. Later fragments likewise carry any chapter or
% section head that falls at the start of, or inside, their range —
% chapter heads live in the fragments, not in this skeleton, so that
% each fragment still replaces exactly one marker.
% --- Forord, PDF 9 ---
Dersom nærværende Skrift tildrager sig nogen Opmærksomhed, vil
Interessen uden Tvivl være af dobbelt Art. Med Spørgsmaalet om
Indholdets mulige Værd vil et Fleertal af Læsere, holdende sig til
Titlens Bogstav, vistnok forene et andet, mere praktisk Spørgsmaal:
\emph{efter hvilke Principer og i hvilket Omfang er den philosophiske
Propædeutik her foredraget for de Studerende?}

Det vil da, med Hensyn paa det sidste Spørgsmaal, ikke være nok med
en Bemærkning som den, at Fremstillingens Motiver, Grundtanker og
hele Anlæg vel svare til de mundtlige Foredrag, men at Adskilligt,
hvad det Enkelte angaaer, er forandret, Noget tilføiet, Andet udeladt
o.~s.~v.; Læseren vil ansee sig berettiget til at vente en tydeligere
Forklaring.

Det bemærkes derfor udtrykkelig, at med Undtagelse af nogle mindre
betydelige Redactionsforandringer er hele første Afdeling S. 1--30 en
ligefrem Gjengivelse af de holdte Forelæsninger. Det Samme gjælder om
anden Afdelings indledende Stykke: \emph{Almindeligt} S. 30--39, om
Philosophiens Forhold til Mathematiken S. 39--63, om Philosophiens
Forhold til Chemien S. 81--100, S. 135--177,
% --- Forord, PDF 10 ---
og om Philosophiens Forhold til Physiologien S. 177--180, S. 199--288,
S. 386--437, S. 475--482.

I Fremstillingen af de øvrige Partier ere derimod betydelige
Forandringer foretagne; Afsnittet om de psychophysiske Undersøgelser
har i de mundtlige Foredrag slet ikke været berørt; skærpede og videre
gaaende Analyser ere her og der indførte, ligesom det da ogsaa
forstaaes af sig selv, at det Allermeste af Anmærkningernes Indhold er
tilføiet under Omarbeidelsen.

Da Stræben efter at give Behandlingen navnlig af de physiologiske
Problemer den Udførlighed, som Undersøgelser af en saa indviklet Natur
ikke vel kunne savne, medførte en betydelig Udvidelse af anden
Afdeling, er tredie Afdeling: Philosophien og de akademiske Studier,
bleven indskrænket til en blot Angivelse af Momenterne og kun tilføiet
som særlig Afdeling for at vise Gangen i det Hele.

\noindent Den 11te Juni 1862.

\begin{flushright}
\textbf{R.~Nielsen.}
\end{flushright}

% (double rule below the signature, closing the Forord as printed)

%% ============================================================
%%  I. PHILOSOPHIENS BEGREB  (pp. 1--29)
%% ============================================================
\chapter*{I.\\ Philosophiens Begreb.}
\addcontentsline{toc}{chapter}{I. Philosophiens Begreb}

% --- p. 1 ---
% (chapter opening; no header numeral on this page, signature "1" at foot)
\section*{A.\\ Modsigelseslære.}
\addcontentsline{toc}{section}{A. Modsigelseslære}

\textbf{Hvad er Philosophie?} „En Videnskab fuld af Modsigelser!“
Denne Definition, hvormeget den end, fra een Side betragtet, synes
anlagt paa at bryde Staven over al Philosophie, er ikke destomindre
begrundet saavel i Videnskabens Historie, som i Sagens Natur. For
enhver anden Videnskab vilde en saadan Tilstaaelse vel være
tilintetgjørende, for Philosophien afgiver den et væsentligt
Udgangspunkt, et paalideligt Støttepunkt, et fast, urokkeligt
Grundlag. At Philosophien maa indeholde Modsigelser, er jo en
Selvfølge, eftersom det netop er den Videnskab, der overalt, saavel i
Tilværelsen, som i den sunde Forstands Opfattelse af Tilværelsen
(Propæd. S. 40--44), opdager og paaviser Modsigelser; Philosophien
bliver just med Hensyn hertil at opfatte som en \emph{almindelig
Modsigelseslære}.

Som Modsigelseslære maa Philosophien imidlertid ikke blot opdage og
paavise Modsigelserne, den maa ogsaa prøve Modsigelserne; thi der er
Forskjel paa Modsigelser. Nogle Modsigelser ere nemlig overfladiske,
unødvendige, lette at undgaae, andre derimod paatrængende og
uafviselige; nogle Modsigelser indsnige sig paa Grund af Tankeløshed,
andre
% --- p. 2 ---
fremkunstles ved Spidsfindighed, andre afsløres ved Skarpsindighed;
nogle Modsigelser spille som i en Tankeleg, andre angive væsentlige
Knuder, vanskelige, tungt opløselige Problemer.

Til at oplyse, hvor nødvendigt det er, at der, hvad Modsigelserne
angaaer, gjøres Forskjel imellem det Overfladiske og det Væsentlige,
det Skarpsindige og det Spidsfindige, kunne følgende Exempler tjene:

\textbf{Løgneren} (ψευδόμενος). Denne Modsigelse fremsættes saaledes:
Naar En siger, at han lyver, lyver han saa med dette sit Udsagn, eller
siger han Sandhed? Her fordres til Svar enten Ja eller Nei. „Ja, han
lyver!“ --- Og siger det selv, nævner Tingen med det rette Navn,
siger, hvad sandt er, lyver altsaa ikke. „Nei, han lyver ikke, han
taler Sandhed!“ --- Man maa altsaa troe, hvad han siger; nu siger han
jo netop, at han lyver; men er det sandt, at han lyver, saa taler han
jo ikke Sandhed.

At her er en Modsigelse, kan den sunde Forstand let fatte; men hvad er
det for en Modsigelse, hvori stikker Feilen?

\textbf{Ligestore Størrelser.} At en Størrelse er liig sig selv, og at
to, tre . . Størrelser ere indbyrdes ligestore, deri seer den sunde
Forstand ingen Modsigelse; dersom det derimod blev paastaaet, at to
Størrelser umulig kunne være ligestore med en og samme tredie, og
følgelig indbyrdes ligestore, fordi det er en Modsigelse, vilde den
sunde Forstand føle sig krænket. Er her da virkelig nogensomhelst
Modsigelse i den anførte Sætning? For at indsee, hvad det egentlig
er, her tales om, maa man udtrykkelig mærke sig, at her ikke tales om
Ting, der have Størrelse, men om Størrelser. Tales her, ikke om
Størrelserne selv, men om Ting, der have Størrelse, t. Ex. et Æble, en
Trekant, en Time o.s.v., da er her ganske vist ikke Spor af
Modsigelse. At to Trekanter, som, hver for sig, ere ligestore med een
og samme
% --- p. 3 ---
tredie, ere indbyrdes ligestore, forstaaes af sig selv. Saalænge vi
have med concrete Størrelser at gjøre, kan Forskjellen og Flerheden
meget vel bestaae med Ligestorheden; men, naar der handles om rene
Størrelser, stiller Sagen sig anderledes. En Flerhed af Størrelser
forudsætter med Nødvendighed, at den ene maa kunne adskilles fra den
anden: uden indbyrdes Forskjel ingen Flerhed. Men hvorved skulle rene
Størrelser vel adskille sig fra hverandre indbyrdes uden ved
Storheden? Ere Størrelserne indbyrdes ligestore, saa er jo ethvert
Glimt af Forskjel forsvunden, dermed er tillige Flerheden forsvunden:
at flere Størrelser ere indbyrdes lige, vil altsaa med andre Ord sige,
at det ikke er flere, men een og samme Størrelse, som, hvormange Gange
den end gjentages, dog ikke mangfoldiggjøres, men vedbliver at være
den selvsamme.

Vil „den sunde Forstand“ afvise denne Conseqvens, fordi den har noget
Anstødeligt ved sig, da kan det dog vel ikke uden videre skee ved et
Magtsprog; den forstandige Afviisning maa lægge an paa at eftervise,
hvori Misligheden bestaaer, og hvorledes den hæves. Kan een og samme
Størrelse ved at nævnes tyve Gange blive til tyve indbyrdes ligestore
Størrelser, saa maa vel een og samme Specie ved at nævnes tyve Gange
ogsaa kunne blive til tyve Specier; men dersom det ikke er nok at
gjentage den samme Størrelse, dersom der under Gjentagelsen tillige
skal tænkes paa en Forskjel, hvorfra skal da Forskjellen komme?

\textbf{Øieblikket.} Dersom det gamle Spørgsmaal: naar døde Sokrates?
besvares ved at henvise til hans Dødsaar, vil den sunde Forstand neppe
heri ane nogen Modsigelse. Det skulde da være, at Sokrates forsaavidt
ikke har noget Dødsaar, som han ikke har været et Aar om at døe: han
var jo frisk og sund endnu det hele sidste Aar, indtil han drak af
Bægeret. Men naar man nu, ifølge en vel bekjendt og hjemlet
Sprogbrug, ved hans Dødsaar betegner det Aar, hvori hans
% --- p. 4 ---
Dødsdag faldt, saa er her jo ingen Modsigelse. Vil man indvende, at
han heller ikke var en heel Dag om at døe, nu, saa tænke man paa hans
Dødstime, og, vil man endelig gaae til det Yderste, paa hans
Dødsøieblik. Altsaa: naar døde Sokrates? „I det Øieblik, han ophørte
at leve!“ Men dette Svar, der kun siger os, at han døde, da han døde,
er jo intet Svar. Spørgsmaalet gjælder det Tidspunkt i Sokratesses
Tilværelse, da hans Død indtraf. For at angive dette Tidspunkt, kan
man nu tænke sig fire Tilfælde: enten døde han, α) imedens han endnu
levede, eller β) da han allerede var død, eller γ) da han tildeels var
levende og tildeels død, eller δ) da han hverken var levende eller
død. Men α) at døe, imedens man endnu lever, at være levende død, er
dog vel en Modsigelse. At døe, naar man β) allerede er død, er da vel
ogsaa en Modsigelse. Det Øieblik, da Sokrates døde, maa altsaa have
været det Øieblik, da han, liggende paa Grændsen mellem Liv og Død,
befandt sig i en saadan Tilstand, at en Deel af ham, t. Ex. hans
Fødder, allerede var død, medens en anden Deel endnu var levende. Ved
denne Forklaring faae vi imidlertid ikke Spørgsmaalet besvaret, men
kun skudt længere tilbage; thi Spørgsmaalet er: naar døde den allerede
døde Deel? Svares hertil: den allerede døde Deel døde, da den endnu
tildeels var levende, tildeels død: saa vil den i Tilfældet γ) skjulte
Modsigelse være iøinefaldende. At tænke sig Sokrates i en Tilstand,
da han hverken var levende eller død, er da endelig ogsaa det Samme
som at tænke sig en Modsigelse (δ); thi hvorlænge lever vel et
Levende? Indtil det døer!

\emph{„En Philosophie, der er i aabenbar Modsigelse med „den sunde
Forstand“, er falsk; en Philosophie, der umiddelbart retter sig efter
den sunde Forstand, er overfladisk. Philosophien vil ikke vilkaarlig
modsige den naturlige Forstand, men kun vise, at den saakaldte sunde
Sands, uden at have
% --- p. 5 ---
nogen Anelse derom, uophørlig modsiger sig selv.“} (Propæd. S. 41).

Forsaavidt de forskjelligartede Modsigelser udvikle sig under
Tænkningens Gang, og deres Oprindelse sees fra den subjective Side,
maa Reflexionens psychologiske Andeel i de for „den sunde Sands“
anstødelige Modsigelser nærmere bestemmes. De herhen hørende
Udskeielser lade sig i Korthed charakterisere som de \emph{eristiske},
de \emph{spidsfindige} og de \emph{sophistiske}.

Eristikens paafaldende Kjendemærke er eensidigt Hang til Kiv og Strid,
Lyst til overalt at glimre ved Indvendinger, gjøre Vanskeligheder,
nedrive det almindeligt Antagne og haardnakket forsvare egne afvigende
Paastande. Hvad den gamle Philosophie angaaer, er det især den med
den eleatiske beslægtede megariske Skole, der har faaet Navn af
eristisk; men det kan heller ikke negtes,
% (a small stray mark stands between "kan" and "heller" in the print —
%  a speck at hyphen height, read by both OCR witnesses as a hyphen;
%  grammatically no hyphen is possible here)
at enhver eiendommelig udviklet Form af Philosophien har sit polemiske
Element, og at Eristiken, trods alle sine Skyggesider, forsaavidt er
af umiskjendelig Betydning i Videnskabens Udvikling. Til Indsigt og
Viden er det nemlig ikke nok, at man danner sig, hvad man saadan
kalder en fornuftig Mening om Tingene, og ved at forene Sandseindtryk
med naturlig Dømmekraft ledes til at gjøre Forskjel imellem rigtige og
urigtige, antagelige og forkastelige Meninger; ved Hjælp af Meninger
naaer man kun det Sandsynlige, ikke det Sande. Den, der vil hæve sig
fra den blotte Menen til Viden, maa væbne sig mod Sandsebedrag, lade
Tænkningen prøve det Antagne, og kun erkjende det for sandt, der lader
sig tænke, men derimod, uden Hensyn til Mængdens iøvrigt uforgribelige
Meninger, erklære alt det, der ikke lader sig tænke, for usandt. At
nu Philosophien paa disse Vilkaar maa komme i mangesidig Conflict med
den sunde Menneskeforstand, er indlysende. (Jvfr. Philos.\ og
Mathem. S. 12. Anm. Mathem.\ og Dialektik S. 34--35. S. 116--117.
S. 131--32 o.s.v.).
% --- p. 6 ---
Til stringent Begrebsanalyse udfordres Skarpsindighed; men den
eensidige, ufrugtbare, ved subtile, men dog kun halvsande eller endog
usande, Distinctioner blændende Skarpsindighed giver Spidsfindighed.
Saa let som det er at angive Forskjellen imellem Skarpsindighed og
Spidsfindighed, naar man blot holder sig ved det Almindelige, saa
umuligt er det at opstille en for alle enkelte Tilfælde gjældende
Regel, da een og samme Yttring kan, naar den sees fra en Side, tage
sig ud som spidsfindig, og dog, fra en anden Side seet, bære Præg af
sund og sand Skarpsindighed. Den i „Løgneren“ anlagte Modsigelse kan
tjene til Exempel. Er Opgaven nemlig at indøve den Sandhed, at man
ikke uden videre skal indlade sig paa at besvare ethvert Spørgsmaal,
fordi der gives Spørgsmaal, hvori den tilsyneladende tydelige og let
fattelige Mening, nærmere beseet, er meningsløs, kan „Løgneren“ staae
som et Mønster paa sand Skarpsindighed; gaaer man derimod troskyldig
ind paa den forlangte Besvarelse, og det endog saaledes, at man, idet
man med Forundring seer Bekræftelsen forvandle sig til Benegtelse, og
Benegtelsen igjen til Bekræftelse, tilsidst indbilder sig at staae
foran en dyb Gaade, en saare indviklet Knude, hvis Løsning maaskee vil
udbrede et hidtil ukjendt Lys over „de evige Sandheder“, da har
Spidsfindigheden overlistet den sunde Forstand og hildet den i sit
Væv.

I hemmeligt Forbund med Spidsfindigheden og Haarkløveriet finde vi
Sophisteriet. Sophisten nøies ikke med at paavise Modsigelser, han
fremkunstler Modsigelser, han leger med Modsigelser, uddrager
Slutninger, og blænder ved falske Slutninger. Et er det nemlig at
udhæve Modsigelser, et Andet at misbruge dem i Sophistikens Tjeneste;
hvor stor Forskjellen er, kan t. Ex. sees af, hvad ovenfor er anført
om „Øieblikket“. At gjøre opmærksom paa de Modsigelser, Tanken med
Nødvendighed maa støde paa, naar den vil gjøre sig Rede for
Øieblikkets Betydning, kan ingen Forstandig ansee
% --- p. 7 ---
for Sophisterie; først, naar Talen er om Modsigelsernes Anvendelse,
bliver Sophisteriet kjendeligt paa det i Slutningerne Falske,
Uholdbare og Blændende. Vilde saaledes Nogen fra de i Tanken om
Dødsøieblikket indeholdte Modsigelser slutte til Dødens Utænkelighed,
fra Dødens Utænkelighed til Dødens Umulighed, og derpaa grunde et nyt,
med eiendommelig Skarpsindighed udstyret Beviis for Sjælens
Udødelighed, maatte vi ansee ham for en Sophist.

% (section rule as printed)
\section*{B.\\ Mulighedslære.}
\addcontentsline{toc}{section}{B. Mulighedslære}

Skal Philosophien som Modsigelseslære ikke opløse sig i Spidsfindighed
og Sophisteri, maa den altid fra en eller anden Side see sig istand
til at magte Modsigelserne; den maa da kunne bestemme deres Forhold
saavel til Tingene som til vore Reflexioner over Tingene, saavel til
den objective Væren som til den subjective Tænkning; den maa, idet den
gjennemskuer Modsigelsens Væsen, opdage Kjendemærkerne paa det i sig
Tænkelige og det i sig Utænkelige, og derved, saa at sige, aftvinge
Modsigelsen en paalidelig Maalestok for Antageligt og Forkasteligt,
Rimeligt og Urimeligt, Muligt og Umuligt: en grundig Modsigelseslære
er uadskillelig fra en \emph{almindelig Mulighedslære}.

\textbf{Hvad er altsaa Philosophie?} „En Lære om det Mulige.“ Men, da
Philosophien jo dog umuligt kan være en Lære om alt Muligt, er denne
Wolfske Definition kun tilfredsstillende paa det Vilkaar, at
Mulighedens Begreb, dets Omfang, dets Indhold og Forhold til
Modsigelsen nøiagtig bestemmes. Den Grundsætning: Muligt er Alt, som
ikke indeholder en Modsigelse (\textit{possibile est, quod non
implicat contradictionem}), hvilket saa igjen forudsætter, at Alt,
hvad der indeholder en
% --- p. 8 ---
Modsigelse, er umuligt, vil, ifølge vort Foregaaende, neppe afgive
nogen paalidelig enten dogmatisk eller kritisk Rettesnor.

For det Første er det af sig selv indlysende, at den opstillede
Grundsætning kun angiver Muligheden formelt, og altsaa er uskikket til
at orientere os angaaende de reale Muligheder. Selv om vi indtil
videre erklære Alt, hvad der indeholder en Modsigelse, for umuligt, og
Alt, hvad der ikke indeholder en Modsigelse, for muligt, vil det
herved givne Mulighedsbegreb dog være saa negativt, saa tomt og
svævende, at det aldeles ikke lader sig bringe i nogensomhelst positiv
Sammenhæng med det Virkelige. For at bestemme de i Virkeligheden
gjældende Muligheder, er det nødvendigt at kjende de virkelige Ting
med deres Betingelser og Love; men til Kundskab om virkelige Ting
naaer man kun ad Erfaringens Vei. Uden Erfaringens Correctiver vil
den dogmatiserende Tænkning med sine tomme Begreber ligesaa let være
fristet til at sætte de reale Muligheder for snævre Grændser, som til
at indrømme dem et for vidt Spillerum. Eller hvem skulde vel endnu
paa hiin Tid, da Wolf fastsatte Grændserne, og definerede
Mulighederne, have indseet Muligheden af, at Skibe skulde kunne gaae
baade imod Vind og Strøm, at Vogne uden Heste skulde kunne løbe, ikke
som Heste, nei hurtigt som flyvende Fugle, eller Muligheden af, at en
Mand i London skulde kunne correspondere med en Mand i Calcutta saa
hurtigt og let, som om de talede sammen. Saalænge Philosophien med
sine Formalbegreber og Definitioner gjør Forsøg paa at udtømme
Tilværelsen, og opkaste sig til Dommer over det i Virkeligheden Mulige
og Umulige, saalænge vil ogsaa den sunde Menneskeforstand faae
mangfoldig Anledning til triumpherende Anvendelse af Hamlets Ord:
\textit{There are more things in heaven and earth, Horatio, than are
dreamt of in your philosophy.}

Men selv om Philosophien opgiver at ville udgrunde de reale
Muligheders positive, empiriske Indhold, kan den
% --- p. 9 ---
alligevel bidrage væsentlig til at rense og forhøie vor Erkjendelse af
det Mulige, dersom den blot kan sikkre os Mulighedernes negative Side
ved at bestemme det Umulige. Lader nu det Umulige sig kun erkjende
som det i sig Utænkelige, og maa det Utænkelige nødvendigviis altid
være kjendeligt paa Modsigelsen, saa paatrænger sig af sig selv det
Grundspørgsmaal, \emph{hvorledes det Modsigende da forholder sig til
det Tænkelige}; her staae vi ved et afgjørende Vendepunkt.

Idet den dogmatiske Philosophie uden videre forudsætter Eenhed af
Tænken og Væren, gaaer den med det Samme ud fra den Antagelse, at det
Modsigende ubetinget er et Utænkeligt, og det Utænkelige et Umuligt.
Men er det Modsigende nu ogsaa i enhver Henseende utænkeligt?

\emph{„Hvad betyder Modsigelsen? Et tilværende Intet er en
Modsigelse; en rund Fiirkant er en Modsigelse, en flyvende Hest er
ogsaa en Modsigelse. Ikke desto mindre ere disse Modsigelser af
forskjellig Art.“} (Propæd. S. 43).

Dersom det i sig selv Modsigende uden videre er utænkeligt, hvorledes
kommer Tanken da til at opfatte dette Utænkelige? hvordan gaaer det
til, at vi ikke blot forstaae, hvad der menes med Ordet: Modsigelse,
men ogsaa klart indsee, hvad der ligger i Modsigelsens Væsen? Sætte
vi t. Ex. „Noget“ imod „Intet“, „det Værende“ imod „det Ikke-Værende“,
da falder det jo Tanken ligesaa let at fatte det Negative, som at
fatte det Positive; stille vi Sætningen: „det Værende er“, imod
Sætningen: „det Værende er ikke“, da mærker man strax, at en Strid
mellem Subject og Prædikat er Tanken ligesaa indlysende, som en
Overeensstemmelse. Men naar nu Tanken ligesaa let opfatter Striden
som Overeensstemmelsen, ligesaa klart gjennemskuer et tilværende Intet
som et tilværende Noget, og altsaa forsaavidt finder hiint ligesaa
tænkeligt som
% --- p. 10 ---
dette, hvorledes kan det i sig Stridige og Modsigende da for Tanken
vise sig som et Utænkeligt, og dette Utænkelige falde sammen med det
Umulige?

For at rede sig ud af Forlegenheden forsøge man engang at lade
Forestillingen understøtte Tænkningen. Den rene, abstracte, fra al
Anskuelse og Forestilling løsrevne Tænken fatter vistnok ligesaa let
det Negative som det Positive, eftersom det Negative netop er
Abstractionens egen abstracte Frembringelse; men naar Anskuelsens Lys
opgaaer, naar Forestillingen træder til, bliver Modsætningen mellem
det Positive og det Negative, mellem Noget og Intet iøinefaldende; thi
Noget lader sig anskue, Intet, det rene Intet, lader sig ikke anskue;
til det Positive svarer en Forestilling, men Intet, det Negative,
frastøder enhver Forestilling, og bliver desaarsag igjen frastødt og
forkastet af Forestillingen.

Hvorledes Tænkningen ved at forbindes med Forestillingen kommer til at
opfatte det Modsigende som det Utænkelige, og det Utænkelige som det
Umulige, kan sees af Exemplet med den runde Fiirkant. Den Sætning:
Fiirkanten er rund, lader sig vel tænke, forsaavidt Prædikatets Strid
med Subjectet lader sig tænke; men den runde Fiirkant lader sig ikke
forestille, efterdi Forestillingen om det Runde fortrænger og
udsletter Forestillingen om det Fiirkantede. Tanken, der indseer, at
den runde Fiirkant gjør en Fordring til Forestillingen, som
Forestillingen maa forkaste, stempler nu Modsigelsen som noget
Forkasteligt, erklærer den runde Fiirkant for et Utænkeligt, og dette
Utænkelige for et Umuligt. At det netop er Anskuelse og Forestilling,
som i slige Tilfælde maa hjælpe Tanken til at opdage det Utænkelige,
sees ikke blot af Sætninger, i hvis Bestanddele Forestillingen
øiensynlig hersker, saasom: Sort er ikke hvidt, Tykt er ikke tyndt,
men tillige af Exempler, der ikke umiddelbart angive Modsigelsen i Ord
og Udtryksmaade, som naar man spørger: er
% --- p. 11 ---
et regulært Hexaeder muligt? er det muligt, at en med Ellipsen
parallel Linie kan være en Ellipse? o. s. v.

Dersom vi nu blot kunde stole paa, at Anskuelse og Forestilling altid
og i alle Tilfælde ville afgive ubedragelige Kjendemærker paa
Tænkeligt og Utænkeligt, saa havde vi dog her noget Fast at holde os
til; men see nu til Exemplet med den flyvende Hest! Det forekommer
dog vist den naturlige Sands, som om her under Benævnelsen: flyvende
Hest, laa skjult en Modsigelse. Men til hvem skal den sunde
Menneskeforstand nu henvende sig for at faae det Modsigende opklaret?
Til Tænkningen! Det hjælper ikke; den rene Tænkning tænker Alt, ogsaa
det Utænkelige. Til Anskuelsen! Hjælper heller ikke; Anskuelsen
lever hos Digteren, og Digteren rider paa Pegasus, og Pegasus forliger
Anskuelsen med alle bevingede Heste. Til en Forening af Tænkning og
Anskuelse! Men hvorledes kan en Forening af tvende Upaalidelige
afgive noget Paalideligt? Til Sandseindtryk! Ja, Sandseindtryk kunne
i Forening med Tænkningen lære os Noget om det Mulige, forsaavidt de
lære os Noget om det Virkelige, forstaaer sig; men hvad lære de os om
det blot Mulige, og hvad kunne de vel lære os om det i sig Modsigende
og Umulige?

For at bestemme Mulighedens Forhold til Modsigelsen, er det fremfor
Alt nødvendigt at fatte \emph{Modsigelsens Grundsætning}, dens Form,
dens Gyldighed, dens Anvendelse.

I den formelle Logik finde vi Identitetens og Modsigelsens Sætninger
saa discret opstillede, at vi derved ikke blot forsynes med to
Grundsætninger, men ovenikjøbet med to selvstændige Principer. Skulde
det nu virkelig have Noget at betyde med denne Principernes Tohed,
maatte vi vel vente en eller anden Oplysning angaaende deres indbyrdes
principielle Forskjellighed. Men hvori skal Forskjellen egentlig
bestaae? For at forklare Identitetssætningen: \textit{A} er, hvad
\textit{A} er, underskyder man Contradictionssætningen, at \textit{A}
netop er,
% --- p. 12 ---
hvad \textit{A} er, fordi \textit{A} ikke er Andet, og ikke kan være
Andet, end hvad \textit{A} er. Har nu Contradictionsprincipet, og det
endog, i logisk Forstand, førend det blev født, maattet gjøre Tjeneste
ved Opstillingen af den rene Identitet, saa er det heller ikke mere
end billigt, at Identitetsprincipet, efter at det allerede er født, og
de logiske Fødselsveer lykkelig ere overstaaede, igjen forhjælper
Contradictionsprincipet til behørig Anerkjendelse. Hvorfor er saa
ingen Ting det, hvad den ikke er? Svar: fordi enhver Ting er, hvad
den er. Og hvorfor er enhver Ting, hvad den er? Svar: fordi ingen
Ting er, hvad den ikke er. Med andre Ord: naar her spørges om
Identiteten, svarer man med Contradictionen, og naar her spørges om
Contradictionen, svarer man med Identiteten. Seer man nøiere til, maa
det unegtelig forekomme den sunde Menneskeforstand, som om her
egentlig ikke var to, men kun een Grundsætning.

Dog sæt endog, at vi, for saavidt muligt at undgaae Conflict med
traditionelle Former, lode den formelle Logik beholde sine to
Grundsætninger ubeskaarede, saa at vi --- for „Vishedens“ og
„Sandhedens“ og „Determinationens“, „Statueringens“, „Hvadhedens“,
„Decisionens“, „Eenstemmighedens“ og „Samstemmighedens“ Skyld ---
bogstavelig fastholdt Identiteten for sig og Contradictionen for sig;
mon en saadan Fastholden for sig saa maaskee skulde kunne forhjælpe
den sunde Forstand til klar og tydelig Indsigt i Forholdet mellem det
Modsigende og det Umulige? Det synes ved første Øiekast, som var her
endnu nogen Vanskelighed. Har nemlig ethvert \textit{A}, naar det
blot er, hvad det er, ifølge Identitetsprincipet Krav paa „et gyldigt
er“: da lader det næsten, som om det Meningsløse og Tankeløse ogsaa
kunde gjøre Krav, gyldigt Krav paa „et gyldigt er“; kan enhver Ting,
naar den bare ikke er, hvad den ikke er, tilfredsstille
Eenstemmigheds- og Samstemmighedsprincipet, hvorfor skulde det
Stridige, det Urimelige, det i sig Modsigende og Taabelige
% --- p. 13 ---
da ikke ogsaa kunne tilfredsstille Eenstemmigheds- og
Samstemmighedsprincipet?

Men at oversee eller endog forkaste Modsigelsens Grundsætning gaaer da
heller ikke paa nogen Maade an; tvertimod, naar Tanken forsøger at
negte Contradictionsprincipets Gyldighed, bliver den netop gjennem
Negtelsen ledet til at bekræfte den. At sætte: „Sort er hvidt“,
istedetfor: „Sort er ikke hvidt“, er ikke alene muligt, men fornuftigt
og nødvendigt, naar den tilsvarende Modsigelse ret skal
tydeliggjøres; men, idet Modsigelsen paa denne Maade er stemplet som
Modsigelse, er Grundsætningen hævdet. Det Utænkelige er i denne
Forstand ikke et Noget, Tanken slet ikke kan tænke, men det Utænkelige
er den Tanke, Tænkningen, ifølge et i en given Sammenhæng givet
Formaal, selv fordømmer, en Tanke, der kun bliver til, for at
forsvinde, en Tanke, ved hvis Ophævelse Tænkningen erkjender sin Lov,
og tilfredsstiller sig selv. Kun naar det Utænkelige saaledes bestemt
er sat imod det Tænkelige, da først kan dette siges at være det
Muliges, hiint det Umuliges logiske Mærke. Umiddelbart kan altsaa
Modsigelsen, logisk seet, ikke afskaffes. Kunde Modsigelsens
Grundsætning een Gang for alle afskaffe Modsigelsen, vilde den med det
Samme gjøre sig selv overflødig. Men naar Tanken ikke kan røre sig,
uden at Modsigelsen, dens Braad og Spore ogsaa rører sig, staaer
Grundsætningen ved Magt, fordi Modsigelsens Medvirkning staaer ved
Magt; for at bestemme det Tænkelige, sigte det Mulige og afvise det
Umulige, maa Contradictionsprincipet uophørlig modsige Modsigelsen,
maa ved Modsigelsens Opløsning hævde Modsigelsen.

I Anvendelsen af Contradictionsprincipet sees Forskjellen mellem den
kritiske og den dogmatiske Philosophie. Enig med Dogmatismen, naar
det gjælder om at stille Modsigelsen udvortes fjendtlig mod Tanken og
Tankebevægelsen, afviger den kritiske Philosophie saa fra den
dogmatiske deri, at den i det Utænkelige kun seer det Umuliges
subjective, ikke tillige dets
% --- p. 14 ---
objective Side. Naar Kant forsøger at tænke sig Verdens Oprindelse,
støder han ganske rigtigt paa Modsigelsen; det Samme gjentager sig,
naar han vil tænke Substantialitetens, Nødvendighedens og Frihedens
Væsen objectivt; den skarpsindige, fra sit Synspunkt conseqvente,
Tænker er her saa langt fra at blive det Modsigende qvit, at han
tvertimod, idet han seer Modsigelsen knytte sine Knuder, og sætte sig
fast i uafviselige Antinomier (Propæd. S. 29--30; Mathem.\ og
Dialekt. S. 125--126), erklærer den for uovervindelig. Verdens
objective Væsen hører altsaa, ifølge Kant, til det for os Utænkelige;
men hvad slutter nu Kant heraf? Ingenlunde, at en objectiv Verdens
Realitet skulde af den Grund ansees for umulig --- det er en Slutning,
der først opkommer hos Fichte (Propæd. S. 33.) --- nei, Forfatteren af
„Kritik der reinen Vernunft“ advarer just imod overilede Slutninger
fra det subjectivt Utænkelige til det objectivt Umulige. Adskillelsen
af Tænken og Væren, den kritiske Philosophies oprindelige
Charakteermærke, grunder sig jo netop paa den Forudsætning, at der for
den menneskelige Forstand gives væsentlige, nødvendige, uopløselige
Modsigelser.

Men ere de kantiske Antinomier da virkelig uopløselige? Det kommer
vel an paa, om man vælger at indrømme Modsigelsen enten en blot
subjectiv eller en baade subjectiv og objectiv Betydning. Et saadant
Valg kan naturligviis ikke være vilkaarligt, det maa fremgaae af
Sagens Natur. Dersom det da nu viser sig, at man, ved at negte
Modsigelsens objective Gyldighed, fastholde den kritiske Adskillelse
af Tænken og Væren, og saaledes afvise ethvert Forsøg paa at erkjende
„\emph{Tingen i sig}“, ingenlunde undgaaer Modsigelsen, men meget mere
indvikler sig i nye og større Modsigelser: hvad saa?

% (short centred rule after this paragraph, closing the discussion as
%  printed at the foot of p. 14; p. 15 opens the next section)

% I. Philosophiens Begreb continues; C. Dialektik region.
% --- p. 15 ---
\setcounter{footnote}{0}
\section*{C.\\ Dialektik.}
\addcontentsline{toc}{section}{C. Dialektik.}

% FOOTNOTE MARK: first footnote of the printing (below) is marked *) both
% in the text and at the note, which is set off by a short rule; the note
% continues at the foot of p. 16 without a repeated mark. Plain \footnote{}
% used pending a decision on \thefootnote.

Hvad er saa Philosophie? „En Tvetydighedslære, en Videnskab om det
Tvetydige“. Rigtigheden af denne Definition synes idetmindste at finde
rigelig Stadfæstelse i vort Foregaaende; thi endskjøndt vi vel skjelnede
imellem Skarpsindighed og Spidsfindighed, maatte vi dog tilstaae, at det
Samme kunde opfattes baade som spidsfindigt og ikke spidsfindigt; da der
handledes om en Adskillelse af Tænkeligt og Utænkeligt, fik vi ud, at det
Utænkelige var Noget, der baade kunde tænkes og dog ikke tænkes; da det
endelig gjaldt om en Modsigelsens Løsning, endte det med, at Modsigelsen
baade kunde løses og dog ikke løses. Vil den sunde Forstand finde en
saadan Forkjærlighed for det Tvetydige mistænkelig, da lægge man Mærke
til, at Udtrykket: „Tvetydighedslære“ selv er tvetydigt. Medens høiere
Anviisning til at gjøre det Klare dunkelt, det Tydelige forblommet, det
Forstandige uforstaaeligt altid bliver at søge hos Sophisterne, maa
Philosophien see at sætte det Tvetydige i et andet Lys, et Lys, der
bringer Tvetydigheden til at forsvinde. Som den Videnskab, der overalt i
det Tvetydige opdager det Tvesidige, og ved Hjælp af Modsigelsen klarer
Fordoblingen, er Philosophien \emph{Dialektik}\footnote{At Philosophiens
Form i sit Væsen er dialektisk, maa ikke misforstaaes, som om
Begrebsudviklingen overalt skulde bære de formelle Antinomier, hvorved
logiske Overgange motiveres, umiddelbart til Skue; tvertimod! den
philosophiske Fremstilling kan og bør ogsaa optage refererende,
beskrivende, kritiserende Bestanddele. Kun dette er afgiørende, at de
Elementer, som meddeles i Form af Beretning, Beskrivelse, kritisk
Drøftelse o.\ dsl., overalt vise sig, ifølge Anlæget og den herskende
Tankestrøm, at tjene det Dialektiske. Der er i den hegelske Methode en
formel Overspændthed, en logisk Overanstrængelse, som fornemmelig
fremkaldes ved den Bestræbelse, at lade Fremstillingen paa alle Punkter
brillere ved
% (the footnote text continues at the foot of p. 16:)
dialektiske Lynglimt. Det gaaer med Dialektikens Lynglimt og de
Elementer, hvori de forberedes, og hvoraf de udvikle sig, paa
tilsvarende Maade, som med det naturlige Lyn; vel maa Lynilden i Reglen
forudsætte Skyer, men derfor ere Skyerne dog ikke Lyn.}.

% --- p. 16 ---
Den Sandhed, at Philosophien efter sin høiere Natur er og maa være
Dialektik, er kategorisk udtalt af Plato. (Jvfr.\ Propædeut.\ S.\
124--129; S.\ 150 o.\ fl.). Hos Plato har nu Fremstillingen af det
Dialektiske (διαλεκτική af διαλέγεσθαι), som bekjendt, antaget Samtalens
Form, og hermed staaer i Forbindelse, at Dialektikeren just betegnes som
Den, der i Samtalen forstaaer at spørge og at svare rigtigt (Τὸν δὲ
ἐρωτᾷν καὶ ἀποκρίνεσθαι ἐπιστάμενον ἄλλό τι σὺ καλεῖς ἢ διαλεκτικόν;
% „ἄλλό“ with double accent so printed
\textit{Krat.}\ 390 c.). At det nu ikke er Samtalen, men den ved
„Sprogets Kunst“ betingede Begrebsudvikling, hvorpaa det i Dialektiken
væsentlig kommer an, forstaaer sig af sig selv; men at Samtalens Form,
naar alt Uvedkommende holdes borte, ikke destomindre er en for den
begyndende Dialektiker meget naturlig Form, kan let oplyses ved nogle
Exempler.

\textbf{Noget og Andet.}
\addcontentsline{toc}{subsection}{Noget og Andet}
A. Dersom Kategorierne „Noget“ og „Andet“ ligefrem udelukkede hinanden,
var der ingen Dialektik, faldt de uden videre sammen, var der heller
ingen Dialektik; det Dialektiske viser sig just deri, at de to
Sætninger: „Noget er, hvad det er, ikke et Andet“, og „Noget er, hvad
det ikke er, altsaa et Andet“, ikke blot ophæve, men tillige forudsætte
og begrunde hinanden gjensidig. B. Hviler det Dialektiske paa en saa
slibrig Grund, da er al Dialektik falsk. A. Hvorledes? B. Den
opstillede Paastand er i aabenbar Modsigelse med den sunde
Menneskeforstand, og en \emph{Philosophie, der er i aabenbar Modsigelse
med „den sunde Forstand“, er falsk}. A. Skal, hvad Du her kalder
\emph{aabenbar Modsigelse}, uden videre affærdiges ved et umiddelbart
Skjøn, kan Philosophien ikke tillægge Kjendelsen nogen Betydning; thi en
\emph{Philosophie, der umiddelbart
% --- p. 17 ---
retter sig efter „den sunde Forstand“, er overfladisk}; men vil „den
sunde Forstand“ nedlade sig til en Undersøgelse, kan Undersøgelsen jo
let komme i Gang. B. Det skulde da være den anden Sætning, under hvis
negative Form en eller anden Tvetydighed havde skjult sig, og som
desaarsag kunde blive Gjenstand for Undersøgelse; i den første Sætning,
Identitetssætningen, den rene Vishedssætning (\textit{principium
certitudinis}), kan der ikke være Spor af Tvetydighed. A. Men dersom
det nu alligevel er den første Sætning, hvori Tvetydigheden,
Ubestemtheden, og forsaavidt Uvisheden, ligger skjult, hvad saa? B. Saa
maa det kunne bevises. A. Er det i Logiken aldeles forbudt at udsige
Andet om Noget, end at det er Noget? B. Ingenlunde, thi hvorledes
skulde man ved blot at gjentage: Noget er Noget, nogensinde kunne
bestemme, \emph{hvad} Noget er; derimod gjør Logiken, der ikke ynder
Tvetydigheder, en bestemt, ligesaa let fattelig som umiskjendelig,
Forskjel imellem det, at bestemme Noget ved Hjælp af et Andet, og det,
at sige Andet om Noget, end hvad Noget er. A. Det er altsaa tilladt at
bestemme, hvad Noget er, ved at bestemme, hvad Noget ikke er? B. Ikke
alene tilladt, men nødvendigt. A. Det er altsaa nødvendigt, at der i
Bestemmelsen af Noget maa indeholdes en Bestemmelse af Andet. B. Ogsaa
dette kan vel indrømmes, dog paa den udtrykkelige Betingelse, at Noget
netop ved Bestemmelsen udelukker Andet. A. Noget er altsaa? B. Et fra
Andet Forskjelligt. A. Og naar vi nu, ifølge Betingelsen, udelukke
Andet, er Noget? B. Et Forskjelligt. A. Men nu er da ogsaa Andet paa
sin Side? B. Et fra Noget Forskjelligt. A. Og naar vi, ligeledes
ifølge Betingelsen, lade Andet udelukke Noget, sige vi at Andet er?
B. Et Forskjelligt. A. Altsaa: Noget er, naar det bestemmes, et
Forskjelligt, Andet er, naar det bestemmes, et Forskjelligt; Noget er
altsaa, hvad Andet er, nemlig et Forskjelligt. B. Denne Tvetydighed er
da let at gjennemskue. Vistnok
% --- p. 18 ---
er Noget et Forskjelligt ligesaavelsom Andet, men Forskjelligt og
Forskjelligt er to Ting, Noget er jo netop et Forskjelligt ved at være
forskjelligt fra Andet, og Andet et Forskjelligt ved at være
forskjelligt fra Noget: Noget og Andet ere altsaa endnu, som før,
indbyrdes forskjellige. A. Ganske vist ere Noget og Andet indbyrdes
forskjellige, kun en Sophist kunde falde paa at ville negte deres
Forskjel. Hvad her skulde bevises, var kun dette, at deres Forskjel er
begrundet i deres Identitet. B. Hvad der er forskjelligt, er ikke
identisk, hvad der er identisk, er ikke forskjelligt: kun ved rene
Kategorier kan man sikkre sig saavel imod skjulte Tvetydigheder som imod
aabenbare Modsigelser.

A. Men, hvis jeg ikke feiler, har Du jo allerede selv ved dine bestemte
og tydelige Svar givet værdifulde Bidrag til at skjærpe vort Blik for
den dialektiske Eenhed af Noget og Andet. B. Det skulde jeg meget
betvivle. A. Du har jo dog fremhævet, at Noget kun lader sig bestemme
som et Forskjelligt, og, for Tydeligheds Skyld, som et fra Andet
Forskjelligt? B. Ja, det har jeg fremhævet. A. Du har da med det
Samme udtrykkelig fordret, at Noget ifølge Forskjelligheden maa opfattes
som et Andet end det, hvorfra det er forskjelligt. B. Naturligviis.
A. Noget er altsaa i dine Tanker noget Andet end Andet, i Modsætning
til Andet, altsaa dog et Andet. B. Dit Ordspil fører til Intet; thi,
naar Noget opfattes som Andet, bliver det jo et andet Andet, nemlig det
Andet, som Noget selv er, i Modsætning til alt Andet, som Noget ikke
er. A. Det kan jeg forstaae; nu taler Du som en sand Dialektiker. Du
har da nu altsaa selv indseet, at der med Hensyn paa den i Noget
indeholdte Bestemmelse gives to Slags Andet, nemlig et Andet, som Noget
selv er, i Modsætning til et andet Andet, som Noget ikke er; at, med
andre Ord, Noget kun kan udelukke Andet ved selv at være et Andet.
B. Ja, saalænge man holder sig til saa almindelige Udtryk som Noget og
Andet,
% --- p. 19 ---
kan man jo vel uden stor Vanskelighed indøve den Tankeleg, at gjøre
Noget til Andet og Andet til Noget; men lader os prøve det med noget
Bestemt! Viis mig i bestemte Exempler, at Noget virkelig er identisk
med Andet, og jeg skal indrømme Dig, at der er „sund Forstand“ i det
Dialektiske. A. Du anseer da et Blad i en Bog for noget Bestemt?
B. Forstaaer sig. A. De to Sider af et Blad ere da vel ogsaa hinanden
saa modsatte, at naar vi kalde Forsiden et Noget, maae vi kalde Bagsiden
et Andet. B. Fortræffeligt! Nu falder det da vel heller Ingen ind at
identificere Forsiden med Bagsiden, og saa paastaae, at Forsiden er,
hvad Bagsiden er, fordi de begge ere Sider, eller fordi Forsiden, naar
Bladet vendes, bliver Bagside, og Bagsiden Forside. A. Men Identiteten
er da vel ogsaa ved denne Leilighed ligesaa iøinefaldende som
Forskjellen? B. Jeg kan ingen Identitet see. A. Hvad mener Du da om
Bladet? B. At Bladet, der har de to Sider, netop bestaaer i Sidernes
Forskjel og Modsætning! thi dersom de to Sider forvandledes til een
Side, fik vi jo kun en Overflade, intet Blad. A. Og dersom de to Sider
saaledes adskiltes, at de derved kom til at tilhøre hver sit Legeme, fik
vi to Overflader, men intet Blad. B. Men det gjør dog en væsentlig
Forskjel, om jeg siger: Forsiden er i Bladet forenet med Bagsiden, eller
jeg siger: Forsiden er, hvad Bagsiden er, Forsiden er Bagside. Det
Første kan jeg indrømme, det Sidste ikke. A. Dersom Bladet, der
forener de to Sider, var et tredie for sig ligeoverfor Siderne
Bestaaende, vare vi vistnok ikke berettigede til at opfatte det som
Sidernes tilværende Eenhed; men da Bladet selv i dets Adskillelse fra
Siderne aldeles ikke er til, saa maa den ene Side jo netop i Bladet vise
sig identisk med den anden. B. Kan man saa maaskee heraf udlede, at
Forsiden selv baade er Forside og Bagside, og Bagsiden igjen baade
Bagside og Forside? A. Det er en Conseqvens, Du neppe vil afvise.
Kalde vi de to Sider α og β, er det jo
% --- p. 20 ---
indlysende, at α, som i og for sig hverken er For- eller Bagside, ved
sin Modsætning til β maa være Forside, forsaavidt β er Bagside, og
omvendt; α er altsaa ifølge sin Identitet med β baade Forside og
Bagside. B. Det er en eiendommelig Udvidelse af Begrebet: Identitet.
Hvad jeg imidlertid har at indvende imod det saaledes udvidede
Identitetsbegreb, vil jeg dog helst tilbageholde, indtil Du har fremsat
et Exempel endnu af en anden Beskaffenhed.

A. Hvad mener Du, om vi kaldte det Større et Noget og det Mindre et
Andet? B. Dette Exempel finder jeg meget passende. A. Og hvad er det
saa, vi skulle bevise? B. Hvad Andet vel, end at Noget baade er et
Større og et Mindre, Andet baade et Mindre og et Større, og en og samme
Størrelse selv baade mindre og større. Og Intet er i Sandhed lettere
end at bevise alt dette. En flygtig Sammenligning af tre Størrelser:
$\alpha > \beta > \gamma$ siger os jo strax, at α er et Større,
forsaavidt β er et Mindre, og β et Større, forsaavidt γ er et Mindre.
Men hvad følger nu heraf? At β i en Henseende er et Mindre, i en anden
Henseende et Større, deri er ingen Modsigelse; men dersom Nogen
paastaaer, at Størrelsen β er paa een Gang og i een og samme Henseende
baade større og mindre end sig selv, kommer han dog nok i Modsigelse,
ikke blot med den naturlige Forstand, men med den sande Logik. At een
og samme Side af samme Blad kan i et Øieblik være Forside, i et andet
Bagside, og i eet og samme Øieblik, naar det, vel at mærke, betragtes
fra to modsatte Synspunkter, baade Forside og Bagside, indeholder ingen
Modsigelse, men en Stadfæstelse af den gamle Regel: \textit{diversi
respectus tollunt omnem contradictionem}. Ignorerer man derimod
\textit{diversi respectus}, overseer man, hvad dog allerede Plato og
Aristoteles, trods deres iøvrigt forskjellige Aandsretning, begge saa
strengt have indskjærpet, at det i slige Undersøgelser Omspurgte hver
Gang maa sees at være τὸ αὐτό, og betragtes: κατὰ τὸ αὐτό,
% --- p. 21 ---
πρὸς τὸ αὐτό, ἐν τῷ αὐτῷ χρόνῳ: saa kan man jo vel ved pikante
Identitetsvendinger og overraskende Modsigelser for et Øieblik blænde de
Uerfarne, men man vil aldrig overbevise den Kyndige, aldrig
tilfredsstille den forstandige, efter Sandhed forskende Tænkning.
A. Vigtige, vægtige Ord! Saa har da nu din „sunde Forstand“
næstendeels faaet Bugt med Modsigelsen og gjort det af med Dialektiken;
jeg siger: \emph{næstendeels}, thi et lille Punkt er endnu tilbage; det
er et Punkt, vi endelig maae have klaret; Du vil vist let kunne klare
det: det angaaer \textit{diversi respectus}, og Du har jo studeret
\textit{diversi respectus}. B. Studeret? Det er aldrig faldet mig
ind, at der i Logikens \textit{diversi respectus} var Noget videre at
studere; jeg har altid antaget, at det var nok, naar man saadan i
Almindelighed bemærkede, at der gives \textit{diversi respectus}.
A. Da er det dog vistnok nødvendigt at gjøre Forskjel paa
\textit{respectus}, siden der nu engang gives \textit{diversi
respectus}. Om samme Blads to Sider, for at vende tilbage til et af
vore Exempler, gjælder unegtelig i en Henseende det Samme, at de nemlig
begge ere Sider, i en anden Henseende det Forskjellige, at den ene er
Forside, den anden Bagside; i en Henseende det Samme, at de begge
t.\ Ex.\ ere beskrevne, i en anden Henseende det Forskjellige, at den
ene t.\ Ex.\ er hvid, den anden rød o.\ s.\ v.\ Ethvert af de
forskjellige Hensyn kan igjen foranledige nye Hensyn; thi, uagtet der om
begge Sider gjælder i en Henseende det Samme, at de begge ere beskrevne,
saa gjælder dog ogsaa, hvad dette Punkt angaaer i en anden Henseende det
Forskjellige, at den ene nemlig er beskreven med sort Blæk, den anden
med rødt Blæk. Disse Hensynsfordoblinger kunne fortsættes i det
Ubestemmelige. Skal nu Tænkningen ikke tabe sig i et uoverskueligt
Virvar af allehaande
% doubtful: a mark above the n makes both OCR witnesses read
% „allehaaude“ here (the parallel word two lines below is plainly
% „allehaande“); the glyph shape matches n and the mark is taken for an
% ink speck — rejected reading: allehaaude
Hensyn til allehaande Henseender og Bemærkninger og Betænkninger og
Betragtninger, maa der i Logiken nødvendigviis gjøres Forskjel, ikke
blot imellem væsentlige og uvæsentlige Hensyn, men
% --- p. 22 ---
tillige mellem saadanne Henseender, der falde ligegyldigt ud fra
hinanden, og saadanne, som indeholdes i hinanden, og desaarsag med indre
Nødvendighed følge af hinanden. At den ene Side af et Blad er
beskreven, medfører naturligviis ikke, at den anden ogsaa maa være det;
man kan tage Hensyn til det Første, uden at tage Hensyn til det Sidste;
det Første kan være, uden at derfor det Sidste er, og der er da heller
Ingen uden Sophister, der kunde falde paa at sætte det Sidste identisk
med det Første. Anderledes forholder det sig derimod med det at være
Forside og det at være Bagside; her er det umuligt at tage Hensyn til
det Første uden at tage Hensyn til det Sidste; thi disse to
Tilværelsesmaader ere saaledes indeholdte i eet og samme Tilværende, at
Hensynet til det Første endog falder sammen med Hensynet til det Sidste:
det Første er trods Forskjellen, ja i Kraft af Forskjellen,
\emph{identisk} med det Sidste. Sand Indsigt i Identitetens,
Forskjellens og Modsigelsens Væsen er betinget af, at man tilstrækkelig
analyserer \textit{diversi respectus}.

B. Men kan der da virkelig mod den Sætning: \textit{A} er, hvad
\textit{A} er, gjøres nogensomhelst begrundet Indvending? A. Blandt
Andet den Indvending, at Sætningen mangler Begrundelse. B. Men, naar
nu Sætningen er Eet med Principet, Identitetsprincipet,
Sandhedsprincipet, Vishedsprincipet, \textit{principium certitudinis},
og selvfølgelig i alle Maader indlysende ved sig selv, behøver den jo
ingen Begrundelse. A. Er det da af Sætningen umiddelbart indlysende,
\emph{at} \textit{A} virkelig er? B. Nei, ikke \emph{at} \textit{A}
virkelig er, men \emph{hvad} \textit{A} er. A. Er det umiddelbart
indlysende, \emph{hvad} \textit{A} da egentlig er? B. Just ikke,
\emph{hvad} \textit{A} egentlig er, men at, \emph{dersom} \textit{A}
overhovedet er, \emph{saa} er \textit{A}, \emph{hvad} \textit{A} er.
A. Men dersom \textit{A} nu slet ikke er? B. Saa er det dog, hvad det
er. A. For at være, hvad det er, behøver \textit{A} altsaa ikke engang
at være? B. Du forstaaer meget godt, hvad jeg mener. A. Her spørges
ikke om Menen, men om Viden;
% --- p. 23 ---
dog, hvad mener Du? B. At selv om \textit{A} er et Ikkeværende, saa er
det dog, hvad det er, nemlig et Ikke-Værende. A. Men det Ikkeværende
er dog vel det Samme, som det vi ellers pleie at kalde Intet, ingen
Ting? B. Naturligviis. A. Altsaa, enten jeg siger: ingen Ting er,
hvad den er, eller jeg siger: enhver Ting er, hvad den er, saa skulde
det strax, naturligviis ifølge Identitetsprincipet, Sandhedsprincipet,
Vishedsprincipet, \textit{principium certitudinis}, være den sunde
Forstand indlysende, at jeg med begge Sætninger mener Eet og det Samme!

B. Men kan der da anføres Nogetsomhelst til Forsvar for en saa
selvmodsigende Sætning som: Noget er, hvad det ikke er? A. Blandt
Andet dette, at en Sætning, der ved sin aabne Modsigelse saa kjækt
udæsker Tænkningen, maa vække Eftertanken og opfordre til Begrundelse,
idet den fordrer en Gjendrivelse. B. Men hvorledes skulde man ved
Grunde gjendrive en Sætning, som umiddelbart gjendriver sig selv?
A. Den Sætning: Noget er, hvad det ikke er, gjendriver ingenlunde sig
selv umiddelbart. B. Dersom det i denne Sætning Utænkelige, Urimelige
og Meningsløse ikke er indlysende ved sig selv, hvad er saa indlysende
ved sig selv, hvad er saa sund Menneskeforstand? A. Den sunde
Menneskeforstand seer ingenlunde det absolut Urimelige i, at Noget
skulde kunne være, hvad det ikke er. B. Altsaa skulde den sunde
Menneskeforstand kunne tænke sig, at \emph{Noget} --- vel at mærke, til
en given Tid og i en ganske bestemt Betydning --- \emph{er, hvad det}
--- til samme Tid og i samme Betydning --- \emph{ikke er}; nei, det er
dog virkelig for galt. A. Ja, det er virkelig for galt; men en saa
urimelig Fordring gjør den Sætning: \emph{Noget er} --- vel at mærke,
naar dets Tilværelse begrundes --- \emph{hvad det ikke er}, --- naar det
eensidig fastholdes i overfladisk, abstract Umiddelbarhed. B. Ja, det
er en anden Sag, naar dit Forsvar beroer paa en Fortolkning.
% --- p. 24 ---
A. Og da vel ogsaa en anden Sag, naar din Gjendrivelse beroer paa en
Fortolkning. B. Saa beviis os da engang ved Hjælp af dine dialektiske
Fortolkninger, at eet og samme Noget kan til een og samme Tid i een og
samme Henseende være baade, hvad det er, og hvad det ikke er, samt at
een og samme Størrelse kan paa een Gang være baade større og mindre end
sig selv. A. Slige Beviser høre under Sophistiken, ikke under
Dialektiken. B. Men til at bevise, at Noget kan i en Henseende være,
hvad det i en anden Henseende ikke er, eller, at een og samme Størrelse
kan i en Relation være et Mindre, i en anden et Større, dertil behøves
kun \textit{diversi respectus}, ingen Dialektik. A. Men til at bevise,
at Noget, skjøndt forskjelligt fra Andet, dog væsentlig er Eet med sit
Andet, fordi begges Tilværen have samme Udspring, samme Rod; til at
bevise, at den samme Størrelse er af een og samme Grund baade et Større
og et Mindre, eftersom de to Relationer, de to Henseender have eet og
samme Punkt, hvori de træffe sammen, nemlig, det at være Størrelse; til
at bevise, at Identiten
% „Identiten“ so printed (for „Identiteten“)
af Noget og Andet, det Større og det Mindre, netop indeholder Grunden
til, at Noget i Tilværelsen er forskjelligt fra Andet, det Større
forskjelligt fra det Mindre, dertil behøves Dialektik; thi dette kan
Dialektiken, men ogsaa kun Dialektiken, bevise.

\textbf{Begreb og Existens.}
\addcontentsline{toc}{subsection}{Begreb og Existens}
C. En dialektisk Udvikling af de aprioriske Begrebers indre Sammenhæng
og Overgang i hinanden er paa Videnskabens nuværende Standpunkt vistnok
i det Hele at anerkjende; men et andet Spørgsmaal er det, om man, fordi
det Dialektiske i Philosophien spiller en betydelig Rolle, nu ogsaa af
den Grund er berettiget til ikke blot at forekaste den formelle Logik,
at den ikke er dialektisk, men endog i Henhold til et platonisk Udtryk,
der dog altid kan tages i mere end een Betydning, at forvandle den hele
Philosophie til lutter Dialektik. A. Havde Du med Opmærksomhed fulgt
den foregaaende Undersøgelse, vilde Du,
% --- p. 25 ---
hvad det første Punkt angaaer, let have overbeviist Dig om, at her ikke
har været tænkt paa at levere en udførlig Kritik af den logiske
Formalisme, men kun paa at afvise den Suffisance, hvormed man ikke
sjældent, under Paaberaabelse af „den sunde Menneskeforstand“, tillader
sig at henregne enhver logisk Indsigelse mod umiddelbare Grundsætningers
umiddelbare Gyldighed til Hjernespind og Sophisterie; hvad det andet
Punkt angaaer, kunne dine Betænkeligheder vel behøve en tydeligere
Udtalelse. C. Aprioriske Begreber, saasom: Noget og Andet, Væsen og
Skin, Grund og Følge o.\ dsl., kunne ved Modsigelsens Anvendelse let
frembringes af hinanden. Den, der sætter Grunden, underforstaaer jo
allerede Følgen, Grunden er altsaa i Grunden baade sig selv og Følgen;
men, naar Grunden først paa Grund af Følgen bliver en egentlig Grund,
saa er jo Grunden Følge, og Følgen Grund. Det Samme er ikke blot i en
Henseende Grund, i en anden Følge, men de to Henseender falde tillige
sammen i een Henseende: at dialektisere Begreberne ud af hinanden, er,
om jeg maa bruge dette Udtryk, det Samme som at dialektisere
\textit{diversi respectus} ind i hinanden. Hvad de aposterioriske
Begreber, Erfaringsbegreberne, derimod angaaer, er det umuligt at
dialektisere Existensen ud af Begrebet; eller hvem vilde vel forsøge at
dialektisere den virkelige Hest ud af Hestens Begreb, eller den
virkelige Snegl ud af Sneglens Begreb? Men er det umuligt at
dialektisere det Existerende ud af dets Begreb, saa er det vel af samme
Grund ogsaa umuligt at dialektisere det ene Existensbegreb ud af det
andet, Hestebegrebet t.\ Ex.\ af Fluebegrebet, dette igjen af
Sneglebegrebet. I Erfaringens Verden falde \textit{diversi respectus}
udenfor hinanden, og hvor \textit{diversi respectus} falde udenfor
hinanden, spiller Dialektiken en meget underordnet Rolle. A. Men naar
det nu ikke vilde lykkes Dig at dialektisere Existensen ud af Begrebet
(Jvfr.\ Propæd.\ S.\ 31), har Du da aldrig forsøgt at gaae den omvendte
Vei, at dialektisere Begrebet ud
% --- p. 26 ---
af dets Existens, som nu t.\ Ex.\ at dialektisere dine Sneglebegreber ud
af deres tilsvarende Snegleexistens? C. Man dialektiserer ikke
Erfaringsbegreber ud af tilværende Gjenstande, man \emph{abstraherer}
dem. A. Der gives altsaa ikke Noget, man kunde kalde
Abstractionsdialektik? C. Ikke, saavidt jeg veed. A. Men de af
Gjenstandene abstraherede Begreber maae dog vel være, ikke blot i vor
Bevidsthed, men ogsaa i Gjenstandene selv? C. Forstaaer sig, thi
dersom Begreberne ikke vare i Gjenstandene, kunde de jo ikke abstraheres
af Gjenstandene. A. Men Gjenstanden, hvoraf Begrebet abstraheres, er
dog vel ikke selv et Begreb? C. Ingenlunde. A. Derimod kunne
Begreberne, forsaavidt de ere i Gjenstandene, vel siges at existere?
C. At Erfaringsbegreberne existere i de existerende Gjenstande, kan
vistnok Ingen betvivle. A. Her er i een og samme Existens altsaa to
Existerende, nemlig den existerende Gjenstand, og det i Gjenstanden
existerende Begreb? C. Nei, saadan kan det ikke tages; det er ikke
Gjenstanden selv, men dens Materiale, der er Begrebet modsat. A. Vi
faae vel altsaa tre Existerende: det existerende Materiale, det
existerende Begreb, og den Gjenstandens Existens, hvori begge existere?
C. Nei, saadan kan man heller ikke tage det. A. Hvordan skal man da
tage det? C. Materialet for sig er ikke et særligt Existerende,
saalidt som Begrebet, men Materialet og Begrebet ere to Sider, nemlig to
Existenssider af det samme existerende Ene, Gjenstanden. A. Det
existerende Ene er altsaa i Modsætning til og ligeoverfor de to
Existenssider det i og for sig Existerende? C. Nei, saadan kan man
heller ikke tage det; thi, dersom vi lade Abstractionen borttage de to
Existenssider, vil den existerende Eenhed jo selv blive et Abstractum,
en Existensside, der ikke existerer. A. Men hvordan skal man da tage
det? C. Saaledes, at den existerende Eenhed sees at existere i de to
Existenssider, som da just ogsaa derved komme til at existere i
uadskillelig Eenhed. A. Altsaa hvormange Dele,
% --- p. 27 ---
Beskaffenheder og Sider den existerende Gjenstand end iøvrigt kan have,
saa tør vi nu, ifølge din Forklaring, antage, at disse Dele med deres
Sider, Beskaffenheder og indbyrdes Forhold alle ere i lige Grad
materielle formedelst Materialet og intellectuelle formedelst Begrebet;
enhver Existensside bliver da paa denne Maade en existentiel Eenhed af
to Sider, den materielle og den intellectuelle. C. Dette tør vi ganske
sikkert antage. A. Men det i Gjenstanden Materielle er dog vel et
Sandseligt, der som sandseligt ikke kan tænkes, og det i Gjenstanden
Intellectuelle et Usandseligt, et Tænkeligt, der ikke kan sandses?
C. Fuldkommen rigtigt. A. I den existerende Eenhed er det Utænkelige
altsaa dog et Tænkeligt, det Usandselige et Sandseligt. C. Det
forstaaer jeg ikke! A. Saa forstaaer Du jo heller ikke den existerende
Eenhed; thi hvorledes skulde vel det Tænkelige i en existerende Eenhed
kunne gjennemtrænge det Sandselige, uden at det Sandselige derved,
formedelst det Tænkelige, selv skulde blive tænkeligt, og det Tænkelige,
formedelst det Sandselige, sandseligt? C. Det lader næsten, som om der
var Noget i, hvad Du siger. A. Existensbegrebet er altsaa et System
ikke blot af intellectuelle, men tillige af materielle Bestemmelser, et
System af intellectuelle Sandsebestemmelser, af sandselige
Intellectualbestemmelser. C. Ja, saaledes er det. A. Men saa
forraader jo Existensbegrebet en høist paafaldende Modsigelse.
C. Hvilken Modsigelse? A. At fordre Existens ved at modsætte sig
Existensen. C. Hvad skal det sige? A. Sættes det hele System af
sandselige Intellectualbestemmelser i den intellectuelle Form, saa
forsvinder det Sandselige som et i det Tænkelige ophævet
Reflexionsmoment, det vil sige: det Existerende forvandler sig til
Existensbegreb, men Existensen gaaer ud. Sættes omvendt det hele System
af intellectuelle Sandsebestemmelser i den sandselige Form, svarende til
umiddelbare Sandseindtryk, saa træder Gjenstanden i Existens, men
Existensbegrebet gaaer ud.... C. Nei, Existensbegrebet gaaer ikke ud,
det fuldkommes
% --- p. 28 ---
i Existensen: det Existerende er jo, ifølge vort Foregaaende, ikke
Andet, end det i alle Henseender fuldstændige Existensbegreb. A. Men,
var det ikke Dig, der, ifølge vort Foregaaende, indrømmede, at
Gjenstanden, hvoraf Begrebet abstraheres, ikke selv var Begrebet?
C. Ja, da tænkte jeg paa det abstracte og forsaavidt ufuldkomne Begreb,
ikke paa det i sig concrete og fuldkomne. A. Det er altsaa dog muligt
at dialektisere Existensen ud af Begrebet, naar man først har forstaaet
at dialektisere Begrebet ud af Existensen? C. At bevise dette i
Gjerningen, vilde vel her være det Samme som at vise det i bestemte
Exempler.

A. Den sunde Menneskeforstand siger os jo, at Hesten er et Dyr, at
Fluen er et Dyr, og Sneglen ligesaa; men hvorledes forholde slige Domme
sig nu egentlig til „Identitets- og Contradictionsprincipet?“ C. Sagen
er jo ligefrem. A. Ja, Sagen er, at dersom vi lade den formelle
Identitet gjælde for Prædikatet, saa fornegte vi den i hvert Subject;
lade vi den derimod gjælde for hvert Subject, saa fornegte vi den i
Prædikatet. C. Hvorledes? A. Sætte vi ubetinget: Dyr er Dyr, saa
bliver Hesten Flue, og Fluen Snegl, sætte vi ubetinget: Hesten er Hest,
Fluen er Flue, Sneglen Snegl, saa er Dyr (Hest) ikke Dyr (Flue),
d.\ v.\ s., saa modsiger Fælledsprædikatet sig selv, og maa bortfalde.
C. Denne Vanskelighed lader sig endda nok overvinde. A. Ved eller
uden Dialektik? C. Uden al Dialektik! man besinde sig blot paa
Existensbegrebernes faktiske Charakteermærker og den disjunktive Doms
Forhold til \textit{diversi respectus}. A. Det maae vi endelig høre!
C. Sammenligne vi Dyrene med Planterne, da sees det, at alle
existerende Dyreindivider have eet Fælleds: det at være Dyr;
sammenligne vi derimod Dyreexistenserne med hverandre indbyrdes, da
viser den ene Dyrerække sig ved givne Charakteermærker forskjellig fra
den anden. Sammenligne vi Existenserne i den ene Række med Existenserne
i den anden, da ere Individerne i hver Række hverandre indbyrdes lige;
men
% --- p. 29 ---
sammenligne vi derimod de under en given Række hørende Existenser med
hverandre indbyrdes, da ville særegne Mærker vise, hvorledes den ene
Klasse adskiller sig fra den anden, o.\ s.\ v.\ Dyret er altsaa, i
Modsætning til Planten, hvad Dyret er; men, ifølge almindelige og
constante Forskjelsmærker, er Dyret, i Modsætning til sig selv, enten
Bløddyr eller Leddyr eller Hvirveldyr. I Modsætning til Leddyr og
Hvirveldyr er Bløddyret altid Bløddyr, men i Modsætning til sig selv er
Bløddyret enten et egentligt Bløddyr eller et Straaledyr. Naar man nu
saaledes ganske ligefrem i hver Række, hver Klasse, hver Orden, hver
Familie, veiledet af de ved Erfaringer constaterede Mærker, vedbliver
med Inddelingerne (\textit{divisiones} og \textit{subdivisiones}),
indtil man gjennem Slægter og Arter har naaet Individerne; hvad
Modsigelse er saa heri, og hvad have slige Classifikationer med
Dialektiken at bestille? A. Det skulde da være, at Dialektiken fik
Adskilligt med \textit{diversi respectus} at bestille. C. Det har vist
ingen Fare. A. Siig os da i Kraft af din gjennemførte Classification:
hvad er egentlig en Hest? C. En Hest (\textit{Equus Caballus}) er,
efter sit Riges Kategorie, et Dyr, efter sin Rækkes Kategorie, et
Hvirveldyr, efter sin Klasses Kategorie, et Pattedyr, efter sin Ordens
Kategorie, en Tykhud, efter sin Familiekategorie, i Slægtskab med Æsler
og Muuldyr, en egen Artskategorie af en egen Slægtskategorie. A. Det
var da en skjøn Samling af Kategorier, som her alle maatte forene sig
for at afgive Hestens concrete Existenskategorie; naar man nu blot ved
Hjælp af \textit{diversi respectus} kunde faae det rigtig oplyst,
hvorledes disse mange og forskjellige Kategorier egentlig ere
sammenvoxne, for ikke at sige: sammenknippede, til een eneste! C. De
omgive naturligviis hinanden saaledes, at den almindeligere altid
omslutter den speciellere. A. Dog vel aldrig saaledes, som naar den
ene Convolut skal omslutte den anden, den ene Ring omgive den anden?
C. Nei, de hvile paa hverandre, idet den speciellere hver Gang har
afsat sig paa
% (p. 29 ends mid-sentence; the sentence continues on p. 30)

% II. Philosophien og de særskilte Videnskaber begins p. 30 (Almindeligt;
% A. Philosophiens Forhold til Mathematiken: Uendelighedsproblemet from p. 39).
% --- p. 30 ---
% (p. 30 opens with the last lines of part I's dialogue; the chapter head
% for II follows mid-page, after a section rule)
den almindeligere. A. Jo almindeligere en Kategorie er, desto
abstractere er den; jo abstractere den er, desto fjernere er den fra at
existere: det Existerende hviler altsaa paa det Ikke-Existerende?
C. Nei, jeg vilde sige: de hvile paa hverandre saaledes, at den
almindeligere altid bæres af en speciellere, denne af en endnu
speciellere, indtil vi komme til den allerspecielleste, der er Eet med
det Existerende selv. A. Det er altsaa ikke det at være Hest, der
beroer paa det at være Dyr, men det, at være Dyr, der beroer paa det,
at være Hest? C. Nei, saadan har jeg dog heller ikke tænkt mig det;
ja, nu kan jeg i dette Øieblik, reentud sagt, ikke sige, hvordan jeg
egentlig har tænkt mig det.

% (section rule as printed)

%% ============================================================
%%  II. PHILOSOPHIEN OG DE SÆRSKILTE VIDENSKABER  (pp. 30--476)
%% ============================================================
\chapter*{II.\\ Philosophien og de særskilte Videnskaber.}
\addcontentsline{toc}{chapter}{II. Philosophien og de særskilte Videnskaber}

% (short rule between chapter head and section head, as printed)
\section*{Almindeligt.}
\addcontentsline{toc}{section}{Almindeligt}

% (block quotation set letterspaced throughout, as printed; the citation
% is set unspaced)
\emph{„Med Hensyn paa det Aposterioriske vil Videns Fremskridt vistnok
være betinget af, at her bestandig indsamles nye Kundskaber, at nye
Kjendsgjerninger prøves, nye Opdagelser gjøres o.~s.~v.; men er nu
Sagen hermed afgjort? Dersom Videns Fremskridt i Forhold til det
Aposterioriske alene opfattes fra denne Side, er det let at afvise
Philosophien. Kan den ikke engang udtømme det Aprioriske, eftersom jo
Mathematiken ingen Tak er den skyldig, hvormeget mindre kan den da
indlade sig med positive Kjendsgjerninger! Vil Historikeren forelægge
den en Opgave: den kan ikke svare, thi en historisk Opgave besvares
% --- p. 31 ---
kun historisk. Vil en Chemiker henvende sig til den: den har intet
% (small baseline speck between "intet" and "Greb" — both OCR witnesses
% read a full stop; rejected on zoom as an ink speck, no stop printed)
Greb paa at experimentere, gjør ingen Opdagelser i Chemien. Vil
Astronomen hjælpes med en Iagttagelse: den forstaaer ham neppe. Vil
Lægen bede den om et Raad: den er ukyndig. Selv det, den paa egen
Haand har udspeculeret angaaende Forholdet imellem Sjæl og Legeme, kan
Physiologen maaskee desværre ikke engang bruge“.} (Propædeut.
S. 99--103.)

I hvilket Omfang det gjælder, at Philosophien er Dialektik, sees af
dens Forhold til Fagvidenskaberne, thi dette Forhold er i strengeste
Forstand dialektisk. Overseer man det Dialektiske, bliver Forholdet
tvetydigt og tvivlsomt; thi medens Fagvidenskaberne af Naturen ere
tilbøielige til at afvise Philosophien, er Philosophien saa igjen,
naar den opfatter sin Bestemmelse dogmatisk, fra sin Side fristet til
at ophøie sig selv paa de øvrige Videnskabers Bekostning; Schellings
Ord: „Speculation er Alt, er i Gud at skue, at betragte det, som er.
Videnskaben selv har kun forsaavidt Værd, som den er speculativ, det
er: Contemplation af Gud, som han er“ --- ere i denne Henseende
charakteristiske. (Propædeut. §~17, S. 81--83).

Hvad der vel i den nyere Tid især har bidraget til at forvirre
Begreberne, er navnlig den Omstændighed, at Hegel, hvis dybtgaaende
Opfattelse af Modsigelsen, af det Negative og af Udviklingens
immanente Methode (Propæd. S. 150--158) maa stille ham saa høit i
Dialektikernes Række, ikke desto mindre i sin Fremstilling af
Philosophiens Forhold til Fagvidenskaberne har overseet Forholdets
Dialektik. Istedetfor at anerkjende de særskilte Videnskabers
videnskabelige Selvstændighed, slutter Hegel sit propædeutiske
Overblik over Videnskabernes Encyklopædie (Propæd. S. 84--85) med den
Erklæring, at Encyklopædien egentlig er „en Fremstilling af
Philosophiens almindelige Indhold; thi hvad der i Videnskaberne
% --- p. 32 ---
er grundet paa Fornuft, afhænger af Philosophien, hvad der derimod er
i dem af det, der beroer paa vilkaarlige og udvortes Bestemmelser,
eller, som det hedder, er positivt og statutarisk, ligger udenfor
den.“ Men saaledes at lade Philosophiens Overlegenhed ligefrem beroe
paa Fagvidenskabernes Underkastelse, er at tage Sagen udialektisk; at
see Philosophiens Selvstændighed grunde sig paa dens Afhængighed af de
selvstændige Fagvidenskaber, er derimod at opfatte Forholdet
dialektisk.

Fagvidenskabernes Selvstændighed ligeoverfor Philosophien indlyser af
deres Udviklingsgang; thi de Methoder og Opdagelser, hvorved disse
Videnskabers store Fremskridt betegnes, ere saa langt fra at kunne
tilskrives Philosophien, at „de inductive Videnskabers Historie“
(Whewell) tvertimod vise, i hvilken Grad selvstændig Forskning og en
til Sagens Natur svarende Fremgangsmaade har været betinget af, at den
Forskende kunde løsrive sig fra eensidig Paavirkning af philosophiske
Speculationer (Propæd. S. 167--68). Hvad der i Fagvidenskaberne „er
grundet paa Fornuft“, afhænger ingenlunde særlig af Philosophien. For
det Første er det jo ikke Fornuften, der skylder Philosophien, men
omvendt Philosophien, der skylder Fornuften sin Oprindelse; dernæst,
og dette er i denne Sammenhæng det Afgjørende, maa det, ikke alene for
Fagvidenskabernes, men fornemmelig for Philosophiens egen Skyld
(Propæd. §~23), paa det Eftertrykkeligste fremhæves, at der gives en
Erkjendelse, som er i strengeste Forstand videnskabelig, uagtet den
ikke er philosophisk, ja en Tænkning, hvis videnskabelige Energie
netop beroer paa, at den ikke er philosophisk. Thi at der i den
empiriske Erkjendelse af det Empiriske er en Sandsningens Vished, en
Forestillingens Alsidighed, en Anskuelsens Tilgift, der ikke kan
erstattes af philosophiske Ideeblik, er dog vel ligesaa uimodsigeligt,
som det er uimodsigeligt, at der gives en exact og kritisk Tænkning,
der ikke lader sig ombytte med den dialektiske (Mathem. og Dialekt.
S. 3). De særskilte Videnskabers rationelle
% --- p. 33 ---
Eiendommelighed sees da ogsaa deraf, at hver enkelt især forudsætter
et eget tilsvarende videnskabeligt Talent.

Men opstiller man nu den Sætning: \emph{Fagvidenskaberne ere
uafhængige af Philosophien, men Philosophien er, navnlig hvad
Tilegnelsen af det positive Indhold angaaer, afhængig af
Fagvidenskaberne}, da kan denne i og for sig sande Forudsætning dog
ogsaa let foranledige en falsk Slutning. Er, siger man, Philosophien
i Alt, hvad det Exacte, det Empiriske og Positive angaaer, saa
afhængig af de exacte og positive Videnskaber, da maa denne
Afhængighed jo aabenbart være et Tegn paa dens Uselvstændighed.
Hvergang Philosophien nærmer sig til Berøring med et eller andet
bestemt positivt Emne, være sig af mathematisk, mechanisk, chemisk
eller historisk Indhold, maa den Philosopherende da naturligviis finde
sig i, ikke blot at vises tilrette med Kjendsgjerningerne, men endog i
Begrebsudviklingen selv at blive tilbørlig hovmestereret af
Fagmanden. Det seer altsaa ud, som vare Begreberne fuldkommen
tilforladelige, klare, tydelige og i alle Conseqvenser sikkre,
saalænge de bevares i den Form, hvori Faget bruger dem, men derimod
priisgivne til Forvanskning, saasnart de gjøres til Gjenstand for
philosophisk Eftertanke; det seer ud, som skulde den Philosopherende
laane Alt hos Fagmanden, medens denne for sit Vedkommende kan betragte
det Philosophiske som et aldeles Ligegyldigt. Denne Synsmaade beroer
paa en Misforstaaelse. Philosophiens Afhængighed af Fagvidenskaberne
er, væsentlig seet, et Tegn paa dens Uafhængighed, dens ideelle
Selvgyldighed.

I Modsætning til Omverdenens materielle Gjenstande ere jo vistnok alle
Videnskaber ideelle, have alle Idealitet; men i Modsætning til
Philosophien, som Idealitetens Repræsentant, maae Fagvidenskaberne
siges at forestille Realiteten. Bestemmes Modsætningen kategorisk,
vil det Reelle i dets Endelighed kjendes derpaa, at det hævder sig ved
at udelukke sit Andet, medens det Ideelle just beviser sin Realitet
derved, at
% (signature "3" at the foot of p. 33)
% --- p. 34 ---
\setcounter{footnote}{0}
det i sig optager Andet, og igjennem Andet forholder sig til sig
selv. At Mathematiken som Mathematik, Chemien som Chemie,
Naturhistorien som Naturhistorie, holde sig hver til
Sit\footnote{„Der er en Hovedregel, der aldrig kan krænkes ustraffet,
og som alligevel bestandig krænkes under vor Famlen fremad mod Lyset.
Det er: \emph{Aldrig at forsøge paa at løse en Videnskabs Problemer
ved den Ideerække, der tilhører en anden Videnskab.} Der er en Klasse
Begreber, der tilhører Physiken, en anden, der tilhører Chemien, en
tredie, der tilhører Physiologien, en fjerde, der tilhører
Psychologien, en femte, der tilhører Samfundsvidenskaben. Medens alle
disse Videnskaber ere nøie forbundne, har dog hver sit uafhængige
Omraade, der maa respecteres.“ G.~H. Lewes: Hverdagslivets
Physiologie, overs. 1ste Deel. Side 66.}, og, selv naar Bidrag hentes
fra andre Fag, dog allernærmest og ligefrem have hver med sit Eget at
gjøre, er just et Kjendetegn paa disse Videnskabers endelige Realitet;
at Philosophien derimod altid maa have sin ideelle Rigdom i reel
Fattigdom, sit Positive i sit Negative, er ved en sindrig Mythe oplyst
af Plato; ved Agathons Gjæstebud lader Plato, som bekjendt, sin
Sokrates fremstille Eros, den ideale Philosoph, som Søn af Poros og
Penia.

Den realistiske Videnskab, engang grundlagt, kan vel i Tidernes Løb
behøve mangfoldig Omdannelse og undergaae mange Forandringer, men dens
Støtter og Grundvolde ville med alt dette forblive urokkede; skal den
idealistiske Videnskab derimod omdanne sin Bygning, vil et
\textit{de omnibus dubitandum} efterhaanden rokke ved alle Støtter;
hele Bygningen maa da nedrives indtil Grunden, at den i sin nye Stiil
kan opføres fra Grunden af. Hvergang et sligt Knudepunkt indtræder,
seer det i Øieblikket jo vistnok ud, som om Realiteten stod paa Spil,
som var det nu aldeles ude med Philosophien; men Historien stadfæster,
hvad da ogsaa ligger i Sagens Natur, at slige gjennemgaaende Kriser
just betegne Vendepunkterne i Philosophiens Gjenfødelse.

% --- p. 35 ---
% (block quotation letterspaced as printed; citation unspaced)
\emph{„Fagvidenskaberne ere solide, have altid fast Arbeide, rigeligt
Udkomme og sikker Anerkjendelse i den lærde Republik; Philosophien er
i Sammenligning hermed som en Omreisende uden endelig Forsørgelse, en
svævende Skikkelse, der synker og stiger, døer hen og lever op igjen,
eftersom Aanden vil det“.} (Propæd. S. 102). --- Men det er dog kun
fra et aldeles udvortes realistisk Synspunkt, at det Ideelle blot
opfattes som det Vage, det Ubestemte, i sig Ubestemmelige; den
dialektiske Form er en i sig bestemt og bestemmelig Form, og
Philosophiens Forhold til Fagvidenskaberne bliver desaarsag at angive
kategorisk.

For at Philosophien skal kunne sætte sig i et indre og væsentligt
Forhold til en særskilt Videnskab, maa den i samme opdage et Problem,
den tillige kan kalde sit Problem. Til at oplyse, hvorledes et
saadant \emph{Fælledsproblem} vælges og behandles, ville vi, for ikke
at sprede Opmærksomheden paa formange Punkter, i det Følgende
indskrænke os til Exempler, hentede fra Mathematiken, Chemien og
Physiologien; vi have da saaledes i Mathematiken:
\emph{Uendelighedsproblemet}, i Chemien: \emph{Forandringsproblemet},
i Physiologien: \emph{Livsproblemet}.

Men selv om slige Fælledsproblemer, naar de betragtes fra
Idealitetssiden, ere nok saa charakteristiske, saa er det dog hermed
ingenlunde afgjort, at de ogsaa vise sig saaledes, naar de opfattes
realistisk; fordi den Philosopherende i en eller anden bestemt
Fagvidenskab har mærket sig et eller andet gjennemgaaende Problem,
derfor er det endnu ikke sagt, at Fagmanden paa tilsvarende Maade vil
vedkjende sig det samme Problem. Den Chemiker, der absolut vilde
modsætte sig en philosophisk Opfattelse, kunde jo t.~Ex. raisonnere
saaledes: „Der gives i Chemien egentlig slet ikke Noget, man kan kalde
et Forandringsproblem, men vel et Forandringsbegreb, der opstaaer og
kommer til Anvendelse, forsaavidt
% (signature "3*" at the foot of p. 35)
% --- p. 36 ---
Chemien undersøger de Naturlove, hvorefter Legemernes indre
Forandringer foregaae. Men Et er det at undersøge de Love, hvorefter
Forandringerne foregaae, et Andet, at ville udgrunde Forandringens
Væsen, speculere over dens Mulighed og Umulighed, en Opgave, der er
Chemien aldeles fremmed.“ I samme Aand kunde en Physiolog, der vilde
afvise Philosophien, raisonnere saaledes: „Vistnok gaaer Physiologien
ud fra den Forudsætning, at der er et Noget, vi maae kalde Liv, siden
der i de levende Organismer er et lovbestemt System af
Livsphænomener; men da Physiologien væsentlig indskrænker sig til at
undersøge Phænomenerne, uden at trænge ind i det bagved liggende
Skjulte, kan man paa ingen Maade indrømme, at her skulde være noget
egentligt Livsproblem; hvad Livet i sig selv er, et Princip, en Kraft,
et Produkt af mechaniske og chemiske Kræfters eiendommelige
Samvirken, ligemeget! man danne sig hvilken individuel Forestilling om
Livets Væsen, man vil; naar man kun ikke derved afholdes fra at indsee
Lovenes Gyldighed og eftervise de physiske Aarsager, er Physiologien
tilfredsstillet.“

De Vanskeligheder, forskjellige Standpunkters naturlige Indflydelse
paa Opfattelsen af et videnskabeligt Problem altid medfører, forhøies
betydeligt derved, at der i Undersøgelserne saa let blander sig noget
Menneskeligt, d.~v.~s. noget Eensidigt, Tilfældigt, Vilkaarligt,
Lidenskabeligt; da nu Philosophien med al sin Idealitet er, som enhver
anden Videnskab, undergiven menneskelige Vilkaar, bliver det i denne
Sammenhæng af Vigtighed ogsaa at tage Sagen fra den reent menneskelige
Side, og saaledes mærke sig de Forpligtelser og Rettigheder,
Videnskaberne under gjensidige Conflicter have at opretholde.

At de omhandlede Fælledsproblemer maae i Philosophien tage sig
anderledes ud, end i Fagvidenskaben, er en Selvfølge.
Forandringsproblemet er, om ogsaa Chemikeren paa sin Maade vil
vedkjende sig det, jo alligevel i Chemien
% --- p. 37 ---
ingenlunde, hvad det er i Philosophien; ligeledes er Livsproblemet,
selv naar Physiologen indrømmer dets Realitet, et Andet i
Physiologien, end det er i Psychologien; den Philosopherende, der
kjender Problemet i Andethedens Form, og er sig bevidst, at det Samme
dog ikke er det Samme, maa da fornemmelig lægge an paa at klare
Udviklingens Betingelser og tydeliggjøre Forvandlingen. Enhver, der
vil søge Veiledning i en særskilt Videnskab, maa naturligviis sætte
sig ind i dens Tankegang. Forsømmer den Philosopherende dette,
misforstaaer han en paagjældende Videnskabs Charakteer enten i det
Hele eller i væsentlige Enkeltheder, være sig nu, at han tillægger den
et Formaal, den ikke har, og ikke kan have, eller at han udriver dens
Læresætninger af deres rette Sammenhæng, og saaledes forvansker dem,
eller at han, ved at oversee de særegne Forudsætninger og Betingelser,
under hvilke et givet Begreb med den tilsvarende Betegnelse har dannet
sig, underlægger Begrebet en falsk Betydning, og paa denne Maade
kommer til at paatvinge Videnskaben Anskuelser, Resultater og
Conseqvenser, den ikke kan vedkjende sig, da er Fagmanden i sin Ret,
naar han tager til Gjenmæle. Denne Ret bliver endog en formelig
Sandhedspligt, naar slige Misgreb fra Philosophiens Side, hvad ikke
saa sjældent har været Tilfældet (Propæd. S. 170--171), ovenikjøbet
ere forenede med den Indbildning, at det er Philosophien, der fra sit
ophøiede Standpunkt maa føle sig kaldet til at reformere
Fagvidenskaberne. Men videre end til en Afviisning af philosophiske
Overgreb kan og bør Fagmanden ikke gaae; gaaer han endnu videre, saa
at han t.~Ex. ikke nøies med at sikkre Begrebet den Betydning, det har
i hans Videnskab, men endog i den Grad betragter det som Faget
tilhørende, at han enten reentud forbyder den Philosopherende at røre
ved det, eller dog paatager sig at foreskrive, hvorvidt, hvorledes og
i hvilke Forbindelser Philosophien tør udvikle det laante, af
% --- p. 38 ---
Fagvidenskaben godhedsfuldt overladte Begreb: da sees heri ikke blot
en Anmasselse, men en latterlig Misforstaaelse.

% (block quotation letterspaced as printed; citation unspaced)
\emph{„Indskrænkede til de forskjellige Bevidsthedsrubrikker ere de
vundne Begreber endnu saa sammenvoxne med endelige Bestanddele, endnu
saa betyngede af partikulære Forudsætninger, at de befinde sig som i
en Tilstand af Ufrihed; de holdes her som Stoffets Fangne indenfor en
bestemt Synskreds, de ere saa at sige Fagets Livegne, der ikke tør
forlade den bestemte Arbeidsplads.“} (Propædeutik S. 113.)

Idet nu Philosophien lægger an paa at frigjøre Realbegreberne fra det
i deres Begrændsning Tilfældige, udfolde Conseqvenserne og bringe de
schematiske Former i dialektisk Strømning, maa den
% (raised speck after "den" — both OCR witnesses read an apostrophe;
% rejected on zoom as an ink speck, nothing printed)
naturligviis paa sit Omraade, naar den først har draget Begreberne ind
under sin Synskreds, stillet Problemerne paa sin Maade, og saaledes
taget hele Behandlingens Ansvar paa sig, have en ubestridelig Ret til
at anlægge sin Maalestok, indføre sine Betegnelser, sin Sprogbrug, og
foretage Udviklingen i sin Aand. Staaer det fast, at enhver
Fagvidenskab i det Væsentlige selv maa indeholde sine egne
Correctiver, da maa vel ogsaa Philosophien indeholde sine
Correctiver. En Bedømmelse af de Feil og Misgreb, en Philosopherende
ved Gjennemførelsen af de dialektiske Analyser muligviis kan have
begaaet, høre i sidste Instans ikke under Fagvidenskabens, men under
Philosophiens Forum.

At negte Philosophien fuld videnskabelig Ret til at behandle, hvad den
fra Fagvidenskaberne har tilegnet sig, paa fri philosophisk Maade, er
at negte den en eiendommelig Begrebsudvikling, og anvende Magtsprog.
Skal det gaae ved Magtsprog, ja da kan en Fagmand vel sagtens afgjøre
Philosophiens Skjæbne ligesaa let, som hiin Kalif afgjorde det
alexandrinske Bibliotheks Skjæbne: „Enten indeholder det kun, hvad der
staaer i Koranen, og da er det overflødigt,
% --- p. 39 ---
eller det indeholder Andet, og da er det falsk;“ --- „enten indeholder
Philosophien kun, hvad der staaer i Fagvidenskaben, og da er det
overflødigt, eller den indeholder Andet, og da er det falsk.“ Det er
et saracenisk Dilemma; men et saracenisk Dilemma kan ikke bringe
Philosophien til at opgive sin Ret.

% (section rule as printed)
\section*{A.\\ Philosophiens Forhold til Mathematiken:\\
Uendelighedsproblemet.}
\addcontentsline{toc}{section}{A. Philosophiens Forhold til
Mathematiken: Uendelighedsproblemet}

% (rule under the section head, as printed; the block quotation is
% letterspaced as printed, and here — unlike pp. 31, 35 and 38 — the
% citation is letterspaced too, so it stays inside the \emph)
\emph{„En philosophisk Ræv eller Maar har listet sig over
Mathematikens fede Fjerkræ, for at fylde sin schematiske Vom; han vil
udsuge deres varme Blod, saa at der kun bliver de døde, saft- og
kraftløse Kroppe tilbage, der da skulle fremstilles som Mathematikens
driftige, praktiske, rørige Begreber“ (Hr.~Professor Steens
„Afregning“ S. 7).}

Disse varslende, ja ildevarslende Ord kunne, især naar de sættes i
Forbindelse med en Erklæring som denne: \emph{„Der er for det Første
ikke Noget, man kan benævne „Uendelighedsproblemet“ i Mathematiken“}
(Jvfr. Begrebet: Nøiagtighed. S. 38), jo vistnok ved første Øiekast
synes at have noget særdeles Afskrækkende; men, naar det saa derhos
indrømmes, at der gives \emph{„et mathematisk Begreb „Uendelig“, der
opstaaer og kommer til Anvendelse, og som er yderst vigtigt og tilmed
vanskeligt, fordi det er af blandet Natur“}, forsvinder det
Afskrækkende, og her viser sig et Indbydende. Er det mathematiske
Begreb „Uendelig“, der opstaaer og kommer til Anvendelse, nemlig,
ifølge Mathematikerens egen Tilstaaelse, \emph{„yderst vigtigt og
tilmed vanskeligt, fordi det er af blandet Natur“}: saa indeholdes jo
heri en for Philosophiens Vedkommende
% --- p. 40 ---
vel begrundet Opfordring til at undersøge, om dette i Mathematiken saa
\emph{yderst vigtige og tilmed vanskelige} Uendelighedsbegreb
% (the line-end "Uende-" of "Uendelighedsbegreb" may also be lightly
% letterspaced; its second half is set tight, so the emphasis is taken
% to end at "vanskelige")
dog ikke paa en eller anden Maade, Mathematikeren maaskee uafvidende,
skulde hænge sammen med et i det Hele gjennemgaaende mathematisk
Problem, om Philosophien i dette gjennemgaaende Problem ikke maaskee
skulde have al Grund til at vedkjende sig sit eget ældgamle
Uendelighedsproblem, og om Mathematiken, uden at en Mathematiker altid
saa lige tænker derpaa, dog ikke skulde kunne siges at afgive
væsentlige og uundværlige Bidrag til dette Problems Løsning.

\subsection*{a) Uendelighedsbegrebet.}
\addcontentsline{toc}{subsection}{a) Uendelighedsbegrebet}

At Begrebet „Uendeligt“, der i Mathematiken „opstaaer og kommer til
Anvendelse“, er „yderst vigtigt“, kunne vi strax indrømme; derimod
maae vi paa det Bestemteste benegte, at Uendelighedsbegrebet skulde,
som Begreb for sig, være enten vanskeligt eller af blandet Natur. Om
det Uendelige kan i Almindelighed kun siges dette, at det er en
Negation af det Endelige, det Ikke-Endelige. Skulde der i en saadan
Negation skjule sig nogen Vanskelighed, da kunde Vanskeligheden dog
ikke umiddelbart hidrøre fra Negationen selv, men maatte væsentlig
være at henføre til det, som ved Negationen blev ophævet. Vil man
altsaa finde Vanskeligheder ved det mathematiske „Uendeligt“, maa man
finde tilsvarende Vanskeligheder ved det mathematiske „Endeligt“; gjør
man ingen Ophævelser over det Sidste, da forraader det kun
Tankeløshed at gjøre Ophævelser over det Første. Skulde
Uendelighedsbegrebet være vanskeligt, „fordi det er af en blandet
Natur“, mon da ikke Endelighedsbegrebet, naar det Endelige særlig
fremhæves, være af en endnu mere blandet Natur? Vil man indvende, at
det mathematiske „Endeligt“, der kun gjælder om Størrelserne, ikke om
Tingene, umuligt kan være af blandet Natur, da bliver man vel ogsaa
nødt
% --- p. 41 ---
til at indrømme, at de Negationer, hvori et saadant „Endeligt“
forsvinder, ligesaalidt kunne være det.

Begreber af blandet Natur have altid en betænkelig Tilsætning af det
Begrebløse, og staae som Blandingsbegreber paa det laveste Trin i
Begrebernes Række. Hvorledes fremkommer et Blandingsbegreb? Derved
at uligeartede, paa udvortes, tilfældig Maade sammenbragte, Dele
henføres under en collectiv Eenhed, der da som Følge heraf kommer til
at indeholde uligeartede, tilfældig sammenbragte Kjendemærker.
Saaledes ere Begreberne: „Muldjord“, „Gruus“, „Mergel“, ægte
Blandingsbegreber; men et Gruusbegreb og et mathematisk
Uendelighedsbegreb --- hvilken Forskjel!

Men, vil maaskee Nogen spørge, naar Begrebet: „Uendeligt“, i hvormange
og forskjellige Forbindelser det end forekommer, hvormange
Forvandlinger det end undergaaer, dog altid bevarer een og samme
uforanderlige Grundbetydning, den nemlig, at være:
„Endelighedsnegation“, hvordan gaaer det da til, at et saa reent, saa
negativt logisk Begreb kan i en Mathematikers Øine faae Udseende af at
være „af en blandet Natur“, eller, hvad der er det Samme, at høre til
Blandingsbegreberne? Svaret er ikke vanskeligt: \emph{man blander
Tilfældigt med Væsentligt, man udfinder ved denne Blanding endeel
forskjelligartede Betydninger; man udleder ikke disse særskilte
Betydninger af Grundbetydningen, efterviser ikke den enes nødvendige
Overgang i den anden, men stiller dem udvortes op ved Siden af
hverandre, gjør dem til Gjenstand for Opregnen, og see, saa gaaer det
som af sig selv.} Stiller man saadanne Betydninger, som: „Negation i
positiv Form“, „Umuligt“, „Overgang med noget Qvalitativt“ o.~s.~v.
udvortes sammen, faaer man strax et Indtryk af det Blandede.
Hvorledes skal man vel kunne adskille en Negation i positiv Form fra
en Negation i negativ Form, uden derved, at hiin udtrykker det
Negative paa middelbar,
% --- p. 42 ---
denne paa umiddelbar Maade; men stilles nu det Endelige imod det
Uendelige, da er det jo ikke det Uendelige, der udtrykker Negationen i
positiv Form, men tvertimod det Endelige, forsaavidt dette, ihvorvel
det har en positiv Form, dog middelbart bestemmes som det
Ikke-Uendelige. At opfatte Negationens eget umiddelbart negative
Udtryk som dens positive Form røber en comisk Forvexling. Vilde
Nogen, for at vise sin Skarpsindighed i det Logiske, t.~Ex. paastaae,
at Sætningen: „$a$ er ikke Noget“, vel var ligefrem benegtende, men
Sætningen: „$a$ er Intet“, derimod bekræftende, idet Negationen er
lagt i Ordet: „Intet“, vilde man vistnok let blive enig om at henregne
slige Skarpsindighedsprøver, formulerede i al Keitethed efter en
Formellogiks \textit{S} er \textit{Non} --- \textit{P}, om ikke
ligefrem til det Pedantiske, saa dog til den Art Umodenhedsforsøg, der
forraader, at den Skarpsindige endnu ikke har lært at distingvere
imellem frugtbare og ufrugtbare Distinctioner. Indrømmer man nu, for
ikke at faae Formeget af det logiske Blandingsgods, at Sætningen „$a$
er Intet“, er ligesaa direct negativ, som „$a$ er ikke Noget“, saa vil
man vel ogsaa indrømme, at Sætningen: „$a$ er Nul“, er ligesaa direct
negativ, som: „$a$ er Intet“. Men, staaer det fast, at Tegnet $0$ er
et ligesaa direct og umiddelbart Negationsudtryk, som Ordet: Intet,
saa hører der dog en ganske abnorm Skarpsindighed til, for at indsee,
at medens den Sætning: „imellem $a$ og $a$ er der ikke nogen
Forskjel“, er ligefrem benegtende, skulde Ligningen $a - a = 0$
derimod være saa nagelfast i det Positive, at kun en Omskrivningens
Skrue var istand til at trænge igjennem til Negationen. (Jvfr. Om
Begrebet: Nøiagtighed. S. 39).

Det er en egen Sag med den logiske Udtydning af mathematiske
Ligninger; jo større Vægt her lægges paa Begrebssiden, desto mere
indlyser det, at Philosophien ogsaa her har et Ord at sige. Ifølge
naturlig mathematisk Tankegang viser Ligningen
$\frac{a}{0} = \infty$ os det velbekjendte Forhold imellem
% --- p. 43 ---
Divisor og Qvotient i dettes yderste Conseqvens: vedbliver Divisor
under Forudsætning af en constant Dividend, at indsvinde i det
Uendelige, maa Qvotienten vedblive at voxe i det Uendelige; her er i
dette Forhold Intet, der minder om det Umulige, Alt tyder tvertimod
paa det Mulige: til den mindst mulige Divisor svarer den størst mulige
Qvotient. Vil man, for at det Uendelige kan komme til at betegne det
Umulige, derimod vælge et andet Synspunkt, og fortolke Ligningen
saaledes: „Naar der spørges, hvormange Gange skal $0$ gjøres større
for at give $a$, saa svares i Ord: det kan man ikke frembringe ved at
multiplicere $0$, men i Tegn: $\frac{a}{0} = \infty$, $0$ skal gjøres
uendelig mange Gange større, for at give $a$“. (Om Begr. Nøiagt.
S. 40.) ---, da er en saadan Fortolkning vel skikket til at befordre
allehaande Indblandinger af et ulogisk Allehaande. Dersom Ligningen
$\frac{a}{0} = \infty$ virkelig og for Alvor var lagt an paa at
besvare det i sig urimelige Spørgsmaal: „hvormange Gange skal $0$, det
absolute $0$, det rene Intet, gjøres større for at give $a$“, da vilde
Betegnelser som: $\frac{a}{0} = 0$, $\frac{a}{0} = a$, have været
langt mere adæqvate. Thi af slige Udtryk sees øieblikkelig det
Umulige, det i Fordringen Absurde, medens Ligningen:
$\frac{a}{0} = \infty$ ved sin ironiske Anviisning paa det Uendelige
kun mimisk antyder det i Spørgsmaalet Urimelige, og saa derhos ganske
rigtig ved „en Negation i positiv Form“ fixerer den Spørgende, fordi
han spørger, ikke som en mathematisk Tænker, men som en logisk Tosse.
Opgiver man det logiske Tosseri at ville udregne, hvormange Gange det
rene Nul skal gjøres større for at frembringe $a$; lader man, for at
Ligningen kan give nogen Mening, $0$ betegne den mindst mulige
Divisor, idet $\infty$ betegner den tilsvarende størst mulige
Qvotient: da sees fra alle Sider,
% --- p. 44 ---
hvorledes det Umulige viger Pladsen for det Mulige. Af Ligningen:
$\frac{a}{0} = \infty$ følger jo ligefrem, i Kraft af det omspurgte
Forholds Begreb, $a = \infty \frac{a}{\infty} = \infty.0$; ligeledes
følger af $\frac{b}{0} = \infty$, $\frac{c}{0} = \infty$,
$\frac{d}{0} = \infty\ldots$ $b = \infty.\frac{b}{\infty} = \infty.0$,
$c = \infty \frac{c}{\infty} = \infty.0$,
$d = \infty.\frac{d}{\infty} = \infty.0\ldots$
% (the multiplication dot after the first ∞ is printed inconsistently:
% present in the b- and d-equations and at "altsaa: a = ∞.a/∞" below,
% absent in the a- and c-equations; transcribed as printed)
Istedetfor Umuligheden af, at $\infty.0$ skulde være istand til at
give $a$, see vi tvertimod Muligheden af, at $\infty.0$ kan give, ikke
blot $a$, men $b$, $c$, $d\ldots$, kort saamange forskjellige
Størrelser, det skal være; istedetfor det korte, afvisende „Umuligt“,
faae vi jo netop Mulighed paa Mulighed, Muligheder i Overflødighed.
Dette Resultat er naturligviis saa langt fra at overraske den
\emph{tænkende} Mathematiker, at det end ikke overrasker den
almindelige sunde Menneskeforstand. Enhver forstaaer jo, at en
hvilkensomhelst endelig Størrelse kan --- det er en Mulighed --- siges
at bestaae af uendelig mange uendelig smaa Dele, altsaa:
$a = \infty.\frac{a}{\infty} = \infty.0$, og derfor ogsaa kan --- det
er en Mulighed --- deles i uendelig mange uendelig smaa Dele, altsaa
$\frac{a}{\infty} = 0$.

Den Fixidee, at Begrebet: „Uendeligt“ skulde, som „Negation i positiv
Form“, være mere skikket til at betegne „det Umulige“, end Begrebet:
„Endeligt“, er let at belyse. „Det Umulige“ er nemlig i og for sig
ligesaa indifferent mod „Uendeligt“ som mod „Endeligt“; thi det
Umulige viser sig, hver Gang en opstillet Fordring tilintetgjør sig
selv; men slige, sig selv tilintetgjørende, Fordringer kunne ligesaa
godt have det Endelige som det Uendelige til Forudsætning. Det er,
for at blive ved det Uendelige, umuligt at \emph{opregne} alle Led i
en uendelig Række; det er, for at gaae til det Endelige,
% (p. 44 ends mid-sentence; the sentence continues on p. 45)

% --- p. 45 ---
et Menneske umuligt at \emph{tælle til} en Qvadrillion; det er, at
vi paa een Gang skulle støtte os baade til det Endelige og
til det Uendelige, umuligt, at en Kugle med en uendelig
Radius kan være en Kugle, eller et Plan en Kugleoverflade
med endelig Radius, og ligesaa umuligt er det, at Ligningen:
$a = 2a = \infty.a = \frac{0}{a}$ kan have Realitet. Vil man i det
mathematiske Tegnsprog vælge en mere ligefrem Betegnelse for
noget i sig Umuligt, da er det ikke $0$ og $\infty$, men $\sqrt{-1}$,
$\sqrt[2m]{-a}$ o.~dsl., man har at holde sig til. (Den af R.~N. gjennemsete
„Afregning“. Bilag. S. 26--27).

I logisk Almindelighed sees det let, at Begrebet „Uendeligt“
ikke som Begreb optager ueensartede Bestanddele, og
altsaa ikke kan regnes til Blandingsbegreberne; at Begrebet
ikke desto mindre forekommer i flere Betydninger, viser sig i
et forskjelligt Lys, og forsaavidt kan, naar det behandles
overfladisk, faae et Anstrøg af at være „af en blandet Natur“,
har sin Grund i noget ganske Andet: at gjennemtænke Begrebet
er at udfolde Problemet.

\textbf{b) Uendelighedsproblemet.}
\addcontentsline{toc}{subsection}{b) Uendelighedsproblemet}

Vilde Nogen begynde en Indsigelse imod en philosophisk
Henviisning til den Maade, hvorpaa Endelighedsproblemet
finder sin Løsning i Mathematiken, med følgende Erklæring:
„Der er for det Første ikke Noget, man kan benævne Endelighedsproblemet“
i Mathematiken, men vel et mathematisk
Begreb „Endelig“, der opstaaer og kommer til Anvendelse,
og som er yderst vigtigt og tilmed vanskeligt, fordi det er af
% printed thus: the objection opens „Der ... and the compositor closes the
% quote both after "Endelighedsproblemet" and after "Natur" below without
% re-opening one before "blandet Natur" — one “ in excess on this page,
% transcribed as printed (balance -1 here, offset by the +1 on p. 46).
blandet Natur“: da vilde man dog vistnok finde, at det var
en underlig Tale. Hvad en endelig Størrelse er, er snart
sagt, men Problemerne, der angaae de endelige Størrelsers
indbyrdes Forhold, ere utallige; istedetfor at sige: der gives
i Mathematiken intet Endelighedsproblem, maatte man da
% --- p. 46 ---
snarere sige: der gives i Mathematiken utallige Endelighedsproblemer.
Dersom nu disse Problemer faldt ud fra hverandre
i løs Mangfoldighed, saa kunde man vel benegte, at
der gaves et igjennem dem alle gaaende Almeenproblem, men
saa maatte man med det Samme benegte, at der gaves en
mathematisk Methode, at, med andre Ord, Mathematiken var
en Videnskab. Er Mathematiken en Videnskab, gives der en
Methode, efter hvilken de forskjellige Endelighedsproblemer
dannes og løses i systematisk Sammenhæng, saa gives der
ogsaa et gjennem dem alle gaaende Almeenproblem, et Methoden
i alle Forgreninger omfattende og bestemmende Grundproblem,
et Problem, der til sin alsidige Løsning udkræver intet Mindre
end hele „det Endeliges Analyse.“

Men hvad der saa uimodsigelig gjælder om „Endelighedsproblemet“
i dets Forhold til „det Endeliges Analyse“, skulde
dette ikke ogsaa gjælde om Forholdet imellem „Uendelighedsproblemet“
og „det Uendeliges Analyse“? hvad der saa uimodsigelig
gjælder om Begrebet „Endelig“ i dets Forhold til
„Endelighedsproblemet“, at det nemlig ikke er Begrebet, men
Problemet, hvori Vanskeligheden ligger, skulde dette ikke ogsaa,
naar vi see nærmere til, befindes at gjælde om Begrebet
% printer's defect: the quote opened before "Uendelighedsproblemet?" is
% never closed — verified on the image; transcribed as printed (balance +1).
„Uendelig“ i dets Forhold til „Uendelighedsproblemet? ---

\emph{α) Begreb og Problem.} Som Særkjende for de
mathematiske Begreber, hvormegen Lighed de end for Resten
have med de almeenlogiske, maa det udtrykkelig fremhæves
(Mathem. og Dial. S. 21--27), at de ere Operationsbegreber;
heraf følger, at, medens de i deres Almindelighed ingen
Vanskelighed frembyde, vil den analyserende Tænkning snart
opdage, at deres eiendommelige Udvikling altid maa foregaae
i nøieste Forbindelse med de Operationer, hvis logiske Væsen
ved dem søges opklaret. Hvad en Logarithme er, kan let
angives i en saa exact Definition, at den sunde Forstand
formaaer at fastholde Begrebet; men de mangfoldige og
charakteristiske Conseqvenser, hvori Begrebets Indhold først egentlig
% --- p. 47 ---
lader sig tilsyne, komme kun frem under Operationernes
Gang. Hvad man forstaaer ved at dele en Linie i yderste
og mellemste Forhold, eller, som det ogsaa hedder, at høidele
Linien, er snart sagt, og let fattet; Begrebet „høidele“ frembyder
ingen Vanskelighed; Vanskeligheden, dersom her er nogen,
er: at løse Høidelingens Problem, og udvikle alle deri indeholdte
Conseqvenser. Hvad der gjælder om de særskilte Operationsbegreber,
gjælder ogsaa om de almene; kun af Endelighedsproblemernes
systematiske Løsning fremlyser det, hvad der
mathematisk ligger i Begrebet: „Endeligt“; kun en Gjennemførelse
af „det Uendeliges Analyse“ kan aabne os et videnskabeligt
Indblik i det mathematiske Uendelighedsbegreb. Af
Operationsconseqvenserne sees endvidere, at det Uendelige ikke blot
falder udenfor, men ogsaa indenfor det Endelige (Jvnfr. Phil.
og Mathem. S. 20--24; Mathem. og Dialekt. S. 28--38),
at det paa mangfoldige Maader udvikler sig af Endeligheden
selv: havde der under „det Endeliges Analyse“ ikke oprindelig
rørt sig og i mange særskilte Glimt yttret sig et gjennemgaaende
Uendelighedsproblem, var Differential- og Integral-Regningen
aldrig bleven uddannet. Da Leibnitz løste „Tangentproblemet“,
var det ikke Tangenten og Secanten, (Jvnfr. Phil. og Mathem.
S. 64), men det articulerende Vexelforhold af Endeligt og
Uendeligt, hvorom det væsentlig gjaldt: i „Tangentproblemet“
sees „Uendelighedsproblemet“.

\emph{β) Mathematisk Opfattelse.} At tænkende Mathematikere,
hvormegen Vægt de end iøvrigt lægge paa Operationsmethodens
tilfredsstillende Resultater, dog desuagtet aldrig
have kunnet faae det i deres Tanker ubegribelige Begreb paa
en tilfredsstillende Maade adskilt fra den paa dette Begreb
hvilende Methode, viser Mathematikens Historie lige fra
Leibnitz og Newton (Jvfr. Mathem. og Dial. S. 90--91,
S. 95 o. fl.) indtil den nærværende Tid. „For det Første
kunne vi, naar Principet er ubekjendt, jo ikke vide, om ogsaa
alle deraf følgende Conseqvenser lade sig indrømme: vi gaae
% --- p. 48 ---
\setcounter{footnote}{0}
som i Blinde, og det turde hænde sig, at vi efter møisommelig
Anstrengelse tilsidst stødte paa Noget, som netop udviste
Principets Natur, d.~v.~s. paa Noget, vi ikke kunde begribe. For
det Andet: er det ikke en skammelig Nedværdigelse af Fornuften,
at lade den følge en Sammenkjædning af Led, en
Række, hvis Begyndelse og Ende den ikke kan gjennemskue?
Det er et stort Onde, der fødes af den Slags Routine:
vant til at operere uden at indsee Grunden taber Aanden al
Øvelse i at tænke, og ender med at udføre Alt ganske
maskinmæssigt“. (Mathem. og Dial. S. III). Disse Savriens
% "Savriens" as printed (presumably for Savérien); not corrected.
Betænkeligheder forklare tilstrækkelig den Uro og Spænding,
hvori det Uendeliges Analyse bestandig har holdt den mathematiske
Reflexion\footnote{\textit{See Carnot: Reflexions sur la
Métaphysique du Calcul Infinitésimal. Paris 1797} (übersetzt von
Hauf. Frankf. a.~M. 1800). Versuch einer Philosophie der Mathematik
--- mit einer Kritik der Aufstellungen Hegel's über den Zweck und die
Natur der höheren Analysis --- von H. Schwartz. Halle 1853.}
(Mathem. og Dial. S. 55--56), hvis Problem
--- saaledes som en af Berlinerakademiet i Aaret 1784 udsat
Priisopgave charakteristisk nok udtrykker det, „wie aus einer
widersprechenden Voraussetzung so viele wahren Sätze haben
können hergeleitet werden“ --- ikke vil ophøre at vende tilbage,
om end „Routinen“, den maskinmæssige Færdighed, til enhver
Tid sikkert vil yde „\textit{habiles Calculateurs}“ saa rigelig
Tilfredsstillelse, at det let kan undgaae en og anden Regnemesters
Opmærksomhed, at der i Mathematiken gives Noget, man kan
benævne „Uendelighedsproblemet“.

\emph{γ) Philosophisk Opfattelse.} „For at gjøre Uendelighedsproblemet
kjendeligt, maa Reflexionen støtte sig til
Operationen; for at løse Problemet, maa den tilegne sig den
i Operationen udtalte Intelligens. At der imellem uendelige
Størrelser kan være et endeligt Forhold, at en Sum af
uendelig mange Led kan i nogle Tilfælde udgjøre en endelig,
% --- p. 49 ---
i andre en uendelig Størrelse, at en Række, hvis Led forholde
sig saaledes, at Grændsen for Forholdet af to paa hinanden
følgende Led giver Qvotienten 1, kan være i nogle
Tilfælde convergerende, i andre divergerende o.~s.~v. Sligt
vil overalt undgaae den Tænkning, der ikke ved Operationens
Vink faaer Øie paa de exacte Bestemmelser. At nu disse
Bestemmelser i philosophisk Form maae blive dialektiske, medens
de i mathematisk Form hverken ere eller kunne være det,
viser en charakteristisk Forskjel imellem Mathematik og Philosophie,
men ingenlunde nogen Strid. Intet Begreb kan,
hvor exact det end defineres, holde sig i en mod de øvrige
Begreber ligegyldig Væren: Begrebernes Gyldighed sees i
deres Overgange; men de Overgange, den exacte Analyse netop
gjør ved Operationen, gjør den Operationsbegreberne gjennemtrængende
Tænkning ved Dialektiken“ (Mathem. og Dial. S. 57.
Jvfr. S. 59--60). Stilles Endeligt umiddelbart imod Uendeligt,
da faaer den sunde Forstand uvilkaarlig et Indtryk,
som var det det Endelige, der i sin Mangfoldighed af Led
og Forholdsbestemmelser indeholdt Tilværelsens Rigdom, medens
det Uendelige derimod er et bundløst Intet, øde og
tomt, som det tomme Rum. Da nu Tanken kun kan bestemme,
hvad der i sig selv er bestemmeligt, og kun det er
bestemmeligt, som paa en eller anden Maade lader sig definere,
og kun det lader sig definere, som lader sig begrændse, og det
Begrændsede er et Endeligt: saa maa det bestandig forekomme
Forstanden, som var det kun det Endelige, der lod sig begribe,
medens det Uendelige, det Ubegrændsede, egentlig slet
ikke var til at begribe, men, ligerviis som de Gamles Hyle
(τὸ ἄπειρον), at erkjende ved en Slags uægte Tænkning
(λογισμὸς νόθος). Hvor denne Misforstaaelse har sat sig
fast, kan der naturligviis ikke være Tale om nogen nærmere
Indgaaen paa Uendelighedsproblemet; for at indsee Problemets
philosophiske Betydning, maa man absolut have idetmindste
en Anelse om det Endeliges og det Uendeliges
% --- p. 50 ---
% degrees are printed as a superscript 0 (not a ring) here and on p. 59;
% transcribed as ^{0} throughout.
dialektiske Vexelforhold. „En Vinkel paa $90^{0}$ kan deles i
uendelig mange uendelig smaa Vinkler, en Vinkel paa $1^{0}$
ligeledes; en Qvadratmiil kan deles i uendelig mange uendelig
smaa Qvadrater, en Qvadrattomme ligeledes: hvad følger
heraf? Lægges Eftertrykket paa det Endeliges Opløsning i
det Uendelige, da synes vistnok Følgen at være, at her ikke
længere er Forskjel imellem Vinklen paa $90^{0}$ og Vinklen paa
$1^{0}$, at Delen bliver liig det Hele, at Grændsen med andre
Ord er udslettet. Men lægges saa Eftertrykket igjen paa
det Uendeliges Tilværen i og ved det Endelige, da bliver
der ikke blot Forskjel imellem Endeligt og Endeligt, imellem
Delen og det Hele, men her viser sig med det Samme en
bestemt Forskjel imellem Uendeligt og Uendeligt saavel med
Hensyn til det uendelig Lille, som til det uendelig Store.
Det ene uendelige Antal uendelig smaa Vinkler kan da være
ligesaa forskjelligt fra det andet, som $1^{0}$ fra $90^{0}$; det ene
uendelige Antal uendelig smaa Qvadrater ligesaa forskjelligt
fra det andet, som en Qvadrattomme fra en Qvadratmiil.
At Grændsen i det Endelige ophæves, betyder, at den sættes
i det Uendelige; at det Uendelige i sig optager Grændsen,
betyder, at det Endelige sættes. Den dialektiske Eenhed viser
sig deri, at Grændsen ligesaameget tilhører det Uendelige som
det Endelige“ (Mathm. og Dial. S. 33). Den Grundtanke,
at det Uendelige ikke blot udsletter Endelighedens Grændser,
men tillige optager de ophævede Grændsebestemmelser som
articulerende Momenter i sin egen Begrebsudvikling, er fra
Operationernes Standpunkt tilstrækkelig betegnet derved, at
det Uendelige analyseres.

% The printed head reads "b)." — a second b) after "b) Uendelighedsproblemet"
% p. 45 — where the Indhold gives "c) Det uendeliges Analyse ... 50".
% Head kept as printed, Indhold's reading recorded here.
\textbf{b). Det Uendeliges Analyse.}
\addcontentsline{toc}{subsection}{b). Det Uendeliges Analyse}

I den Opgave at „analysere det Uendelige“ have altsaa
Mathematiken og Philosophien et omfattende Fælledsproblem;
at Problemet imidlertid, idet det optages af Philosophien,
maa undergaae en Forvandling, at den philosophiske Analyse
% --- p. 51 ---
hverken kan eller skal gjennemføres paa samme Maade, som
den mathematiske, er en Selvfølge: Et er det nemlig at uddanne
et Begreb til Brug ved derunder hørende Operationer,
et Andet at gjennemreflectere Operationerne for at tydeliggjøre
det deri virksomme Begreb. „Hvor flere Videnskaber
behandle det Samme, vil enhver gjøre Sit tydeligt; men de
ville aldrig gjøre det Samme tydeligt paa samme Maade.
Hvad en Udvikling ifølge sit Øiemed maa bestemme med
største Nøiagtighed, tør en anden efter sit Øiemed kun antyde
i flygtige svævende Træk, just fordi den skal lade Nøiagtighedens
Eftertryk falde paa andre Punkter.“ (Begreb. Nøiagt.
S. 51--52).

Da Mathematiken som Fagvidenskab umiddelbart seer
paa sit Eget, kan det vistnok i Reglen ikke forudsættes, at
Mathematikerne som Fagmænd just skulde have enten særlig
Interesse af eller synderlig Færdighed i at følge Philosophiens
dialektiserende Analyser; men vel kan det forudsættes, ja i
Videnskabelighedens Navn ligefrem fordres, at den Fagmand,
der føler sig kaldet til at fælde Dom over det i en philosophisk
Fremgangsmaade Eiendommelige, dog idetmindste værdiger
Sagen saamegen Opmærksomhed, at han fastholder de
allervæsentligste Udgangspunkter, følger Tankegangen og veed,
hvorom Talen er. Naar det, for at anføre et Par Exempler,
udtrykkelig skal udhæves, hvilken Forskjel det gjør, om det er
Kategorierne: \emph{Deling} og \emph{Sammensætning}, eller:
\emph{Sammenstykning} og \emph{Bevægelse}, der modsættes hinanden, og
dette saa oplyses ved en Sammenligning mellem den brudte
og den krumme Linie saaledes: „For den brudte Linie gjælder
en Sammensætning, thi dens Dele ere endelige; den er et
Stykværk, som i Eet og Alt tilhører Endeligheden. For den
% printer's defect: "Sammensætuing" (turned n) for "Sammensætning" —
% verified on the image; transcribed as printed.
krumme Linie gjælder ingensomhelst endelig Sammensætuing,
ingen, end ikke den allerfineste, Sammenstykning; den skylder
kun Bevægelsen sin Tilblivelse, thi den tilhører Uendeligheden“
(Phil. og Mathm. S. 76); mon da Nogen, der fatter,
% --- p. 52 ---
hvorom Talen er, ikke strax skulde indsee Umuligheden af at
charakterisere \emph{Sammenstykningens} Kategorie, uden at forudsætte
de endelige Deles umiddelbart givne Tilværelse? Skal
her reflecteres paa, hvorledes alle de Linier, hvoraf den brudte
Linie sammensættes, de Dele, hvoraf et Hele vorder, selv igjen
have deres Vorden, da sees alene paa Vorden, ikke paa
Sammenstykning. Vistnok kan man ogsaa her, forsaavidt det
ikke gjælder om at klare en Tanke, men kun om at anbringe
„en Kritik“, paa Fagets Vegne anbringe Følgende: „den rette
og brudte Linie skylder ikke mindre og ikke mere end den
krumme Linie Bevægelsen sin Tilblivelse; den krumme Linie
er sammensat mathematisk af endelige Længder, saavel som
den brudte og den rette.“ Men hvad vinder Videnskaben
ved en saadan Kritik? At lade Alt løbe sammen, Dannelse
ved Sammenstykning og Tilblivelse ved den jævne Overgang, at
lade Alt komme ud paa Eet: Discretion og Continuitet, Endeligt
og Uendeligt, Brudt og Ubrudt, Krumt og Lige, Ret og
Kroget, er dog vist ligesaalidt mathematisk, som philosophisk!

Uvidenhed og Tankeløshed ere ingenlunde altid det Samme;
man kan være uvidende, uden derfor at være tankeløs; man
kan være i endelig Forstand vidende, og dog i paafaldende
Grad tankeløs. Exemplerne ere mangfoldige. Der spørges
i en philosophisk Undersøgelse, hvorledes man egentlig skal
tænke sig Liniens Oprindelse i Forhold til Punktet; thi at
man i Grunden altid har Linien, naar man først har Fladen,
og Fladen, naar man først har Legemet, forstaaer sig af sig
selv. „I Forhold til Linien som det Positive er Punktet,
Liniens Grændse, det Negative; naar Linien er en Forestilling,
er Punktet en Tanke. I Forhold til Fladen som
det Positive er Linien, Fladens Grændse, det Negative.
Fladen forestilles, naar Linien tænkes.“ (Phil. og Math. S.
65--67). Naar saaledes Spørgsmaalet er, hvorledes Liniens
Oprindelse egentlig skal tænkes, da det, at lade Linien fremkomme
ved „Sammensætning af endelige Længder“, kun er at
% --- p. 53 ---
forudsætte, hvad først skulde vises, eftersom de endelige Længder
jo selv ere Linier, og det, at lade Liniens Tilblivelse beroe
paa Flader, heller ikke gaaer an, da Fladens egen Oprindelse
nødvendigviis forudsætter Linien; naar nu dette udtrykkelig
er Spørgsmaalet, og saa en Mand af Faget, en
Sagkyndig, der kan alle Tilfældene, baade med Linier og
Flader, med Skjæring og uden Skjæring, møder med den
Oplysning: „\emph{mathematisk opstaae krumme Linier
ogsaa ved Fladers indbyrdes Skjæring}“, saa faaer
man ikke blot en mathematisk Oplysning, men tillige en
æsthetisk Opmuntring.

\emph{α) Grændser.} Forskjellen imellem den exacte og den
dialektiske Tænkning er især kjendelig paa en forskjellig Brug
af Grændsen. „Enten Noget eller Intet“, „enten Positivt
eller Negativt“, „enten Endeligt eller Uendeligt“, „enten et
Opnaaeligt eller et Uopnaaeligt“: slige Sondringsdomme forudsætte
en skarp Betoning af det i Grændsen Adskillende.
Den skarpt adskillende Grændse, den Grændse, der holder de
Adskilte saa aldeles udenfor hinanden, at det Ene maa benegtes,
naar det Andet bekræftes, er \emph{qvalitativ}. Men jo
mere den contradictoriske Modsætning skjærpes, jo mere udelukkende
det Qvalitative gjælder, desto mere iøinefaldende bliver
det, at Grændsen ved at være det Adskillende netop er det
Forenende. Sætter man saaledes: enten $+a$ eller $-a$,
saa holdes det Positive vistnok herved udenfor det Negative;
men da Adskillelsen er saa skarp, at her slet ikke er noget
Mellemværende, bringes det Positive og det Negative jo netop
i den adskillende Grændse, Nulpunktet, til umiddelbar Berøring.
Nulpunktet er altsaa Foreningspunkt for det Positive
og det Negative, just fordi det er deres Sondringspunkt.
Sætte vi nu istedetfor det Positive det Endelige, istedetfor det
Negative det Uendelige, altsaa: enten Endeligt eller Uendeligt,
saa gjentager sig det Samme. Ved at stirre paa Negationen,
det i Grændsen Adskillende, faaer man et Indtryk,
% --- p. 54 ---
som om det Endelige og det Uendelige holdtes adskilte ved
et svælgende Dyb: „enten Endeligt eller Uendeligt! man kan
altsaa blive ved med det Endelige, saalænge man vil, et
Endeligt er og bliver et Endeligt, aldrig et Uendeligt“. Ved
at besinde sig paa det i Grændsen Forenende indseer man
derimod let, at de saaledes Adskilte alligevel maae ligge hinanden
uendelig nær: „enten Endeligt eller Uendeligt! altsaa,
naar man ikke har det Endelige, saa har man strax uden
videre det Uendelige“; hvor det Endelige ophører, begynder det
Uendelige; Grændsen, Negationen, er altsaa selv det logiske
Berøringspunkt, det Punkt, hvori de Adskilte træffe sammen:
her er intet Mellemværende. „Men naar det Endelige og det
Uendelige paa denne Maade ere hinandens Grændser, maa
Grændsebestemmelsen blive dialektisk. Har nemlig det Uendelige
sin Grændse i det Endelige, saa har det jo en Endelighedsgrændse,
en endelig Grændse, og er selv endeligt. Naaer
det Endelige omvendt først sin Grændse i det Uendelige, da
har det jo ingen endelig Grændse, og er følgelig uendeligt“.
(Phil. og Math. S. 22). Det Uendelige er et Ubegrændset;
hvorledes kan nu det, der slet ingen Grændser har, selv være
Grændse? Det Endelige, der overalt har Grændser, forsaavidt
det ene Endelige begrændser det Andet, kan vel med
Rette siges at være et i sig Begrændset; men hvorledes kan
nu det i sig Begrændsede have et Ubegrændset til Grændse?
Eller gives her maaskee tre Slags Grændser: en Grændse,
der gjælder for Endeligt, ikke for Uendeligt, en Grændse, der
gjælder for Uendeligt, ikke for Endeligt, og saa en Grændse,
der er fælleds for Endeligt og Uendeligt? Men naar dette
antages, hvorledes forholder det sig da med enhver af disse
Grændser i Modsætning til det, hvis Grændse den er?

Hvad der med Hensyn paa den mathematiske Udtryksmaade
let giver Grændsebegrebet et tvetydigt Anstrøg,
er den Omstændighed, at baade Størrelser og Størrelsesnegationer
sættes som Grændser, og at $0$ og $\infty$
% --- p. 55 ---
\setcounter{footnote}{0}
ophæve hinanden. Ved $0$ bliver den endelige Grændse sat,
ved $\infty$ bliver enhver endelig Grændse hævet; en nærmere
Betragtning vil da let gjøre det indlysende, at begge Negationer
høre med til en fuldstændig Opfattelse af den qvantitative
Grændse. Uden Qvalitet ingen Qvantitet, uden Nulgrændse
ingen Størrelse; men uden en Qvantiteren, der igjen
ophæver den umiddelbare Nulgrændse, og forvandler det Qvalitative
til et i Formerelse og Formindskelse svindende Moment,
gives der heller ingen Størrelse\footnote{„Alt hvad der kan tænkes
forøget eller formindsket, og som kan opløses i eensartede Dele, kaldes
Størrelse (\textit{quantum}) og siges at besidde Storhed
(\textit{quantitas}).“ \ Ramus.}. For at indsee dette, behøver
man blot at sammenstille $0$, $1$ og $\infty$, thi medens $0$
betegner den satte Grændse, det abstract Qvalitative, og $\infty$ alle
mulige Grændsers Ophævelse, den rene Qvantiteren, er Taleenheden
selv betinget af begge. Sige vi, at en Størrelse $x$
falder t.~Ex. imellem Grændserne $0$ og $1$, da er det egentlig
ikke Størrelsen $1$, men det Nulpunkt, hvormed $1$ overskrides,
der udgjør Grændsen, nemlig Slutningsgrændsen. Kalde vi
Slutningsgrændsen $0_2$, Begyndelsesgrændsen $0_1$, saa betegner
$1$, at Afstanden imellem de to Grændser er sat som en i al
Almindelighed bestemt d.~e. begrændset Størrelse, som Afstandseenhed.
At $x$, forsaavidt det falder imellem $0$ og $1$,
kan med uforandret Begyndelsesgrændse variere ved utallige
Slutningsgrændser, eller, som det hedder, have uendelig mange
Værdier, lader sig da udtrykke saaledes: $\frac{1}{\infty} = 0$, kun ved
at dele Eenheden i det Uendelige naaes den virkelige Grændse;
$\frac{1}{0} = \infty$, indenfor enhver endelig Eenhed falder et uendeligt
Antal Grændser; eller: $1 = \infty. 0$, den endelige Eenhed
fremkommer derved, at Grændsen sættes og hæves uendelig
mange Gange.
% --- p. 56 ---

Det bestemte Qvantum forudsætter en Qvantiteren indenfor
qvalitativt bestemmende Nulgrændser. Ifølge de bestemmende
Negationer ere de endelige Størrelser saa forskjellige
indbyrdes, at Forskjellen ikke blot kan angives, men endog
angives exact; ikke nok, at $a$ er større end $b$, her spørges:
hvormeget større? Men da Størrelsesforskjellen udelukkende
beroer paa fastsatte Grændser, og ingen endelig Størrelsesgrændse
staaer saa fast, at den jo, hvad $\infty$ minder om, igjen
kan hæves, saa kunne de endelige Størrelser forandres, idet
Grændserne forandres: enhver Størrelse kan som Størrelse
formeres og formindskes i det Uendelige.

Den dialektiske Eenhed af $0$ og $\infty$, der har sin Bestaaen
i den endelige Størrelse, $1 = \infty. 0$, udvikler sig exact,
idet der idelig vexles mellem Fastsætten af Grændser og Overskriden
af Grændser; saaledes danner der sig en afgjørende Modsætning
imellem \emph{faste} og \emph{flydende}, \emph{værende} og
\emph{vordende} Grændser. Den umiddelbare Nulgrændse har ingen Vorden,
den bliver til, idet den sættes; men da enhver Grændse kan
overskrides, og Grændseforandringen vedblive saaledes, at den
ene Grændse uophørlig falder enten indenfor eller udenfor
den anden, kan Grændsen befinde sig i en Vorden. Grændsen
for Grændsens Vorden, Slutningsgrændsen, lader sig igjen
betragte deels som sat, deels som hævet, sat nemlig, forsaavidt
Størrelsen er endelig, hævet, forsaavidt en endelig
Størrelse kan i sin Almindelighed være saa stor og saa lille,
man vil: idet en bestemt Slutningsgrændse sættes, ere utallige
forudgaaende Slutningsgrændser allerede overskredne, og forsaavidt
hævede.

\emph{β) Overgange.} I Grændsen reflecteres det Begrændsede.
At tænke Grændsen uden at underforstaae et
derved Begrændset, er ligesaa umuligt, som at fatte Betydningen
af en Negtelse uden at tænke paa det, som derved
% printer's defect: "ogsao" for "ogsaa" — final o verified on the image
% (both OCR witnesses agree); transcribed as printed.
bliver negtet. Den umiddelbare Nulgrændse bliver da ogsao
ligefrem sat under den Forudsætning, at der i Anskuelsen er
% --- p. 57 ---
et Givet, hvis Grændse den er. At saadanne Grændser som
Punkter, Linier og Flader, uagtet de i og for sig ere rene
Negationer, dog \emph{anskues} i positiv Form: Punktet som
Liniens mindste Deel, Linien som Fladens Rand, Fladen
som en Side af Legemet, er et slaaende Beviis for, at
Grændsen gaaer i Eet med det Begrændsede. Men naar
Grændsen nu ikke destomindre adskiller sig fra det Begrændsede
som Negationen fra det Positive, naar denne Adskillelse
er i lige Grad gyldig for Tanken og Anskuelsen: saa
indeholdes heri den Opgave at følge det, som skal begrændses,
til Grændsen. Hvorledes naaer man nu fra Grændse til
Grændse? Ved en Bevægelse! Her er oprindelig ingen anden
Udvei, intet andet Middel. Forsøger man i Tanken at afsætte
t.~Ex. paa en Linie to Grændsepunkter, $A$ og $B$, i en
vis Afstand, da maa man, om end ubevidst, foretage en Bevægelse,
idet man, for at anskue den vilkaarlig valgte Afstand,
uvilkaarlig lader Blikket gjennemløbe Veien fra $A$ til $B$.
Blikket maa fæste sig, hvor det vil, det møder paa sin Vei
kun Grændse, thi hvert Punkt i Linien er i og for sig et
Grændsepunkt. Idet Blikket bevæger sig hen ad Linien, seer
det ud, som om det kun fulgte den sig bevægende Grændse,
det sig bevægende Punkt; herpaa grunder sig den mathematiske
Forklaring, at Linien fremkommer ved Punktets Bevægelse.
Da denne Forklaring væsentlig beroer paa et optisk
Bedrag, er det intet Under, at den falder noget tvetydig ud.
Hvorledes kan det rene Punkt, Grændsen, det abstract Negative
bevæge sig? Thi at Anskuelsen, der i sin Umiddelbarhed forudsætter
det Positive, kan, idet Tanken sætter Negationen,
overskride enhver Grændse, gaae fra Grændse til Grændse,
er endnu ingenlunde tilstrækkeligt til at forklare det mathematiske
Mirakel, at Punktet, den rene Negation, skulde kunne
bevæge sig, forandre Sted i Rummet, uagtet det ikke indtager noget
Rum, eller ligesom strække sig, og danne noget saa Positivt,
som en Længdeeenhed, en Linie. Kunde Linien fremkomme
% --- p. 58 ---
\setcounter{footnote}{0}
ved en Bevægelse af det rene abstracte Nulpunkt, da maatte
den ogsaa kunne tænkes sammensat af lutter Nulpunkter, af
$0 + 0\,....\,+ 0$\footnote{At denne Sætning ingenlunde staaer
i Strid med Ligningen $\frac{a}{\infty} = 0$ opfattet som
Mulighedsudtryk, indlyser, naar man blot erindrer, at Muligheden just
beroer derpaa, at $0$ og $\infty$ opfattes, ikke som tomme Negationer,
men som Symboler paa den størst og mindst mulige Størrelse.}.
At sætte Grændse umiddelbart ved
Grændse, er umuligt, thi idet der abstraheres fra det Begrændsede,
falde de to Grændser sammen; Herbarts „die
% FOOTNOTE MARK: two footnotes on this page; the first is printed *) as
% usual, the second (below) is printed **) in text and at the note.
starre Linie“ er et ligesaa logisk som mathematisk
Monstrum\footnote{Jvfr. Den Herbartske Philosophie af P.~S.~V.
Heegaard. Kbhvn. 1861. S. 21--45.};
kun i Reflexion af et Positivt kan det Negative komme udenfor
sig selv, kun formedelst et Begrændset kan Grændse følge
paa Grændse.

Imellem to paa hinanden følgende Grændser er her
altsaa et Begrændset. Dersom nu dette Begrændsede faldt
umiddelbart sammen med Grændserne, saa at her intet, slet
intet Mellemværende var, maatte Grændserne selv falde sammen
og kunde ikke følge paa hinanden; faldt det Begrændsede som
et umiddelbart Tilværende imellem de derfra forskjellige
Grændser, saa var her en endelig Afstand, altsaa utallige
Grændser mellem to paa hinanden umiddelbart følgende
Grændser, hvilket er en Modsigelse: i den dialektiske Eenhed
af Grændsen og det Begrændsede er al Bevægelse begrundet.
Formedelst Grændsen er Bevægelsen Overgang; den rene
Grændse er forsaavidt altid Overgangsform. Vistnok gjør
det en væsentlig Forskjel, om Overgangen anskues i Rummet
eller blot forestilles som arithmetisk Overgang fra Endeligt
til $0$, fra Endeligt til $\infty$; men dette staaer fast: uden
Reflexer af Rum og Tid ingen Bevægelse, uden Bevægelse
% --- p. 59 ---
ingen Qvantiteren, uden Qvantiteren intet Qvantum, uden
Qvantum ingen Mathematik. At udelukke Bevægelseskategorien
fra den rene Mathematik og endda beholde Kategorierne:
„Grændse“, „Overgang“, „Tilnærmelse“, er en Modsigelse.
Hvorledes Kategorierne i Mathematiken vælges, betones og
anvendes, forsaavidt det technisk gjælder om den meest hensigtssvarende
Udførelse af exacte Operationer, maa naturligviis
afgjøres i Mathematiken selv; men Operationsbegrebernes
logiske Sammenhæng og dialektiske Overgange undersøges i
Philosophien. Hændes det saa --- thi det kan jo hændes ---
at en Fagmand, en Routineret, hvis Begrebsmechanik forlængst
er gaaet i Eet med hans Operationsmechanik, i en
saadan Anledning bliver høirøstet, og ved høirøstet Anvendelse
af det saraceniske Dilemma udtrykkelig fordrer, at Philosophien
skal opgive Dialektiken og ubetinget hylde Begrebsmechaniken:
da kan en saadan Hændelse dog ikke videre afbryde,
men snarere fremskynde den dialektiske Udvikling.

Havde Nogen, for at bevise den Sætning, at „Overgangen
fra Endeligt til Endeligt, falder sammen med Overgangen
fra Endeligt til Uendeligt“, beraabt sig paa en Ligning
som: $\mathrm{tg}\,90^{0} - \mathrm{tg}\,(90^{0} \pm \varepsilon)
= 0$, idet $\varepsilon$ betegner
et Endelighedspunkt, da havde en Mathematiker her
ganske vist havt Anledning til at gjøre opmærksom paa Feiltagelsen,
og imod den falske Ligning at stille den sande:
$\mathrm{tg}\,90^{0} - \mathrm{tg}\,(90^{0} \pm \varepsilon) =
\infty$. Men naar der nu i en
philosophisk Analyse (Jvfr. Phil. og Math. S. 23--24) aldeles
ikke er Tale om nogen særlig Bestemmelse af Forskjellen
imellem $\mathrm{tg}\,90^{0}$ og $\mathrm{tg}\,(90^{0} \pm
\varepsilon)$, men kun om den vel
bekjendte Kjendsgjerning, at $\mathrm{tg}\,x$ gaaer over fra at være
endelig til at være uendelig, idet $x$ gaaer over fra en
% printer's defect: "Mathemathikeren" (Mathe-mathikeren over the line
% break) for "Mathematikeren" — verified on the image; as printed.
Værdi, der er mindre end $90^{0}$, til $90^{0}$, hvor vil Mathemathikeren
% printed "?." — question mark followed by full stop, as printed.
(Jvfr. Begr. Nøiagt. S. 47, „Afregning“, gjennems.
S. 7) saa hen med sit $\mathrm{tg}\,90^{0} -
\mathrm{tg}\,(90^{0} \pm \varepsilon)$?. Er $\varepsilon$
af en, om end iøvrigt nok saa lille endelig Værdi, altsaa et
% (p. 59 ends mid-sentence; the sentence continues on p. 60)

% --- p. 60 ---
\setcounter{footnote}{0}
Endelighedspunkt\footnote{Man har t.~Ex. $\mathrm{tg}\,90^{0} = \infty$,
man har $\mathrm{tg}\,(90^{0} \pm \varepsilon)$ endelig; hvorledes kommer
man nu fra $\mathrm{tg}\,90^{0}$ til $\mathrm{tg}\,(90^{0} \pm
\varepsilon)$? Aabenbart derved, at $\varepsilon$ som Deel af Buen
gjennemløbes continuerligt; under denne Bevægelse skeer det saa, at
Tangens gaaer over fra at være endelig til at være uendelig. Hvorledes
er nu en saadan Overgang egentlig at tænke? „Den kan“, paastaaer man,
„kun skee ved et Spring“. Hvad er det da, her overspringes? „Et
Endeligt!“ Men saa er her jo kun en endelig Afstand mellem Endeligt og
Uendeligt. Skal nu desuagtet Overgangen fra Endeligt til Uendeligt skee
ved et Spring, saa maa Overgangen fra Endeligt til Endeligt jo ogsaa,
naar det gjælder en endelig Afstand, skee ved et Spring; men dersom
Overgangen gjennem Endeligt kun kan skee ved et Spring, saa er jo al
continuerlig Overgang, og dermed al mulig Bevægelse, ja al mulig
Overgang, negtet. „Et Uendeligt!“ Men at komme til det Uendelige ved at
overspringe det, er dog vel en utaalelig og meningsløs Modsigelse.
„Hverken et Endeligt eller et Uendeligt!“ Men hvad er saa det for en
Størrelse, der hverken er endelig eller uendelig? den størst mulige
endelige Størrelse kan det dog umulig være, thi den størst mulige
endelige Størrelse har jo netop den Egenskab, at den ved at være
\emph{hverken} endelig \emph{ei heller} uendelig netop er paa een Gang
\emph{baade} endelig og uendelig. „Intet!“ Men hvad er det saa for et
Slags Spring, hvorved Intet overspringes? (Jvfr. Begr. Nøiagtighed.
S. 45. Den gjennemsete „Afregning“ S. 7).}, saa er
$\mathrm{tg}\,(90^{0} \pm \varepsilon)$ jo endelig, og Overgangen
fra Endeligt ($90^{0}$) gjennem et Endelighedspunkt ($\varepsilon$)
til Endeligt ($90^{0} \pm \varepsilon$) falder altsaa sammen med en
Overgang fra Uendeligt ($\mathrm{tg}\,90^{0}$) til Endeligt
($\mathrm{tg}\,(90^{0} \pm \varepsilon)$); er $\varepsilon$ uendelig
lille, altsaa et Uendelighedspunkt, saa er $\mathrm{tg}\,(90^{0} \pm
\varepsilon)$ uendelig, og Overgangen fra Endeligt gjennem et
Uendelighedspunkt til Endeligt (thi $90^{0} \pm \varepsilon$ er dog vel
ligesaavel af en endelig Værdi som $90^{0}$) falder altsaa sammen med
Overgangen fra Uendeligt ($\mathrm{tg}\,90^{0} = \infty$) til Uendeligt
($\mathrm{tg}\,(90^{0} \pm \varepsilon) = \infty$). Hvor stor
Differensen imellem $\infty$ og $\infty$ forøvrigt
% --- p. 61 ---
\setcounter{footnote}{0}
er, vedkommer ikke Overgangen som Overgang\footnote{Hvor uendelig
dialektisk Størrelsesbegrebet er, indlyser, naar man t.~Ex. spørger:
hvor stor er den størst mulige endelige Størrelse? eller: hvor lille er
den mindst mulige endelige Størrelse? Den i disse Spørgsmaal udtalte
Fordring indeholder en væsentlig Modsigelse. Der gives ingen endelig
Størrelse saa stor, at jo dens Grændse kan overskrides, saa at en endnu
større altid er mulig; heller ikke gives der nogen endelig Størrelse
saa lille, at en mindre jo altid er mulig; den størst mulige endelige
Størrelse er, med andre Ord, den uendelig store endelige, og den mindst
mulige endelige den uendelig lille endelige; ogsaa her viser det sig,
at det Uendelige ligger, ikke i en Umulighed, men i en Mulighed.
Hvorledes kan nu Modsigelsen hæves? Ved at see det i Størrelsens
Grændse Dobbeltsidige, altsaa ved at opfatte Grændsen deels som
værende og deels som vordende. Størrelsen med værende Grændse, den
værende Størrelse, maa være saa stor og saa lille, den vil, den er
altid endelig; med vordende Grændse bevæger Størrelsen, den vordende
Størrelse, sig derimod enten imod det Større, og fremstiller sig da
under den \emph{uendelige} Vorden som det uendelig Store, eller imod
det Mindre, og opfattes da under den \emph{uendelige} Vorden som det
% printer's defect: "nendelig" for "uendelig" (turned sort) — verified on
% the image at 3x; both OCR witnesses agree; transcribed as printed.
uendelig Lille: det nendelig Store og det uendelig Lille ere altsaa
Størrelser, hvis værende Grændse opløses i Vorden som i uophørlig Voxen
og uophørlig Aftagen. Hvad er Betegnelsen for den størst mulige
% printed with a full stop after "Størrelse", though the sentence is a
% question (the parallel clause that follows has "?"); as printed.
endelige Størrelse. $\mathrm{Lim}\ x = \infty$! og for den mindst
mulige? $\mathrm{lim}\ x = 0$. Den Grændse, der adskiller det Endelige
fra $\infty$ og fra $0$, er, saa at sige, ikke et Haar bredere, end
den, der adskiller det Endelige fra det Endelige. „Ligger det ene
Endelighedspunkt det andet uendelig nær, da kan nu det Samme siges om
$0$ og Endeligt, om $\infty$ og Endeligt“. (Philos. og Math. S. 24).
Afviser man Dialektiken, forsøger man ved kritisk Sondring: „enten et
Endeligt eller et Uendeligt“, at fjerne Modsigelsen, bliver den først
for Alvor paatrængende; forsøger man saa ved allehaande techniske
Vendinger, Spring og Kunstgreb at tilhylle den overalt lurende
Modsigelse, vil man ved engang imellem at indlade sig med det Logiske
let kunne geraade udi adskillige komiske Forlegenheder.
% (this footnote runs over: the paragraph below is printed, unmarked, at
% the foot of p. 62 and concludes at the top of the p. 63 notes)
\par Fastholder man Dilemmaet: „enten er $\varepsilon$ endeligt eller
det er $0$“, eensidig kritisk, altsaa udialektisk, altsaa:
„$\varepsilon - 0$ enten endeligt eller $0$“, saa tør man da ligesaa
lidt om $0$, som om $\infty$, sige, at det ligger det Endelige uendelig
nær. Vil Mathematikeren hertil svare: „ganske rigtigt! $0$ ligger
heller ikke Endeligt uendelig nær“, saa maa man igien spørge: hvad
ligger der da imellem $0$ og Endeligt? Det uendelig Lille tør
Mathematikeren ikke anføre, thi han har jo udtrykkelig selv sagt:
„enten er $\varepsilon$ $0$ eller endeligt“. Der ligger altsaa Intet
imellem $0$ og Endeligt, og dog skulde $0$ ikke ligge det Endelige
uendelig nær! Paa tilsvarende Maade forholder det sig med $\infty$ og
Endeligt. Udialektisk fastholder man det Dilemma: enten er $x$ endeligt
eller uendeligt, \textit{tertium non datur}; altsaa $\infty - x$ enten
$\infty$ eller endeligt eller $0$; man tør da ligesaa lidt om $\infty$
som om $0$ sige, at det ligger det Endelige uendelig nær. Men hvad
ligger der da imellem Endeligt og $\infty$? Vil Mathematikeren svare:
„der ligger $\infty$!“ saa svarer han jo netop det Modsatte af, hvad
han egentlig tænkte. Han tænkte nemlig, at der ifølge Ligningen
$\infty - a = \infty$, for $x = a$, absolut maatte være en uendelig
Afstand imellem det Endelige og $\infty$, og see, saa forvandler den
formeentlige Afstand, det \emph{Mellemværende}, sig just til det, der
igjennem Afstanden skulde naaes, nemlig det Uendelige selv; med andre
Ord: det Endeliges yderste Grændse kan umuligt overskrides, uden at det
Uendelige naaes. Det, der ved „den uendelige Afstand“ skulde holdes
uendelig langt borte fra det Endelige, kommer altsaa dog dette Endelige
uendelig nær. „Enten Endeligt eller Uendeligt!“ Men saa ligger her jo
Intet, slet Intet imellem Endeligt og Uendeligt; og dog skal
$\infty - a = \infty$, $\mathrm{tg}\,90^{0} - \mathrm{tg}\,(90^{0} \pm
\varepsilon) = \frac{1}{\cos 90^{0}} = \infty$ \emph{være et Beviis, et
fuldgyldigt mathematisk Beviis for, at der imellem Endeligt og
Uendeligt er en uendelig Afstand!} Skulde en mathematisk Kritiker, ved
at buldre med en saadan Logik, virkelig vinde Tilhængere, saa maatte
man da rigtignok paa Logikens Vegne anvende det gamle Ord: \textit{un
sot trouve toujours un plus sot, qui l'admire}.}. Den Indvending, at
Udtrykket: „falde sammen med“, kun betegner „et tilfældigt
Sammentræf“, er ubeføiet; thi Tangensforandringer
% --- p. 62 ---
% (no new footnote mark on this page; its foot carries the unmarked
% run-over of the p. 61 note, transcribed at the p. 61 anchor above)
foregaae ikke blot \emph{samtidigt} med, men ere \emph{nødvendige
Følger} af tilsvarende Vinkelforandringer: at forvexle Conseqvensens
Fremgang af Forudsætningen med tilfældigt Sammentræf af noget
Samtidigt, forraader kun Tankeløshed. Ifølge $\frac{1}{a-x} = y$ ville
Forandringerne i $y$ jo Punkt for Punkt være afhængige af
Forandringerne i $x$; naar nu $y$, idet $x = a$, gaaer over fra at være
endelig til at være uendelig, hvo vil da i en saadan Overgang see et
blot udvortes tilfældigt Sammentræf med den $x$ tillagte Værdi? Er
$\varepsilon = a - x$ den mindst mulige endelige Differens, saa er det
Vendepunkt, da $y$ gaaer over fra Endeligt til Uendeligt, ogsaa
uendelig nær; men heraf følger ingenlunde, at Forskjellen imellem
Endeligt og $\infty$ skulde lade sig udstykke enten i endelige eller i
uendelige Subtractionsdifferenser. Skulle Continuitetsspørgsmaalene
afgjøres ved Beviser af \emph{umiddelbare} Subtractionsdifferenser,
maae Beviserne rigtignok blive derefter. Ligningen:
% the minus in the equation below is printed with the Danish
% division-style minus sign ÷ (on p. 63 the same equation is printed
% with —); transcribed as printed, \div here, - there.
$\mathrm{tg}\,90^{0} \div \mathrm{tg}\,(90^{0} \pm \varepsilon) =
\frac{1}{\cos 90^{0}} = \infty$
% --- p. 63 ---
\setcounter{footnote}{0}
gjælder jo ogsaa for $\varepsilon = 0$. Ved at gjøre den umiddelbare
Subtractionsdifferens til Maalestok for det Continuerede kommer man da
til det i logisk Forstand saa overmaade skjønne Resultat, at, efterdi
$\mathrm{tg}\,90^{0} - \mathrm{tg}\,(90^{0} \pm 0) =
\frac{1}{\cos 90^{0}} = \infty$, saa maa der nødvendigviis være en
uendelig Afstand imellem $\mathrm{tg}\,90^{0}$ og $\mathrm{tg}\,90^{0}$,
eller, hvad her er det Samme, $\mathrm{tg}\,90^{0}$ være uendelig langt
% (this footnote runs over: its final clause and the further unmarked
% paragraph "Endnu en Bemærkning ..." are printed at the foot of p. 64)
borte fra sig selv\footnote{Den Udflugt, at $\infty$ ikke er en
Størrelse, siger Intet; var $\infty$ ikke en Størrelse, hvorledes kunde
man da have $\infty - \infty = n$, en bestemt \emph{Størrelsesforskjel}
imellem to \emph{Ikke-Størrelser}? Hvad det vil sige, at $\infty$ er en
Størrelse, nemlig en vordende Størrelse, sees t.~Ex. af $\frac{1}{1-3}
= 1 + 3 + 3^{2} + 3^{3}\,....\,+ 3^{n} + \frac{3^{n+1}}{1-3} =
\infty - \infty = -\,\frac{1}{2}$, for $n = \infty$. Den Udflugt, at
$\mathrm{tg}\,90^{0} - \mathrm{tg}\,(90^{0} \pm \varepsilon)$ jo
% only "Functions" is letterspaced in the print; the word is set
% "Functions differens" with the gap of the letterspacing.
egentlig er en \emph{Functions}differens, hjælper heller ikke; thi
dersom Functionsdifferensen skal faae Betydning i Modsætning til den
umiddelbare Subtractionsdifferens, maa $\varepsilon$ forvandles til
Element, og derved Vorden, continuerlig Forandring, Overgang, altsaa
Bevægelsen være det Afgjørende.
\par Endnu en Bemærkning til Overveielse. Der opstilles en fra flere
Sider (Jvfr. Phil. og Math. S. 23--24, Begrebet Nøiagtighed.
% "Bilag. S. 29": initial verified as Fraktur B (identical sort to
% "Begrebet" on the same page); both OCR witnesses misread it as V.
S. 45--47, „Afregn.“ gjennems. S. 7, Bilag. S. 29) belyst Sætning,
hvori Grundtanken er denne: For at tænke den continuerlige Overgang fra
Endeligt til $\infty$, er det nødvendigt at see paa Bevægelsen;
abstraherer man fra Bevægelsen, og holder sig til umiddelbare
Subtractionsdifferenser, bliver den continuerlige Overgang ufattelig.
Herimod optræder saa en Mathematiker, som tillige er overordentlig
Logiker, med den Bekjendtgjørelse, at den opstillede Sætning er falsk,
grundfalsk. Hvorledes bærer den mathematiske Logiker sig nu ad med at
bevise sin Paastand? Han paaviser, at Sætningens mathematiske Feil er
„uendelig stor“, idet han ved Ligningen:
% displayed equation, printed on its own line in the note; the numerator
% carries ∓ (minus-over-plus), as printed
\[ \mathrm{tg}\,90^{0} - \mathrm{tg}\,(90^{0} \pm \varepsilon) =
\frac{\sin\,(\mp\,\varepsilon)}{\cos 90^{0}\,\cos\,(90^{0} \pm
\varepsilon)} = \frac{1}{\cos 90^{0}} = \infty \]
udtrykkelig abstraherer fra Bevægelsen, og ved at abstrahere fra
Bevægelsen tydelig og uimodsigelig beviser, hvad den angrebne Sætning
netop stadfæster, at nemlig Beviset for en continuerlig Overgang
nødvendigviis maa, naar der abstraheres fra Bevægelsen, svigte. Den
forsøgte Gjendrivelse forvandler sig herved til en udtrykkelig
Stadfæstelse af det, der skulde gjendrives; altsaa: man vil paavise en
Sætnings „uendelige mathematiske Feil“, og see, saa kommer man ved et,
ja det forstaaer sig, naturligviis ubetydeligt mathematisk-logisk
Feilgreb netop til at bevise Sætningens mathematiske Gyldighed:
ordentlig Mathematik, overordentlig Logik!}.
% --- p. 64 ---
% (no new footnote mark on this page; its foot carries the unmarked
% conclusion of the p. 63 note, transcribed at the p. 63 anchor above)
At $\infty$, det i sig Ubegrændsede, kan være Grændse, beroer
væsentlig paa $0$, den bestemmende Negation, Grændsen selv. At $0$
ligger det Endelige uendelig nær, vil med andre Ord sige, at Grændsen
ligger det Begrændsede uendelig nær; men hvad der, med Hensyn paa den
opnaaelige Grændse, i en Ligning, som $\frac{1}{a-x} = y$, gjælder om
$a - x = 0$, gjælder jo ogsaa om $y = \infty$. Grændsen $\infty$ er
altsaa, seet fra dette Synspunkt, ligesaa opnaaelig, som Grændsen $0$.
Heraf
% --- p. 65 ---
\setcounter{footnote}{0}
forklares det saa igjen, at medens ingen endelig Størrelse kan betegne
nogen Overgang fra $+$ til $-$, kan $\infty$ derimod betegne denne
Overgang, ligesaavel som $0$, just fordi $0$ kan betegne den. „Da nu en
Overgang, hvad Modsætningen af Positivt og Negativt angaaer, kun faaer
Betydning derved, at det Positive gaaer over i det Negative, og
omvendt det Negative i det Positive, altsaa derved, at det Positive og
det Negative begge ere i Overgangen og desaarsag høre med i
Overgangens Begreb: saa burde saadanne Ligninger, som $\sin 0 = 0$,
% "cot." with a full stop at its first occurrence, "cot" without at the
% second; as printed.
$\mathrm{cot.}\ 0 = \infty$, strengt taget, skrives: $\sin\,(\pm\,0) =
\pm\,0$, $\mathrm{cot}\,(\pm\,0) = \pm\,\infty$, o.~s.~v. Ved denne
Adskillelse af $+$ og $-$ i $0$ og $\infty$ opkommer der i disse
Grændsebetegnelser selv en indre Forskjellighed, der er i høieste Grad
dialektisk. Betegnelsen $\pm\,0$ udtrykker jo, naar vi see nærmere til,
at det ene Nul har skilt sig ad i tre forskjellige Nuller, et $+0$, et
% the minus in "Negation af + og ÷" and in "hvori + og ÷ ere
% forsvundne" is printed with the Danish division-style minus ÷ (the
% other minuses on this page are printed —); transcribed as printed.
$-0$, og det rene neutrale $0$, den skarpe Negation af $+$ og $\div$;
den samme Adskillelse gjentager sig, hvad $\pm\,\infty$ angaaer; ogsaa
her er et $+\,\infty$, et $-\,\infty$, og et neutralt $\infty$, hvori
$+$ og $\div$ ere forsvundne. Naar det altsaa hedder, at $0$ og
$\infty$ ere Overgangsbetegnelser, kan dette ikke gjælde om det rene
$0$ eller det neutrale $\infty$, men kun om $\pm\,0$ og
$\pm\,\infty$\footnote{Kun ved at betragte Nullet fra 3 Synspunkter:
1) som det Positives Udgangspunkt, 2) som det Negatives Udgangspunkt,
og 3) som begges adskillende Eenhedspunkt, og lade disse 3
Bestemmelser reflectere sig i een Grundbestemmelse, tænker man
Nulbegrebet, det rene Grændsebegreb, som Overgangsbegreb.}. Den Maade,
% "1^∞": the superscript on 1 is a small ∞, verified on the image at 3x
% (both OCR witnesses fail on it: "1*" / "1°").
hvorpaa de saakaldte ubestemte Former: $\frac{0}{0}$, $0.\ \infty$,
$0^{0}$, $\infty^{0}$ og $1^{\infty}$, behandles i Mathematiken, viser
tydelig nok, hvad Overgangsformer ville sige.“ (Math. og Dialekt.
S. 72).

% head printed letterspaced, run-in, matching the Indhold's
% "γ) Articulationer" (~p. 65 per the Indhold; printed here on p. 65)
\emph{γ) Articulationer.} Som Overgangsformer betegne $0$ og $\infty$
altsaa Grændsens dialektiske Eenhed med det
% --- p. 66 ---
% (printed: "Be-" ends p. 65, "grændsede" begins p. 66; the word is
% completed here)
Begrændsede, og derved Bevægelsen: Overgang forudsætter Bevægelse.
Forsaavidt de endelige Størrelser bestemmes ved givne Tegn, kunne de,
uden at nogen Bevægelse mærkes, opfattes af Forestillingen og
fastholdes i deres Endelighed, endog i den Grad, at det seer ud, som
var det Uendelige kun en tom Abstraction, en udvortes, oversvævende
Negation; forsaavidt de endelige Størrelser derimod bestemmes ved
Relationer og opfattes som Begreber, viser Negationen sig som
Endeligheden iboende, og det Uendelige articuleres. „Saalænge
Forestillingen uvilkaarlig hæfter ved Tegnet 1, tager det sig ud, som
var Eenheden et umiddelbart Noget, et Tilværende, man kunde anskue og
ligefrem paapege; denne Forestillingens Umiddelbarhed forsvinder, naar
man begynder at opfatte Eenheden som et Forhold: $\frac{2}{2} =
\frac{3}{3} = \frac{4}{4}$.“ (Mathem. og Dial. S. 79). Hvad er altsaa
Eenheden? I og for sig er enhver Størrelse Eenheden: $a = a$, $b = b$,
$c = c\,....$; saamange forskjellige Størrelser, saamange forskjellige
Eenheder! Men ingen af de forskjellige Eenheder er Eenheden selv;
hvorledes komme vi da til den alle Eenheder bestemmende Eenhed? Ved den
i de ligestore Størrelser udtrykte Ligestorhed! Af $\frac{a}{a} = 1$,
$\frac{b}{b} = 1$, $\frac{c}{c} = 1\,....$ sees udtrykkelig, at
Eenheden netop er enhver Størrelses Ligestorhed med sig selv; ved
særlig at fremhæve den for hver Størrelse gjældende Ligestorhed, og
igjen sætte denne som Størrelse, have vi den almindelige Eenhed. Som
Ligestorhedens bestemte Udtryk, som Qvotient af lige Forholdsled, er
Eenheden uafhængig af Leddenes umiddelbare Forskjellighed; saalænge
hvert Led blot vedbliver at være sig selv liigt, kunne Leddene
forandres, voxe og aftage i det Uendelige, uden at slige Forandringer i
mindste Maade forandre Eenheden. Har man altsaa formedelst det
Begrændsede ($a = a$, $b = b\,....$): $\frac{a}{a} + \frac{b}{b} +
\frac{c}{c} + ... =$
% --- p. 67 ---
$1 + 1 + 1 + ...$, saa har man ogsaa formedelst det Grændseløse
% "Lim." capitalised at the first limit, "lim." lower-case at the
% second; as printed.
$\left(\mathrm{Lim.}\,\frac{x}{x} = \frac{\infty}{\infty}\right)$:
$\frac{\infty}{\infty} + \frac{\infty}{\infty} + \frac{\infty}{\infty}
+ ... = 1 + 1 + 1 + ...$ og formedelst Grændsen
$\left(\mathrm{lim.}\,\frac{x}{x} = \frac{0}{0}\right)$: $\frac{0}{0} +
\frac{0}{0} + \frac{0}{0} + ... = 1 + 1 + 1 + ...$ Af disse
Forandringernes Strømme hæver Eenheden sig frem som det Uforanderlige,
som Qvantiteringens rene Qvantum, det i Bevægelsen Ubevægelige.

Paa Eenheden af den almindelige og særskilte Eenhed, den dialektiske
Eenhed af Størrelsens Begreb og Størrelsens Værdi beroer Functionen;
enhver Størrelse er en Function af sig selv og Eenheden: $x = f(x) =
1. x$ (Jvfr. Math. og Dial. S. 6--7). Af $\frac{f(x)}{x} = 1$,
Udtrykket for Størrelsens Eenhed med sig selv i Almindelighed, den
\emph{almindelige Eenhed} ($f(x) = x$), følger
% in the print the quotient of integrals stands large in the running
% line, with the limits set above and below the integral signs
$\frac{\int_{f(x)=0}^{f(x)=1.a} df(x)}{\int_{x=0}^{x=a} dx} =
\frac{a}{a} = 1$, ɔ: Størrelsens Eenhed med sig selv som særskilt
Størrelse, den \emph{særskilte Eenhed} ($f(a) = a$).

Fremstilles Overgangen, idet $d.f(x) = dx$, ved $f(x) =
\int_{x=0}^{x=\infty} dx = \infty$, maa det for Alting fastholdes, at
her aldeles ikke kan tænkes paa en ligefrem Summation, men kun paa
Vorden, Tilbliven, intellectuel Bevægelse. Intet har fra Begyndelsen af
mere bidraget til at forrykke det afgjørende Synspunkt og
tilbagetrænge den naturlige Begrebsudvikling, end den Omstændighed, at
Operationsmechaniken altfor meget har behandlet Differentialet snart
som et umiddelbart Størrelsesatom, snart som et mystisk Nul, istedetfor
at fatte det
% --- p. 68 ---
\setcounter{footnote}{0}
som Overgangsmoment, og paa tilsvarende Maade betragtet Integrationen
blot som en Art Summation (Math. og Dial. S. 99--104), istedetfor at
bestemme den som en Restitution af Endeligt gjennem Uendeligt
formedelst den for Qvantiteringen gjældende Bevægelse. Hvor umuligt det
er at forestille sig $\int dx = dx + dx\,....\,= \infty\ dx$ som en
egentlig Sum\footnote{Vistnok kan $x$, seet under Synspunktet af
uendelig Discretion, \emph{tænkes opløst} i $dx + dx +\,....$, men Et
er det at \emph{tænke} den uendelige Discretion, et Andet at
\emph{addere} $dx$ til $dx$ --- i det Uendelige.}, sees ligefrem deraf,
at $dx$ som dialektisk Mellembestemmelse imellem $0$ og $x$ umulig kan
være Addend. For at være Addend maatte $dx$ have umiddelbar endelig
Tilværen, --- saaledes kan ikke blot $a$, men $\frac{a}{10000}$ være
Addend; men opfattet som tilværende er $dx$ en saare tvetydig
Størrelse, for stor til at være Intet, for lille til at være Noget.
Dette Tvetydige forvandler sig til et i Begrebet Tvesidigt, naar $dx$
opfattes som den Overgangsbestemmelse, hvorpaa al Bevægelse væsentlig
beroer. Som Overgangsbestemmelse kan $dx$ ikke i noget Øieblik staae
stille; dets Væren er Vorden, det forandrer sig og bliver enten $0$
eller $x$. Vil Nogen indvende, at man dog ikke saaledes ved et Spring
kan komme fra $dx$ til $x$, men først maa have $2dx$, $3dx$, $....$
indtil $ndx$, for $n = \infty$, da er det ikke vanskeligt at paavise
den Begrebsforvirring, som skjuler sig herunder.

For at articulere Bevægelsen bliver det nemlig nødvendigt at skjelne
imellem en enkelt og en dobbelt Vorden. I den enkelte Vorden er $n$ i
$ndx$ enten endeligt, eller uendeligt. Forsaavidt $n$ er endeligt, er
$ndx$ et Element, og som Element et almindeligt $dx$; skal $ndx$ blive
en endelig Størrelse, maa $n$ være $\infty$. Saalænge her endnu kun
spørges om det i
% --- p. 69 ---
Overgangen Almeenvæsentlige, sees her overalt paa den dialektiske
Forskjelseenhed af $dx$ og $x$, men derimod ikke paa Forskjellen
imellem $dx$ og $ndx$; ethvert $ndx$ er under Bevægelsen som Element i
Begreb med at blive et Endeligt, og som Endeligt at voxe ved et
Element. Hvor stort Elementet er, kan ved denne Sammenligning, efterdi
Sammenligningen er umiddelbar, ikke komme i Betragtning, hvor lille den
endelige Størrelse er, heller ikke; Overgangen beroer just derpaa, at
det størst mulige Element og den mindst mulige endelige Størrelse ere
Vexelbestemmelser. Bevægelsen gaaer i den enkelte Vorden fra $0$ til
$x$ ikke igjennem Myriader af $dx$, men igjennem $dx$. Ved at
forestille sig den i $\int dx = x$ liggende Opgave løst saaledes: $0 +
dx + dx +\,....\,+ dx = ndx = x$, hilder man sig i den Fixidee, at
Rækken skulde indeholde et særligt $dx$, hvis Plads netop betegnede den
værende Grændse, hvor Elementernes Sum for første Gang gav en endelig
Størrelse. Forskjellen mellem Addition og Bevægelse er her
paafaldende. Ved at addere $dx$ til $dx$, naaer man aldrig $x$; $x -
dx = x - 2dx = x - 3dx = x - ndx$, saalænge $n$ er endeligt, og under
Additionen er $n$ altid endeligt; ved Bevægelsen naaes derimod
umiddelbart og øieblikkelig fra $dx$ til $x$, idet $x$ er den mindst
mulige \emph{endelige} Størrelse. Det $x$, hvortil Bevægelsen
forvandler det dialektiske $dx$, er igjen dialektisk som Bevægelsen
selv; det i $x$ Dialektiske følger ligefrem deraf, at enhver Bevægelse
maa begynde, før den kan ende. Med den begyndende Bevægelse er $dx$
umiddelbart blevet et $x$, thi Bevægelsen er i sin Umiddelbarhed selv
Relation, Reflexion; men da Bevægelsen dog endnu kun er i sin
Begyndelse, er $x$ det mindst Mulige. Standser nu Bevægelsen, just som
det Begyndte skulde fuldføres, saa falder det Endelige igjen tilbage
til Elementet: det mindst mulige $x$ er da selv et $dx$. Fuldføres
derimod den begyndende Bevægelse, da naaes Resultatet, da voxer det
mindst
% --- p. 70 ---
\setcounter{footnote}{0}
mulige $x$ ud over et mindst Muligt, rodfæster sig i det Endelige og
bliver derved qvalitativt forskjelligt fra sit Element. Det $x$, der
fremkommer ved $\int dx = x$, har altsaa, hvor lille det end er (Math.
og Dial. S. 111--112), den fuldendte Bevægelse til Forudsætning, er af
en bestemt endelig Værdi, kan som et Endeligt deles i uendelig mange
uendelig smaa Dele og frembringer saa ved denne Mulighed det Skin at
være en Sum af utallige $dx$.

De eleatiske Indsigelser imod Bevægelsens Mulighed\footnote{Jvfr.
Math. og Dialekt. S. 131 o. fl.} lade sig i nærværende Sammenhæng
henføre til en udialektisk Identificeren af Begyndelse og Fuldendelse.
For at en Bevægelse kan fuldendes, maa den naturligviis være begyndt;
sættes nu den begyndende Bevægelse identisk med den fuldendte, saa
have vi som Modsigelsernes Cardinalpunkt denne vistnok urimelige
Sætning: „for at en Bevægelse kan fuldendes, maa den allerede forud
være fuldendt.“ Man fordrer, med andre Ord, en absolut Begyndelse; men
at fordre en absolut Begyndelse er en Selvmodsigelse, thi ingen
Begyndelse er eller kan være absolut. Under Bevægelsen er enhver
Begyndelse Fortsættelse, og enhver Fortsættelse en ny Begyndelse.
Begyndelsens og Fortsættelsens dialektiske Eenhed er Bevægelsen selv.
Under Discretionens Synspunkt viser enhver Bevægelse sig derfor ogsaa
som en Sum af uendelig mange Bevægelser. Seet fra dette Synspunkt
ligger der unegtelig imellem $x$ og $dx$ uendelig mange $dx$, og
ligesaa imellem $dx$ og $0$ uendelig mange $\mathrm{Lim.}\,
\frac{dx}{x}$; men for at indsee, hvorledes en saadan Discretion
opgaaer i Continuiteten, maae vi særlig betragte den fordoblede
Vorden, den Vorden, hvori Bevægelse bestemmer Bevægelse.

% --- p. 71 ---
Haves $f(x) = y = ax$, altsaa (Math. og Dial. S. 84--85) $dy = adx$,
da viser Forholdet imellem $dy$ og $dx$ hen paa en dobbelt Vorden, en
Vorden af $x$ ved Tilvæxten $dx = dx$, og en herved bestemt Vorden,
hvorunder det af $x$ afhængige $y$ voxer ved Tilvæxten $dy = adx$.
Ihvorvel $dy$ nu er et Element, og som Element en uendelig lille, i
Sammenligning med en hvilkensomhelst endelig Størrelse forsvindende,
Værdi, ligesaavel som $dx$, saa sees dog heraf, at det ene Element kan
være $a$ Gange større end det andet; er f.~Ex. $a = 2$, saa er $dy =
2dx$, og $dx = \frac{1}{2}dy$; er $a = 3$, saa er $dy = 3dx$, og $dx =
\frac{1}{3}dy$ o.~s.~fr. At der imellem Elementerne indbyrdes
imidlertid ikke kan anstilles nogen umiddelbar Sammenligning, at her af
slige Elementer Intet lader sig danne enten ved Sammenstykning eller
Udstykning, er let at indsee. Istedetfor af $dy = 2dx$ at aflede $dy -
dx = dx$, opgiver man Subtractionen som en ufrugtbar Symbolik,
anvender Divisionen og erholder $\frac{dy}{dx} = 2$,
Differentialcoefficienten, den Factor, hvormed $dx$ skal multipliceres
for at give $dy$ (Math. og Dial. S. 81). Hvilken afgjørende Betydning
for Articulationen af Uendeligt ved Endeligt det har, at Multiplication
og Division træde istedetfor Addition og Subtraction, lader sig uden
Vanskelighed oplyse. Størrelsen af en endelig Tidseenhed kan
umiddelbart angives, ligeledes kan Størrelsen af det Rum, der med en
vis endelig Hastighed gjennemløbes i en endelig Tid, umiddelbart
angives; da nu Størrelserne hver for sig ere umiddelbart givne, kan
Sammenligningen imellem dem ogsaa anstilles paa en saa umiddelbar
Maade, at Leddene selv bevare en vis Selvstændighed ligeoverfor den
Relation, hvori de befinde sig til hinanden. Er Hastigheden paa en
Bane 4 Miil, paa en anden 8 Miil i Timen, da falder den herved givne
Subtractionsdifferens som af sig selv. Anderledes, naar Tids- og
Rumselementer træde istedetfor de endelige Tids- og Rumseenheder.
% --- p. 72 ---
\setcounter{footnote}{0}
% (printed: "Hvor-" ends p. 71, "langt" begins p. 72; the word is
% completed here)
Hvorlangt er vel et Øieblik? Ja, det kan ikke umiddelbart angives, og
ligesaa lidt kan det umiddelbart angives, hvor langt et Stykke Vei der
med en given endelig Hastighed kan gjennemløbes i et eneste Øieblik:
istedetfor at holde sig i udvortes Relation, maae Elementerne forvandle
sig til Momenter i en gjennemgaaende Reflexion; hertil behøves
Multiplication og Division, Produkt og Qvotient. I Produkt og Qvotient
have Leddene allerede opgivet deres udvortes mod hinanden ligegyldige
Tilværen og forvandlet sig til reflecterede Forholdsmomenter. Saalænge
Produkt og Qvotient imidlertid kun dannes af endelige Størrelser, er
den tilsvarende Reflexion vistnok endnu indskrænket af den blot
forestillende Opfattelse\footnote{For den endelige Reflexion, den blot
forestillende Tænkning, tager det sig da ogsaa ud, som om Summen $a +
b = (a + b)$, var ligesaa reflecteret, som Produktet $a.b = ab$,
Differensen $a - b = c$, ligesaa reflecteret som Qvotienten
$\frac{a}{b} = c$.}; men idet Elementerne træde i de endelige
Størrelsers Sted, maa den sidste Rest af Forestillingens Umiddelbarhed
forsvinde, og Begrebet, det rene ved Articulationen tydeliggjorte
Begreb, gjælde. Lade vi $s = 4t$, $ds = 4dt$, $\frac{s}{t} =
\frac{ds}{dt} = 4$ være forskjellige Udtryk for een og samme
Hastighed, nemlig 4 Miil i Timen, saa indlyser det, at, uagtet
Hastigheden er den samme, staaer Hastighedsudtrykket $\frac{ds}{dt}$
dog i et ganske andet Forhold til Tænkningen, end Udtrykket:
$\frac{s}{t}$. Qvotienten $\frac{s}{t}$ holder sig endnu i den til
Forestillingen støttede Relation, men i $\frac{ds}{dt}$ beroer den
hele Indsigt paa gjennemgaaende Reflexion; i det ved dem bestemte
Forhold ere $ds$ og $dt$ Intet, slet Intet i og for sig, men Reflexer
af hinanden,
% --- p. 73 ---
og i denne Gjennemsigtighed bestemte ved hinanden, saaledes at de ikke
yde Forestillingen nogetsomhelst endeligt Tilhold.

Den i sig fordoblede Vorden bliver endnu rigere articuleret, naar den
til Bevægelsen svarende Hastighed med hvert Øieblik forandres. En
Bevægelse kan nemlig, hvad Hastigheden angaaer, være uendelig mange
Forandringer underkastet, og dog, naar ingen af disse Forandringer
pludselig indtræder, vise sig continuerlig. Continuiteten er betinget
deraf, at Forandringerne alle ere een og samme høiere Lov undergivne
og altsaa befinde sig i indbyrdes Afhængighedsforhold af hinanden.

% (the quotation opened here runs past p. 74; the closing mark falls in
% a later batch — this fragment carries one surplus opening quote on
% purpose)
„Forsaavidt $s$ er et i Tiden $t$ allerede gjennemløbet Rum, kan man
endnu ikke vide, om dette Rum er gjennemløbet med foranderlig eller
uforanderlig Hastighed. For at bestemme Hastigheden maa man give Tiden
en Tilvæxt, det vil sige: skaffe Tid til fortsat Bevægelse. Kaldes nu
Tilvæxten til Tiden $h$, haves, idet $s = f(t)$ Rækken:
% displayed equation, printed on its own line; denominators 1.2 and
% 1.2.3 as printed
\[ f(t + h) = s + \frac{ds}{dt}\,\frac{h}{1} +
\frac{d^{2}s}{dt^{2}}\,\frac{h^{2}}{1.2} +
\frac{d^{3}s}{dt^{3}}\,\frac{h^{3}}{1.2.3} +\,.... \]
Er Tilvæxten $h$ en endelig Størrelse og Hastigheden constant, vil
hele Rækken indskrænke sig til $f(t + h) = s + \frac{ds}{dt}\,
\frac{h}{1}$; er $h$ uendelig lille, liig $dt$, vil det Samme gjælde,
selv naar Hastigheden er variabel. Derimod ville, naar $h$ er endelig,
Rækkens Led kunne forøges i det Uendelige. Seet fra Begrebets
Synspunkt, angiver Rækken en Totalitet af objective
Reflexionsbestemmelser. De enkelte Led kunne, idet $s$ betyder Rummet,
$v$ Hastigheden og $\varphi$ Kraften, gjengives ved forskjellige til
hinanden svarende Udtryk:
% three-line displayed block, printed with the second and third lines
% aligned on =; set here as successive displays (no amsmath in the
% preamble)
\[ f(t + h) = s + \frac{ds}{dt}\,\frac{h}{1} +
\frac{d^{2}s}{dt^{2}}\,\frac{h^{2}}{1.2} +
\frac{d^{3}s}{dt^{3}}\,\frac{h^{3}}{1.2.3} +\,.... \]
\[ = .. + v\,\frac{h}{1} + \frac{dv}{dt}\,\frac{h^{2}}{1.2} +
\frac{d^{2}v}{dt^{2}}\,\frac{h^{3}}{1.2.3} +\,.... \]
\[ = .. +\,....\,+ \varphi\,\frac{h^{2}}{1.2} +
\frac{d.\varphi}{dt}\,\frac{h^{3}}{1.2.3} +\,.... \]
% --- p. 74 ---
\setcounter{footnote}{0}
Kategorierne: \emph{Bevægelse}, \emph{Hastighed}, \emph{Kraft} og
\emph{Forandring i Kraft} blive saaledes bestemte ved objectiv
Reflexion\footnote{I hvilken Grad en ellers udviklet Reflexion kan
forfeile Udviklingen af et Operationsbegreb, naar de Vink, den exacte
Analyse meddeler, aldeles oversees, har Hegel viist i sin speculative
Udtydning af Faldbevægelsen (Jvfr. Phil. Propæd. S. 161--167). Den
bekjendte Sætning, at de med jævnt voxende Hastighed gjennemløbne Rum
forholde sig som Qvadraterne af de anvendte Tider, udtyder Hegel
nemlig saaledes, at Qvotienten $\frac{s}{t^{2}}$ bliver en dialektisk
Conseqvens af Rums- og Tidsbegrebets abstract-aprioriske Forhold: „Die
der Einheit, als der Form der Zeit, entgegengesetzte Form des
Außereinander des Raums, und zwar ohne daß irgend eine andere
% "Qvadrat" (Danish-style Qv also in the German quotation): initial
% verified as Fraktur Q on the image; the German is set in Fraktur like
% the Danish, hence no antiqua marking.
Bestimmtheit sich einmischt, ist das Qvadrat: die Größe außer sich
% printer's defect: full stop after "vermehrend" before lower-case
% "aber" (comma expected; possibly a broken sort) — verified on the
% image at 2x; transcribed as printed.
kommend in eine zweite Dimension sich setzend, sich somit vermehrend.
aber nach \emph{keiner andern als ihrer eignen} Bestimmtheit ---
diesem Erweitern sich selbst zur Grenze machend, und in ihrem
Anderswerden so sich nur auf sich beziehend“. Det undgaaer da i dette
Tilfælde Hegels Opmærksomhed, at om ogsaa Bevægelsen i al
Almindelighed lader sig deducere af Forholdet mellem det rene Rum og
den rene Tid, en Deduction, der dog altid maa være betinget af
understøttende Biforestillinger, saa kan Tidens Qvadrat jo dog umulig
fremkomme uden en varierende Hastighed, og en varierende Hastighed
umulig uden en tilsvarende Aarsag, her en constant Kraft, der
paavirker den sig i Rummet efter Inertiens Lov bevægende Materie.}.
Umiddelbart betegner $s$ en Rumslængde; men i Ligningen $s = f(t)$
sees denne Rumslængde i Reflexion af en Tidslængde; den skylder jo
Tiden sin Tilværen, forsaavidt den nemlig er tilbagelagt i en given
Tid. Umiddelbart betegner $v$ en Hastighed, men i Udtrykket
$\frac{ds}{dt}$ er Hastighedsbegrebet bestemt ved objectiv Reflexion
af Rum og Tid i det givne Forhold. Umiddelbart er Kraften betegnet ved
$\varphi$, dernæst ved Reflexion af Hastigheden og Tiden i det
bestemte Forhold $\frac{dv}{dt}$, endelig, idet $v = \frac{ds}{dt}$
gaaer op i Reflexion
% (p. 74 ends mid-sentence; the sentence continues on p. 75)

% B. Philosophiens Forhold til Chemien: Forandringsproblemet begins ~p. 81.
% --- p. 75 ---
\setcounter{footnote}{0}
af Rum og Tid, ved en fordoblet Reflexion i Forholdet
$\frac{d^2s}{dt^2}$. Paa denne Maade kunne Kraftens Ændringer med
Hensyn paa Forholdet af Rum og Tid bestemmes som Reflexioner
af Reflexioner i det Uendelige“. (Phil. Propæd. S. 165--167).
% (the closing „…“ above ends a quotation opened on p. 74, in the
%  preceding batch; this fragment therefore closes one more quote
%  than it opens)

Continuitets- og Discretionsbegrebets dialektiske Vexelbestemmelse
fremlyser nu af den Maade, hvorpaa den continuerlige
Bevægelse articuleres, idet Elementerne selv ikke
blot kunne deles og deles, men endog paa givne Betingelser
% (the footnote below is marked *) on p. 75 and continues, without a
%  repeated mark, over the whole foot of p. 76)
opløses i Elementer af høiere og stedse høiere Orden\footnote{At
raisonnere saaledes: „Størrelsen $y$ kan tænkes deelt i uendelig
mange $dy$, lad gaae! men da $dy$ paa denne Maade selv bliver
den mindst mulige Deel af en Størrelse, kan $dy$ selv dog umulig
være deelbart, eftersom man ved at dele $dy$ jo maatte faae Dele,
der vare mindre end den mindst mulige“, --- gaaer ikke an; thi
$dy = a\,dx$, $dx = \frac{dy}{a}$, og, dersom vi vælge Functionen
$y = f(x) = x^2$, $dy = 2x\,dx$, $dx = \frac{dy}{2x}$, viser
tydelig nok, hvad det vil sige, at $dy$ er deelbart. Men dersom
et Element af første Orden ikke er absolut udeelbart, uagtet det
qva Element er \emph{uendelig lille}, saa følger det jo af sig
selv, at Elementer af anden Orden $d^2y$, tredie Orden $d^3y$,
4de Orden $d^4y$, $n$\textsuperscript{te} Orden $d^ny$ heller
ikke paa Grund af deres \emph{uendelige Lidenhed} kunne være
absolut udeelbare. Vistnok kan man i al Almindelighed sætte
$\frac{d^2y}{dx} = 0$, $\frac{d^3y}{dx^2} = 0$,
$\frac{d^4y}{dx^3} = 0$, $\frac{d^ny}{dx^{n-1}} = 0$, og hvad
hermed er eenstydigt $\frac{dx}{d^2y} = \infty$,
$\frac{dx^2}{d^3y} = \infty$, $\frac{dx^3}{d^4y} = \infty$,
$\frac{dx^{n-1}}{d^ny} = \infty$, idet man saa ved $0$ og
$\infty$ betegner, at Elementerne af forskjellig Orden ere
„\textit{incomparable}“ (Jvfr. Math. og Dialekt. S. 96); men
heraf følger naturligviis ingenlunde, at $\frac{d^2y}{dx^2}$
skulde være mindre end $\frac{dy}{dx}$, $\frac{d^ny}{dx^n}$
mindre end $\frac{d^{n-1}y}{dx^{n-1}}$; tvertimod! Relationerne
kunne endogsaa medføre $\frac{d^4y}{dx^4} > \frac{d^3y}{dx^3}$,
$\frac{d^3y}{dx^3} > \frac{d^2y}{dx^2}$,
$\frac{d^2y}{dx^2} > \frac{dy}{dx}$,
$\frac{d^ny}{dx^n} > \frac{d^{n-1}y}{dx^{n-1}}$. Jo større
Qvotienten er, i desto flere Dele maa Dividenden jo være deelt:
Dividendernes Eiendommelighed bestemmes ved de som Divisorer
valgte Maalestokke. Imellem Elementerne $dy$ og $d^2y$, $d^2y$
og $d^3y$ o. s. v. i de anførte Qvotienter er her ingen
umiddelbart bestemt Relation; den ifølge $y = f(x)$ bestemte
Relation fremkommer først derved, at Maalestokkene $dx^2$,
$dx^3$, $dx^4\ldots$ alle ere bestemte ved $dx = dx$. Er
saaledes $\frac{dy}{dx}$ constant, t. Ex. $= a$, saa er ogsaa
$\frac{dy}{dx} = \frac{\int dy}{\int dx} = \frac{y}{x} = a$, og
omvendt. Er $\frac{d^2y}{dx^2}$ constant, t. Ex. $= 2a$, saa er
ogsaa $\frac{\int\int d^2y}{\int\int dx^2}
= \frac{\int dy}{\int x\,dx} = \frac{y}{\frac{1}{2}x^2} = 2a$,
og omvendt; er $\frac{d^3y}{dx^3}$ constant t. Ex. $= 6a$ saa
er $\frac{\int\int\int d^3y}{\int\int\int dx^3}
= \frac{\int\int d^2y}{\int\int x\,dx^2}
= \frac{\int dy}{\int \frac{1}{2}x^2\,dx}
= \frac{y}{\frac{1}{6}x^3} = 6a$ og omvendt. Slige Relationer
kunne fortsættes i det Uendelige.}. Dersom
Rummet $s$ t. Ex. ikke efter Mulighed lod sig dele i lutter
$ds = ds$, ligesom Tiden $t$ i lutter $dt = dt$, vilde
Continuiteten selv og følgelig enhver continuerlig Bevægelse
være ufattelig, Forholdet imellem Hvile og Bevægelse,
hurtigere og langsommere Bevægelse aldeles ubegribeligt.

„Forskjellen mellem langsom og hurtig Bevægelse maa
aabenbart beroe derpaa, at den første har mere, den anden
% --- p. 76 ---
\setcounter{footnote}{0}
mindre Hvile i sig, at den første altsaa er Hvilen nærmere
end den anden; en Bevægelse med en uendelig lille Hastighed
er jo netop Hvilen saa uendelig nær, at den kan betragtes
som Hvile\footnote{Jvfr. Math. og Dialekt. S. 55--56.}. Ikke
desto mindre er Modsætningen mellem
Hvile og Bevægelse dog qvalitativ: hvad der hviler, bevæger
% --- p. 77 ---
sig ikke, hvad der bevæger sig, hviler ikke. Forkaste vi nu
de uendelige Størrelser, kan Hastighedsforskjellen kun forklares
deraf, at den rene Bevægelse endelig afbrydes af Hvilen. At
Legemet paa $A$ i Tidseenheden $t$ kun gjennemløber $s$, medens
Legemet paa $B$ gjennemløber $s_1 = 8s$, maa altsaa forklares
saaledes, at hiint Legeme kun bevæger sig i $\frac{1}{8}t$, men
derimod hviler i $\frac{7}{8}t$. Saalænge Tidseenheden, hvori
en saadan Afvexling foregaaer, er en Time, er Hvilens Afvexling
med Bevægelsen iøinefaldende; er Tidseenheden 1 Secund, er
Afvexlingen mindre paafaldende og tager sig ud som en stødviis
Bevægelse; i en Tidseenhed af $\frac{1}{100}$ Secund vil
Afvexlingen slet ikke mærkes, men Phænomenet vise sig som en
jævn Bevægelse \ldots\ Skal Bevægelsen paa $A$ være
continuerlig, maae Tiden og Rummet være deelbare i det
Uendelige; i dette Punkt har Aristoteles seet rigtigt. Men hvad
Aristoteles ikke har seet og ikke har kunnet see, fordi han
manglede det Støttepunkt, som kun det Uendeliges Analyse
kan give, er, at her ved denne Tidens og Rummets uendelige
Deelbarhed netop danne sig uendelig smaa Størrelser med
en saadan indbyrdes Forskjel, at de indgaae i bestemte
Endelighedsforhold med hinanden“. (Math. og Dialekt. S. 144--45).
Til den constante Hastighed: $v = a = \frac{ds}{dt}$ svarer en
eenfold Continuitet, en Continuitet, som vi, da den netop betegnes
derved, at Differentialcoefficienten af første Orden er constant,
her ville kalde \emph{Continuiteten af første Orden}.

Kan den eenfolde Continuitet kun begribes af en tilsvarende
eenfold Discretion, der gaaer i det Uendelige, saa
kan den i sig fordoblede Continuitet da heller ikke begribes
uden en tilsvarende Discretionsfordobling. Giver den med
hvert Øieblik forandrede Hastighed os Rækken: $v_1$, $v_2$, \ldots
% --- p. 78 ---
\setcounter{footnote}{0}
$v_{n-1}$, $v_n$ i $n$ paa hinanden følgende Øieblikke, saa at
$v_p - v_{p-1} = dv$, hvorved Rumstilvæxten $ds$ i hvert $dt$
successivt forvandles til $ds_1$, $ds_2 \ldots ds_n$, da er her,
efterdi $v_n \mathrel{{>}\atop{<}} v_{n-1}$ og
$ds_n \mathrel{{>}\atop{<}} ds_{n-1}$, i disse to Rækker,
forsaavidt de betragtes \emph{umiddelbart}, ingen Continuitet at
opdage. Continuiteten, som fordrer, at det Samme maa uden
Afbrydelse gjentage sig i hvert Øieblik, viser sig først, naar
vi saavel i Hastigheden som i Rummet opdage Muligheden af et
under Forandringerne uforanderligt Element. Er det saaledes
givet, at Kraften $\varphi$ er constant, saa bliver ogsaa den derved
i hvert Øieblik, $dt = dt$, bevirkede Hastighedsforandring $dv$
constant; men det hertil svarende constante Rumselement er
ikke $ds$; det er, ifølge $dv\,dt = d^2s$, netop Elementet
% (as printed: no full stop after the footnote mark of "d²s*)" before
%  "Den ved Functionsloven" — the compositor dropped the period;
%  transcribed as printed)
$d^2s$\footnote{Ved at anvende Summationssymboliken føres man da
ogsaa gjennem en umiddelbar Summation af $d^2s$ til Udtrykket
for $s$ som Function af Tiden.\\
I 1ste $dt$ er Hastigheden $dv$, og bevirker Rumsforandringen $d^2s$\\
i 2det $dt$ \dotfill\ $2dv$ \dotfill\ $2d^2s$\\
i 3die $dt$ \dotfill\ $3dv$ \dotfill\ $3d^2s$\\
\leavevmode\dotfill\\
\leavevmode\dotfill\\
i $n$\textsuperscript{te} $dt$ \dotfill\ $ndv$ \dotfill\ $nd^2s$\\
Efter Forløbet af Tiden $t$ er her da, idet $t = ndt$,
$n = \frac{t}{dt}$, tilbagelagt et Rum:
$(1 + 2 + 3 \ldots + n)d^2s = \frac{1}{2}n(n+1)d^2s
= \frac{t}{dt}\left(1 + \frac{t}{dt}\right)\frac{d^2s}{2}$,
som, idet $\frac{t}{dt}$ er uendelig lille i Sammenligning med
$\left(\frac{t}{dt}\right)^2$ er liig
$\frac{t^2}{2}\frac{d^2s}{dt^2} = \frac{gt^2}{2}$, det
sædvanlige Udtryk for Faldloven.} Den ved Functionsloven,
$s = at^2$, articulerede Continuitet,
der forudsætter, at Differentialcoefficienten af anden
Orden er constant, kalde vi da med Hensyn hertil en
\emph{Continuitet af anden Orden}.
% --- p. 79 ---
\setcounter{footnote}{0}

Fordrer den fordoblede Continuitet en fordoblet Discretion,
saa fordrer paa tilsvarende Maade den tredobbelte
Continuitet, den nemlig, der svarer til $s = at^3$, en tredobbelt
Discretion, et Rumselement af tredie Orden. Her er
det nemlig ikke Kraften selv, $\varphi$, men Kraftforandringens
Aarsag, $\psi$, der er constant. Kraftforandringen, $d\varphi$, er da
naturligviis ogsaa constant. Det constante $d\varphi$ bevirker nu i
hvert Øieblik, $dt$, en constant Forandring i Hastighedstilvæxten,
som igjen bevirker en constant Forandring, $d^3s = d^3s$,
i Rumstilvæxtens Tilvæxt; $d^3s$ er altsaa i Rummet
det til denne Continuitetslov svarende Discretionselement.
Det er tredie Differentialcoefficient, $\frac{d^3s}{dt^3}$, som
her er constant. \emph{Continuiteten er af tredie Orden.}

Af de saaledes paaviste Articulationer sees, at Continuiteten
ingenlunde afbrydes derved, at flere og stedse flere Forandringer
inddrages under Continuitetsloven. Skulde Continuiteten
nogensinde briste under Vægten af Forandringernes
Mængde, da maatte der for Elementernes stigende Orden
omsider vise sig en uoverstigelig Skranke, uden at dette dog
hidrørte fra nogen constant Differentialcoefficient; men da
dette er en Modsigelse, da Elementernes Orden jo kan stige
og stige i det Uendelige, maa den til en constant
$n$\textsuperscript{te} Differentialcoefficient,
$\frac{d^ns}{dt^n} = C$, svarende Continuitet af
$n$\textsuperscript{te} Orden
gjælde, ogsaa for $n = \infty$\footnote{Er $s$ en bruden,
irrational eller transcendent Function af $t$,
t. Ex. $s = \cos t$, da haves her ikke nogen constant
Differentialcoefficient; men dette betyder ingenlunde, at
Continuiteten er brudt, men tvertimod, at den er saa uendelig
articuleret, at man ikke kan paavise det til en saadan
Continuitet svarende Discretionsatom isoleret.}.

Samle vi de her vundne Resultater under eet Overblik,
blive følgende Hovedpunkter især at fremhæve: 1) den
% --- p. 80 ---
% (p. 79 ends in the divided word "Om- / stændighed")
Omstændighed, at der imellem $dx$ og $x$ kan ligge utallige $dx$,
og, forsaavidt $x$ er afhængig variabel, imellem $0$ og $dx$,
utallige $d^2x$, imellem $0$ og $d^2x$ utallige $d^3x$ o.s.v.,
betegner vel, at der imellem de sammenhørende og samtidig
forsvindende Tilvæxter til den afhængige og uafhængige Variable
gives varierende Forhold, hvoraf da Differentialer af høiere
Orden med Nødvendighed fremgaae; men ingenlunde, at $0$,
$dx$ og $x$ (mindste Endeligt) skulde, hvad Overgangsbegrebet
i dets Almindelighed angaaer, være adskilte ved utallige
Mellemled: $dx$ kan, naar det i en til Uafhængigheden svarende
Umiddelbarhed skal repræsentere den mindst mulige Størrelse,
umuligt blive mindre uden at blive $0$, og heller ikke større,
uden at blive $x$. I sin \emph{Umiddelbarhed} er Elementet en
yderste Discretionsgrændse, et Størrelsesatom, og som saadant
udeelbart; reflecteret i sine Relationer til andre Elementer kan
det derimod vise sig uendelig deelbart. Paa Dialektiken af det
i sig udeelbare-deelbare Element beroer Bevægelsen. 2) Bevægelsens
Begyndelse skeer ikke umiddelbart med $0$, saalidt som med $x$
for sig, eller med $dx$, men med en negativ-concret Eenhed af
alle tre Momenter. „I det absolute Nu er Intet og skeer
Intet. I det relative Nu sættes det Uendelige som det Endeliges
Begyndelse: det Endeliges Begyndelse falder altsaa
sammen med, er i det mathematiske Øieblik samtidigt med
det Uendeliges Begyndelse. I det empiriske Øieblik sættes
det Endelige som fuldendt; men i og med denne det Endeliges
Fuldendelse er en Uendelighed selv fuldendt: det Endeliges
Fuldendelse falder altsaa sammen med, er i hvert virkeligt
Øieblik samtidigt med en Uendeligheds Fuldendelse.“ (Math. og
Dialekt. S. 134--35). 3) Endskjøndt Bevægelser af enhver
Art altid maae forudsætte idetmindste Biforestillinger, Reflexer,
af Tid og Rum, saa er Bevægelsen dog i streng mathematisk
Forstand at indskrænke til den rene Qvantiteren. Da nu en
Qvantiteren ligesaa nødvendigt forudsætter et qvantiterende
Qvantum, som Bevægelsen et Noget, der kan bevæge sig,
% --- p. 81 ---
medens dog det abstracte Qvantum, der afgiver Qvantiteringens
umiddelbare Forudsætning, just skylder Qvantiteringen
sin Tilværen, og derved kommer til at forudsætte det, hvis
Forudsætning det selv skulde være: saa indsees heraf, at Qvantiteten
som abstract Begreb først kan faae sin sande Betydning
i og ved et Tilværende, hvis Qvantitet den er.

Om den rene Mathematik hedder det, at den finder
„Anvendelse“ paa virkelige Gjenstande (i Mechanik, Astronomie,
% (as printed: no full stop after the closing parenthesis of
%  "Optik o. desl.)" — the sentence ends with "desl.)" and
%  "Væsentlig" begins the next)
Optik o. desl.) Væsentlig forstaaet forudsætter en saadan
„Anvendelse“, at det Mathematiske oprindelig er virkeliggjort
i de Gjenstande, hvorpaa Mathematiken anvendes. En Anvendelse
af Mathematik paa virkelige Gjenstande vilde være,
om ikke en reen Umulighed, saa dog en subjectiv Vilkaarlighed,
dersom det Mathematiske ikke selv med indre objectiv
Nødvendighed hørte med i Virkelighedens Væsen; men som
Moment i Virkelighedens Væsen hører det Mathematiske
tillige med i Philosophien.

% (section rule as printed; signature "6" at the foot of p. 81)
\section*{B.\\ Philosophiens Forhold til Chemien:\\
Forandringsproblemet.}
\addcontentsline{toc}{section}{B. Philosophiens Forhold til
Chemien: Forandringsproblemet}
% (short rule under the section head, as printed; the printed head
%  divides "Forandrings-" / "problemet." across its two lines and
%  matches the Indhold's reading)

For at fremstille Philosophiens Forhold til Chemien, vil
det i nærværende Sammenhæng være nødvendigt: a) at
charakterisere Chemien som særegen Videnskab, b) at paavise,
hvorledes Forandringsproblemet bliver til under den chemiske
Analyse, og endelig c) hvorledes Philosophien tilegner sig det
i Chemien udviklede Problem.

At Philosophien, hvormeget den end fordyber sig i et
Andet, dog igjennem dette Andet forholder sig til sig selv,
og forsaavidt bevæger sig paa sine egne Enemærker, vil
% --- p. 82 ---
da under Fremstillingens Gang være kjendeligt derpaa, at
intet Stofagtigt umiddelbart optages udenfra, men alle positive
Bestanddele tilegnes saaledes, at de til Væsenserkjendelsen,
hver paa sin Maade, afgive ideelle Bidrag og tjene
Begrebsudviklingens Dialektik.

\subsection*{a) Chemien som særegen Videnskab.}
\addcontentsline{toc}{subsection}{a) Chemien som særegen Videnskab}
% (head as printed, centred; matches the Indhold)

% (the block quotation below is letterspaced as printed; the citation
%  "(Phil. Propæd. S. 70)" is set normally and stays outside the \emph)
\emph{„Da Videnskaben har sin Rod i Menneskelivet,
hvor Phantasiens Idealer gaae forud for Tænkningens,
er det naturligt, at de forskjellige Videnskaber
fra først af komme til Verden i Indbildningens
Svøb og ere behæftede med uvidenskabelige
Tilsætninger. Saaledes var Astronomien
længe sammenvoxet med Astrologie, Chemien med
Alchemie, Philosophien med Mythologie. Vistnok
ligger det i Videnskabens Drift at hæve Forpupningen,
bryde sit Hylster og udrense de fremmedartede
Bestanddele. Men Dannelsen koster Tid,
Udrensningen Anstrengelse; de Hindringer, Menneskeaanden
her har at overvinde, ere af forskjellig
Art.“} (Phil. Propæd. S. 70).

Det er Videnskabens Opgave at trænge ind i Tingenes
Indre og igjennem Phænomenet at opdage Væsenet; men saa
anerkjendt denne Fordring end er i abstract Almindelighed,
saa vanskelig er dens Opfyldelse, naar Opmærksomheden vender
sig mod det Særegne, det Concrete, det Virkelige. Hvilket
Phænomen er vel mere bekjendt og mere iøinefaldende end
Ildphænomenet, og dog, hvad er det Indre i dette Ydre,
hvad er Flammens Væsen, hvad er Ilden? De Fleste have
vel, idetmindste som i Forbigaaende, hørt tale om, at der i
den nyere Tid ved Kunst er udskilt en Mængde Metal-Substanser,
man forhen ikke kjendte; blandt disse Substanser
hører da ogsaa Kalium, et tinhvidt, stærkt glindsende, meget
% --- p. 83 ---
\setcounter{footnote}{0}
blødt og smidigt Metal. Kaster man et Stykke af dette
Metal i Vandet, synker det ikke tilbunds, men bevæger sig
omkring paa Overfladen, antændes og brænder med en violet
Flamme; derefter bliver det, som det synes, aldeles borte:
hvad er nu her foregaaet? Er det Kalium, der har brændt
Vandet, eller Vandet, der har brændt Kalium? \emph{Hvad er i
Forbrændingens Phænomen det Brændende, og
hvad det Brændbare?} Til Besvarelsen af et saadant
Spørgsmaal har Videnskaben, naar vi, for at begynde med
Begyndelsen, søge vore første Chemikere iblandt Ægyptens
Præster, maattet forberede sig i henved 3000 Aar.

\subsection*{α) Chemiens Oprindelse\protect\footnote{\textit{Leçons
sur la philosophie chimique par Dumas, recueillies par
Bineau.} \textit{Bruxelles} 1839. Geschichte der Chemie von
Dr.~H. Kopp. 4 Bd. Braunschweig 1843--47.}.}
\addcontentsline{toc}{subsection}{α) Chemiens Oprindelse}
% (head as printed, centred, with the footnote mark *) before the
%  full stop; matches the Indhold. The note's French title through
%  "Bruxelles 1839." is antiqua; "Geschichte der Chemie …
%  Braunschweig 1843--47." is Fraktur. Signature "6*" at the foot
%  of p. 83)

Hos Ægypterne synes Chemiens Dyrkelse at have udgjort
en Bestanddeel af de Mysterier, der skjultes i Templerne;
hos Grækerne gik den chemiske Videns Elementer ligesom
andre Erfaringskundskaber over i Philosophien; i Middelalderen
var Chemien som Alchemie en Anviisning til at udgrunde
Naturens Gaade, og finde de Vises Steen (\textit{lapis
philosophorum}).

Da Indsigt i Naturkræfternes Sammenhæng ikke mindre
er af praktisk end theoretisk Interesse og ligesaa meget beroer
paa rig Erfaring som paa vedholdende Eftertanke, er det
forklarligt, at Chemiens Studium længe maatte være sammenvoxet
med forskjellige fremmedartede Bestræbelser, og denne
Green af Videnskaben længe tjene som Middel, inden den
kunde frigjøres til at være sit eget Maal og uddannes i reen
videnskabelig Aand. \emph{Mystisk Speculation, Famlen
efter Methode og Methodens Opdagelse} ere da de
Stadier, som her betegne Udviklingens Vei.
% --- p. 84 ---
\emph{αα) Mystisk Speculation: Oldtid og Middelalder.}
Uden Erfaring, ingen Tænkning; men naar Synstingene
kun flygtig overskues, fordi Iagttagelsens Kunst endnu
ikke er uddannet, forener den abstraherende Tænkning sig med
Phantasien og danner saa af faa tilfældige og upaalidelige
Erfaringer et Begrebssystem, hvis Indhold, hvad det Tilværende
angaaer, mere tilhører Indbildningen end Virkeligheden.
Vistnok kunne, hvad den græske Philosophie fra saa
mange Sider tilstrækkelig viser, Begrebsformer i abstracte
Retninger udvikles med Klarhed og skarp Conseqvens, uagtet
Erfaringens Indhold mangler; men spørges her om Existensbegreber,
de egentlige Realbegreber, indlyser det som af sig
selv, at Formudviklingen overalt maa være betinget af en
tilsvarende Indholdsudvikling.

Den Maade, hvorpaa Aristoteles begrunder Tilværelsen
af fire Elementer, viser just, hvor umuligt det er at bestemme
Tingens Væsen uden at kjende Tingen. De elementære
Egenskaber: hed, kold, tør og fugtig, fremhæves nemlig
saaledes, at hver af dem sættes som Væsenhed for sig, en
substantiel Form, der bestemmer Materien; saaledes giver da
en samtidig Forening af Tørhed og Varme \emph{Ildelementet},
af Varme og Fugtighed \emph{Luftelementet}, af Fugtighed og
Kulde \emph{Vandelementet}, af Kulde og Tørhed
\emph{Jordelementet}. Det er en mærkelig Forvexling af
Virkning med Aarsag, af det Afledede med det Oprindelige, vi
her møde. Tingen selv, det Existerende, fremkommer jo dog
ikke ved Egenskabernes, altsaa heller ikke ved Grundegenskabernes
Forening, men omvendt bliver Egenskabernes og da ogsaa
Grundegenskabernes Forening sat i og med og formedelst Tingen.
Her er jo ikke et System af almindelige, for sig bestaaende
Egenskaber, der indvirke paa Tingene, men et System af
Ting, der igjennem deres Egenskaber indvirke paa hinanden
gjensidig; det er f. Ex. ikke Egenskaben: Tørhed, der tørrer
Tingen, saalidt som det er Egenskaben: Varme, der opvarmer
% --- p. 85 ---
% (a small stray mark stands after "Aarsag" in the first line — read
%  as an ink speck, not a printed stop; rejected reading: "Aarsag.")
den; Varmens og Tørhedens Aarsag er at søge, ikke i en
almindelig Væsenhed, men i en betinget Virken. Hvilken
Indflydelse en saadan fra Philosophien, navnlig fra Aristoteles,
nedarvet Misviisning maatte udøve paa den i dybe Anelser
og dunkle men dristige Bestræbelser gjærende Middelalder, er
det, hvad Chemien særligt angaaer, da heller ikke vanskeligt
at forestille sig.

Saadanne Stofsubstanser, som Qviksølv og Svovl, maae,
forsaavidt de høre blandt de udvortes virkelige Ting, jo vistnok
vise sig som Erfaringsgjenstande, der kræve en tilsvarende
Iagttagelse; men da Speculationen i slige Stoffer allerhelst
kun seer en Sandseliggjørelse af visse usandselige, i Materien
afprægede Væsenheder, skete det efterhaanden, at Alchemisternes
\textit{Mercurius} og \textit{Sulphur} kom til at betegne, ikke nærmest det
i Virkeligheden tilværende Qviksølv og Svovl, men fornemmelig
disse Stoffers metaphysiske Charakterer. Er nu Tingen,
der efter sin Natur skulde være en Iagttagelsens Gjenstand,
forvandlet til speculativ Væsenhed, saa ligger det nær at betragte
det Sandselige som flygtigt Billede af det Oversandselige,
det Naturlige som Symbol paa det Aandelige. Det
sorte Pulver, der fremkommer, naar Svovl og Qviksølv rives
sammen, gaaer ved Glødning over til en rød Masse, nemlig
Zinnober; i eet og samme Stof ere altsaa Rødt og Sort,
Symbolerne paa Lys og Mørke, det gode og det onde Princip,
synlige. Under den Forudsætning, at der i alle Metaller
fandtes \textit{Mercurius}, det Bestandiges, det Uopløseliges
Charakteer, lysende i Metalglandsen, og \textit{Sulphur}, det
Brændbare, Flygtige, Foranderlige, og at Metallernes indbyrdes
Forskjellighed fornemmelig beroede paa en særegen Blanding og
„Fixation“ af disse Naturer, førtes Alchemisterne stedse videre
i deres Bestræbelse for at omdanne, forvandle og forædle
Metallerne.

„Hvad dengang blev betragtet som Magikernes høieste
videnskabelige Stræben, gik ud paa at rense det
% --- p. 86 ---
\setcounter{footnote}{0}
% (p. 85 ends "Guddommeligt-", opening the hyphenated compound that
%  p. 86 completes)
Guddommeligt-Ligeartede, som forefandtes i den syndige Natur, forenet ved
modstræbende Elementer; paa det at det, der laae skjult som
det Guddommelige og Opholdende, i enhver indvortes Form
kunde virke frit. Dette er Alchemien --- ingen tilfældig
vilkaarligt udtænkt, men en aldeles væsentlig Bestanddeel af den
herskende Physik. Alle Physikere søgte de Vises Steen, thi
der gaves dengang ingen anden Physik, og kunde ingen anden
opstaae. Fremvirkningen af denne al Tilværelses ædleste
Kjærne var ligesaavel af religiøs Art, som et physisk Experiment,
og denne almeenherskende Bestræbelse giver det mest
slaaende Beviis for Aandens Bundenhed til det Jordiske.
Det saaledes Rensede, hvori den oprindelige Skabningskraft
concentrerede sig, maatte, naar den anvendtes paa Makrokosmus,
fremvirke de ædleste Materier, Ædelstene, og fremfor
Alt Guldet; men, anvendt paa Mikrokosmus, maatte den paa
samme Grund, da nemlig det oprindelige opholdende Princip
% (as printed: a stray full stop stands after "Sundheden" before
%  "og" — transcribed as printed; rejected reading: a comma with a
%  broken tail)
i begge var det samme, befordre Sundheden. og forlænge Livet.“
(H. Steffens)\footnote{Jvfr. „\emph{Naturvidenskabens Forhold
til Tidsaldere og deres Philosophie}“, en bedømmende Anmeldelse
af H. C. Ørsted. Aanden i Naturen. 2. S. 161--62.}.

Mystisk, vidunderlig og gaadefuld maatte den Substans
jo vistnok være, der ved den blotte Berøring skulde betvinge
\textit{Sulphur}, fortrylle \textit{Mercurius}, \textit{Metallorum
Draco}, og forvandle
ethvert, endog det ringeste Metal til det pure Guld.
De Vises Steen, den store Elixir, \textit{Magisterium magnum},
den røde Tinctur, der har Evnen til at forædle Metallerne
efter en saa uhyre Maalestok, --- \textit{mare tingerem, si
mercurius esset}, siger Raimundus Lullus, --- og helbrede alle
Sygdomme, dannes da heller ikke som et Aggregat, men voxer
indenfra, organisk, som et levende Væsen. Intet Under, at
Opdagelsen af en saadan Hemmelighed omsider (Basilius
Valentinus) viste sig at være en ligesaa religiøs som videnskabelig
% --- p. 87 ---
Opgave. Stenen, der lyste med et Naadens som med et
Naturens Lys, udbredte et magisk Skjær over den hele
Tilværelse; det jordiske Liv med dets Sorger og Lidelser var da,
seet i dette Lys, som en „Digestion“, en Renselse ved
% (printer's defect: "shvor" for "hvor" — a stray long s (Fraktur)
%  stands fused before the h; transcribed as printed)
„Fermentation“, Graven det Sted, shvor en „Putrefaction“, der
kun forstyrrede det menneskelige Legemes uædle Bestanddele,
skulde virke, medens det i Sjælen Udødelige fremgik som ved en
„Sublimation“ af det ædlere Væsen, en \textit{Quinta Essentia}.

\emph{ββ) Famlen efter Methode: qvalitativ Analyse.}
Alchemisternes Drømmerier pegede i to Retninger: den
\emph{metallurgiske} og den \emph{medicinale}. At de Vises Steen,
naar den engang opdagedes, skulde helbrede alle Sygdomme,
var naturligviis en Phantasie, der umulig kunde give en
Opgave; men under denne Phantasie ulmede dog en Gnist af
noget Rationelt, en Tanke, der, om end i phantastisk Form,
kom til Gjennembrud i den saakaldte Iatrochemie, Lægechemien;
Tanken er, at de organiske Functioner beroe paa chemiske
Processer. Vistnok vare det 16de og 17de Aarhundredes
Chemikere --- Paracelsus, Van Helmont, Sylvius --- endnu
ingenlunde istand til at danne sig noget klart Begreb enten
om de chemiske Processer eller de organiske Functioner, og
altsaa heller ikke istand til at udfinde nogen bestemt paalidelig
eller følgerigtig Fremgangsmaade: hvad en Sylvius forklarer
om den abnorme Fremhersken af Syrer og Ludsalte i Legemets
Safter, er i Grunden ligesaa utilforladeligt, som hvad en
Paracelsus aabenbarer os i den Forkyndelse, at enhver Deel
af det menneskelige Legeme har sit eiendommelige Svovl, sin
egen Merkurius og sit eget Salt, at Overvægt af Svovl
giver Feber og Pest, Overvægt af Qviksølv Tungsind og
Lamhed, Overskud af Salt Vattersot o. s. v.; men den nye
Retning gav dog ikkedestomindre Undersøgelserne et nyt
Opsving. At trække „Qvintessensen“ ud af Lægeplanterne,
at tilberede de mange Slags Tincturer, Essenser, Extracter,
% --- p. 88 ---
Decocter og Safter vare Bestræbelser, der vel maatte egne
sig til at skjærpe Iagttagelsesaanden, mangfoldiggjøre Forsøgene
og lære den enkelte Forsker at øve sit Blik og sin
Dømmekraft. Det ligger i Videnskabens Natur, at den kun
i samme Grad kan blive sig sit Formaal bevidst, som den
kan blive sig den til Formaalet svarende Fremgangsmaade,
Methoden, bevidst, og omvendt. Lægechemiens Formaal, skjøndt
i og for sig fornuftigt, maa, saalænge Forskningen ikke formaaer
at tilveiebringe Midlerne, saalænge den Undersøgelsens
Vei, der fører til Maalet, endnu ikke er funden, altid vise
sig fra en mere eller mindre phantastisk Side; men
Fremskridtsdialektiken er, at der ogsaa paa Vildfarelsens Vei
findes Sandheder, blandt Andet den Sandhed, at her er en
Vildfarelse. Da Iatrochemien var saavidt gjennemført, at
dens Huulhed og Charlatanerier ret kunde vise sig, var
Chemien selv ved at følge det andet Spor, det ældgamle
metallurgiske, kommen til det Vendepunkt, at den omsider kunde
optræde selvstændig og begynde at operere paa egen Haand.

Chemiens Emancipation er kjendelig paa den forandrede
Opgave; thi Opgaven er nu i Hovedsagen hverken at forvandle
Metallerne til Guld, ei heller at finde Medicinens
Arcanum, et Universalmiddel mod alle Sygdomme, men Opgaven
er at undersøge de for Stoffernes Adskillelse og Forbindelse
gjældende Betingelser og Love og fremfor Alt at besvare
det Spørgsmaal: hvorpaa beroer Forbrændingen? Men
saa videnskabelig og derhos eiendommelig chemisk denne
Opgave end stiller sig, saa var det dog umuligt Andet, end at
de første Forsøg med dens Løsning maatte falde noget
tvivlsomt ud. Den fra Alchemisterne nedarvede Forudsætning, at
der i alle Metaller gaves et hypothetisk Grundstof, som
betingede Forbrændingen, vedblev endnu at gjælde; heller ikke
kunde man saa let blive den Fordom qvit, at Forbrændingen
absolut maatte være en Forstyrrelse, en Opløsning, en
Adskillelse, saa at der af det brændende Legeme gik Noget bort,
% --- p. 89 ---
\setcounter{footnote}{0}
det nemlig, der gik op i Flammen. Deri er imidlertid
Forbrændingsproblemet væsentlig forskjelligt fra de iatrochemiske
% (doubtful: a faint dot stands above the comma after "Phantasier",
%  which would make it a semicolon; read as a comma — the scan's
%  text layer agrees, and the syntax expects it. Rejected: ";")
Phantasier, at det, hvormeget Phantastisk der end fra
Begyndelsen af indblander sig deri, dog ligger Forskningen saa nær,
at det ved sin egen rationelle Charakteer selv efterhaanden kan
anvise Veien og føre ind i Methoden.

De Forandringer, alle brændbare Legemer, de organiske
saavelsom de uorganiske, undergaae i Ilden, sammenfattede
Stahl\footnote{Georg Ernst Stahl, født i Anspach 1660, blev
1687 udnævnt til Livlæge hos Hertugen af Weimar, og, efter
Friedr. Hofmanns Tilskyndelse, 1694 Professor i Medicin ved
Universitetet i Halle.} under et eneste almeengjældende
Synspunkt, idet
han udledte Forbrændingen af et alle Legemer iboende Stof,
et Brændbarhedsstof, nemlig: Phlogiston. Saaledes dannede
sig den fra sidste Halvdeel af 17de til henimod Slutningen
af 18de Aarhundrede --- Stahl, Priestley --- charakteristiske
Lære om Phlogiston, den phlogistiske Theorie. Ved at prøve
Analogierne, fremhæve det Almene, drage theoretiske Slutninger
af Phænomenernes Iagttagelse, og anbringe mere
omfattende rationelle Synsmaader samlede Videnskaben Kræfter
og hævede sig saa Trin for Trin i sin eiendommelige
Skikkelse.

Men trods den lange Række af Forsøg, som alene
Forbrændingsspørgsmaalet maatte fremkalde, trods den
Mangfoldighed af Discussioner, det hypothetiske Stof, Phlogiston,
om det nu, hvad Stahl benegtede, var Svovl, om det ikke
snarere var en Lysmaterie, eller om det, hvad Andre fandt
rimeligt, maaskee var Brint, maatte foranledige --- Discussioner,
som med al deres Svæven dog altid adskilte
sig fra tomme Speculationer derved, at de alle dreiede sig
om virkelige Kjendsgjerninger og tilsvarende Forudsætninger:
saa vare Undersøgelserne dog endnu eensidige og ufuldkomne,
saalænge der blot lagdes Eftertryk paa Qvaliteten af de givne

% [text to be added: pp. 90--104]

% [text to be added: pp. 105--119]

% [text to be added: pp. 120--134]

% [text to be added: pp. 135--149]

% [text to be added: pp. 150--166]

% Resultat begins p. 167.
% [text to be added: pp. 167--176]

% C. Philosophiens Forhold til Physiologien: Livsproblemet begins p. 177.
% [text to be added: pp. 177--191]

% a) Livets Væsen region (α Organismen from ~p. 208).
% [text to be added: pp. 192--206]

% [text to be added: pp. 207--221]

% [text to be added: pp. 222--236]

% β) Den organiske Type region.
% [text to be added: pp. 237--251]

% [text to be added: pp. 252--266]

% γ) Det organiske Formaal region.
% [text to be added: pp. 267--281]

% b) Livets Trin begins ~p. 286; α) Plantelivet from ~p. 291.
% [text to be added: pp. 282--296]

% [text to be added: pp. 297--311]

% β) Dyrelivet begins ~p. 328; the Darwin passage falls on pp. 326--328
% (cf. texts/nielsen/propaedeutik-darwin/, transcribed earlier from the
% British Library scan — collate rather than re-derive).
% [text to be added: pp. 312--326]

% [text to be added: pp. 327--341]

% [text to be added: pp. 342--356]

% [text to be added: pp. 357--371]

% [text to be added: pp. 372--386]

% γ) Menneskelivet begins ~p. 399.
% [text to be added: pp. 387--401]

% [text to be added: pp. 402--416]

% c) Sjæl og Legeme begins ~p. 423.
% [text to be added: pp. 417--431]

% β) Psychophysiske Undersøgelser region (Phrenologie ~p. 438,
% Psychophysik ~p. 444, Mathematisk Psychologie ~p. 453).
% [text to be added: pp. 432--446]

% [text to be added: pp. 447--461]

% γ) Den menneskelige Videns Grændser begins ~p. 468.
% [text to be added: pp. 462--476]

% III. Philosophien og de akademiske Studier begins p. 477
% (Videnskabelig Tænkning 477; Videnskabernes Sammenhæng 479;
% Videnskabelig Dannelse 481). Chapter head carried by this fragment.
% [text to be added: pp. 477--482]

%% ============================================================
%%  INDHOLD (unpaginated, PDF 11--12) — transcribed LAST, from the image,
%%  per playbook §6: do not trust the contents OCR, and record any
%%  divergence between the Indhold and the printed heads as a datum.
%% ============================================================
% [Indhold: to be added when the body is complete]

\end{document}
